% =====================================================================
% model_appendix.tex — replaces Appendix B ("Details on the Structural Model")
% of the paper, sections B.1 and B.2. B.3 (untargeted moment derivation) and
% the EV eco-score appendix are unchanged and should follow this file.
%
% % =====================================================================
% model_appendix.tex — replaces Appendix B ("Details on the Structural Model")
% of the paper, sections B.1 and B.2. B.3 (untargeted moment derivation) and
% the EV eco-score appendix are unchanged and should follow this file.
%
% % =====================================================================
% model_appendix.tex — replaces Appendix B ("Details on the Structural Model")
% of the paper, sections B.1 and B.2. B.3 (untargeted moment derivation) and
% the EV eco-score appendix are unchanged and should follow this file.
%
% % =====================================================================
% model_appendix.tex — replaces Appendix B ("Details on the Structural Model")
% of the paper, sections B.1 and B.2. B.3 (untargeted moment derivation) and
% the EV eco-score appendix are unchanged and should follow this file.
%
% \input{model_appendix} from the main file. Requires amsmath, amssymb, amsthm,
% booktabs, enumitem and natbib. Add to the preamble if not already present:
%
%   \newtheorem{proposition}{Proposition}
%   \newtheorem{lemma}{Lemma}
%   \newtheorem{assumption}{Assumption}
%   \newtheorem{remark}{Remark}
%   \DeclareMathOperator*{\argmin}{arg\,min}
%
% Labels used by core_model.tex: app:model-derivation, app:model-foreign,
% app:model-calibration, app:model-prop, app:model-zeros, app:model-bounds.
% =====================================================================

\section{Details on the Structural Model}
\label{app:model}

\subsection{Derivation}
\label{app:model-derivation}

\paragraph{Setup.} The economy is divided into $R+1$ regions, populated by two types of firms:
producers of a final good and producers of intermediate goods used as inputs in the production
of the final good. The $(R+1)$th region represents the rest of the world. The analysis primarily
focuses on the matching of firms across domestic regions, controlling for competition from
foreign firms using their shares in total sales as reference. We assume that each domestic
region hosts at most one representative firm in the downstream sector, identified by its region
of origin. Intermediate goods are produced across $S$ sectors, with each region potentially
hosting many suppliers in each sector. We define a market $(s,r')$ as the collection of firms
located in region $r'$ that belong to upstream sector $s$.

Domestic regions are partitioned into attraction areas,
$\mathcal{R}=\bigsqcup_{a\in\mathcal{A}}\mathcal{R}_{a}$, where $\mathcal{R}_{a}$ is the set of
regions whose nearest downstream producer is the anchor of area $a$, and $a(r')$ denotes the
area of region $r'$.

\paragraph{Demand for final manufacturing products.} Firms in the downstream sector each own the
blueprint to produce a unique differentiated variety of the final good. The market for the final
good is characterized by monopolistic competition and iso-elastic demand functions:
\begin{equation}
q_{r}=\left(\frac{p_{r}}{P}\right)^{\varepsilon-1}\frac{E}{P}\,\delta_{r},
\end{equation}
where $q_{r}$ denotes the quantity demanded of the variety produced in region $r$ at price
$p_{r}$, given nominal consumption $E$ and the downstream sector's price index $P$. The
parameter $\varepsilon$ measures the price elasticity of sales, and $\delta_{r}$ is an exogenous
component of demand, varying across downstream producers and later used to simulate demand
shocks and their diffusion along production linkages. We abstract from foreign demand and assume
that the domestic market for final consumption goods is perfectly integrated.

\paragraph{Technology in the downstream sector.} Production of final goods requires the assembly
of intermediates, combined with labor, under a nested CES structure:
\begin{equation}
y_{r}=A_{r}\Big[\Omega_{L}^{\frac{1}{\lambda}}\ell_{r}^{\frac{\lambda-1}{\lambda}}
+(1-\Omega_{L})^{\frac{1}{\lambda}}\big(q^{I}_{r}\big)^{\frac{\lambda-1}{\lambda}}\Big]^{\frac{\lambda}{\lambda-1}},
\end{equation}
where $\ell_{r}$ is labor demand and $q^{I}_{r}$ the firm's intermediate input consumption. The
parameter $\Omega_{L}$ governs the share of labor in total production costs, while $\lambda$ is
the elasticity of substitution between labor and intermediates. Firm productivity $A_{r}$ is
exogenous. Intermediate consumption is itself a nested CES aggregator across sectors and, within
each sector, across varieties:
\begin{equation}
q^{I}_{r}=\Big[\sum_{s=1}^{S}\Omega_{s}^{\frac{1}{\nu}}q_{rs}^{\frac{\nu-1}{\nu}}\Big]^{\frac{\nu}{\nu-1}},
\qquad
q_{rs}=\Big[\frac{1}{N_{s}}\sum_{\rho=1}^{N_{s}}q_{rs}(\rho)^{\frac{\nu_{s}-1}{\nu_{s}}}\Big]^{\frac{\nu_{s}}{\nu_{s}-1}},
\label{eq:app-ces}
\end{equation}
where $\Omega_{s}$ governs the share of sector $s$ in intermediate consumption, $\nu$ is the
elasticity of substitution across sectors and $\nu_{s}$ that across varieties within a sector.

The only change relative to the continuum formulation is the inner aggregator: the integral over
$j\in[0,1]$ is replaced by an average over the $N_{s}$ varieties. Writing it as an
\emph{average} rather than a sum is a normalization of the love-of-variety level, discussed and
justified in Appendix~\ref{app:model-price}; it leaves every sourcing share, and hence the
entire extensive margin, unchanged.

Given the nested CES structure, the nominal demand for intermediate products is
\begin{equation}
x_{rs}(\rho)=\Omega_{s}(1-\Omega_{L})
\left(\frac{p_{rs}(\rho)}{p_{rs}}\right)^{1-\nu_{s}}
\left(\frac{p_{rs}}{p_{r}}\right)^{1-\nu}
\left(\frac{p_{r}}{c_{r}}\right)^{1-\lambda}\frac{c_{r}y_{r}}{A_{r}},
\end{equation}
with $p_{rs}\equiv\big[\tfrac{1}{N_{s}}\sum_{\rho}p_{rs}(\rho)^{1-\nu_{s}}\big]^{\frac{1}{1-\nu_{s}}}$
the cost index for intermediates in sector $s$,
$p_{r}\equiv\big[\sum_{s}\Omega_{s}p_{rs}^{1-\nu}\big]^{\frac{1}{1-\nu}}$ the cost index for
intermediates, and
$c_{r}\equiv\big[\Omega_{L}w_{r}^{1-\lambda}+(1-\Omega_{L})p_{r}^{1-\lambda}\big]^{\frac{1}{1-\lambda}}$
the unit cost of the firm. In equilibrium $y_{r}=q_{r}$.

\paragraph{Technology in upstream sectors.} Intermediate goods are produced with labor using a
constant-returns-to-scale technology. Wages in sector $s$ and region $r'$ are denoted $w_{r's}$
and taken as given. As in Eaton and Kortum (2002) we assume that varieties produced within and
across regions are perfect substitutes and that productivity draws are independent. The
departure from the continuum case is that each sector contains $N_{s}$ varieties.

\begin{assumption}[Granularity]
\label{as:granular}
Each upstream sector $s$ consists of $N_{s}\in\mathbb{N}$ horizontally differentiated varieties
$\rho=1,\dots,N_{s}$. Within each variety $\rho$ and each region $r'$ there is a continuum of
firms whose productivities above $z$ form a Poisson process with mean $T_{r's}z^{-\theta}$.
Draws are independent across varieties, regions and sectors.
\end{assumption}

The continuum of firms is what keeps the model tractable. The number of firms above $z$ in
$(r',s,\rho)$ is Poisson with mean $T_{r's}z^{-\theta}$, so the most productive one satisfies
\begin{equation}
\Pr\big(\varphi_{r'\rho s}\le z\big)=\Pr(\text{no firm above }z)=e^{-T_{r's}z^{-\theta}},
\label{eq:app-champion}
\end{equation}
that is, $\varphi_{r'\rho s}$ is exactly $\mathrm{Fr\acute{e}chet}(T_{r's},\theta)$. One
Fr\'echet draw per (region, variety) pair \emph{is} the champion of a continuum of firms, not an
approximation to it. Since the iceberg cost scales all firms of region $r'$ identically, that
champion is the unique region-$r'$ firm that ever serves any buyer sourcing $\rho$ from $r'$:
\textbf{a firm is a (region, variety) champion}, with delivered cost
$c_{r'\rho s}(r)=w_{r's}\tau_{r'rs}/\varphi_{r'\rho s}$. This is the ``fixed number of
independent Ricardian draws'' form of granular Eaton--Kortum \citep[cf.][]{eks2012,bejk2003}.

\begin{assumption}[Area-level comparative advantage]
\label{as:geo}
Comparative advantage is area-specific: $T_{r's}=T_{a(r')s}$ for every $r'\in\mathcal{R}_{a}$,
with one parameter per (sector, area). Productivity draws remain independent across regions ---
two regions in the same area share the Fr\'echet scale, not the realization of their champions.
\end{assumption}

Assumption~\ref{as:geo} does not make regions within an area interchangeable. Each runs its own
Ricardian competition for each variety, and two regions at different distances from a buyer have
different win probabilities. What it removes is unobserved technological heterogeneity within an
area, which is what allows the model to explain the same variation the extensive-margin
regression uses.

\paragraph{Equilibrium.} The resolution follows the same steps as in Eaton and Kortum (2002).
Given the Fr\'echet assumption, the minimum price at which suppliers from $(r',s)$ can serve
market $r$ is distributed
\begin{equation}
G_{r'rs}(p)=\Pr\Big(\frac{w_{r's}\tau_{r'rs}}{\varphi_{r'\rho s}}\le p\Big)
=1-e^{-\left[T_{a(r')s}(w_{r's}\tau_{r'rs})^{-\theta}\right]p^{\theta}},
\end{equation}
so the distribution of input prices for downstream firms in region $r$ is
\begin{equation}
G_{rs}(p)=1-\prod_{r'=1}^{R+1}\big[1-G_{r'rs}(p)\big]=1-e^{-\Phi_{rs}p^{\theta}},
\qquad
\Phi_{rs}=\sum_{a'}T_{a's}\sum_{r'\in\mathcal{R}_{a'}}(w_{r's}\tau_{r'rs})^{-\theta}.
\label{eq:app-Phi}
\end{equation}
The probability that a downstream firm in location $r$ picks a supplier in market $(r',s)$ is
thus
\begin{equation}
\gamma_{r'rs}=\frac{T_{a(r')s}(w_{r's}\tau_{r'rs})^{-\theta}}{\Phi_{rs}} .
\label{eq:app-share}
\end{equation}
Taking the perspective of the supplier, the probability of being the lowest-cost supplier in
market $(r,s)$, given location $r'$ and productivity $z$, is
\begin{equation}
\rho_{r'rs}(z)=e^{-\Phi_{rs}(w_{r's}\tau_{r'rs})^{\theta}z^{-\theta}}
=e^{-T_{a(r')s}\gamma^{-1}_{r'rs}z^{-\theta}} .
\label{eq:app-rho}
\end{equation}
Because the conditional distribution of prices is not origin-specific,
$G_{r'rs}(p)=G_{rs}(p)$, and $\gamma_{r'rs}$ is also the expected share of intermediates from
$r'$ purchased by a downstream firm in $r$. The expected cost indices are
\begin{equation}
p_{rs}=\Big[\Gamma\Big(\frac{\theta+1-\nu_{s}}{\theta}\Big)\Big]^{\frac{1}{1-\nu_{s}}}\Phi_{rs}^{-1/\theta},
\qquad
p_{r}=\Big[\sum_{s}\Omega_{s}p_{rs}^{1-\nu}\Big]^{\frac{1}{1-\nu}},
\qquad
c_{r}=\Big[\Omega_{L}w_{r}^{1-\lambda}+(1-\Omega_{L})p_{r}^{1-\lambda}\Big]^{\frac{1}{1-\lambda}},
\end{equation}
where the first expression is the certainty-equivalent index of Appendix~\ref{app:model-price}.
Bilateral flows are
\begin{equation}
X_{r'rs}=\gamma_{r'rs}\,x_{rs}
=\gamma_{r'rs}\,\Omega_{s}(1-\Omega_{L})
\left(\frac{p_{rs}}{p_{r}}\right)^{1-\nu}\left(\frac{p_{r}}{c_{r}}\right)^{1-\lambda}\frac{c_{r}y_{r}}{A_{r}} .
\end{equation}
Since we do not observe these bilateral flows directly, we cannot use this equation to estimate
the model's structural parameters and rely instead on simulated moments.

\subsection{The extensive margin and the two classes of zeros}
\label{app:model-zeros}

Across the $N_{s}$ varieties the win events are i.i.d., so with $K_{r'rs}$ the number of
sector-$s$ varieties buyer $r$ sources from $r'$,
\begin{equation}
K_{r'rs}\sim\mathrm{Bin}\big(N_{s},\gamma_{r'rs}\big),
\qquad
\Pr(K_{r'rs}=0)=(1-\gamma_{r'rs})^{N_{s}}>0 .
\end{equation}
Under a continuum this probability is $\lim_{N\to\infty}(1-\gamma_{r'rs})^{N}=0$ for every
origin with $T_{r's}>0$: the extensive margin is degenerate and the observed zeros are not a
model outcome. Finiteness is exactly what restores it.

\begin{remark}[Dependence across buyers]
\label{rem:dependence}
For a given variety the champion costs are common across buyers; only $\tau_{r'rs}$ differs. Win
events are therefore positively dependent across buyers, and the union probability $q_{r's}$ is
not the product of the per-market marginals. The per-market marginal above is nonetheless exact,
and $q_{r's}$ is evaluated by simulation. This matters because the survey records whether a firm
supplies the industry, not which buyer it serves: the empirical outcome is the union, and the
model counterpart is computed as the same union.
\end{remark}

Assumption~\ref{as:geo} sorts observed zeros into two economically distinct classes, and the
distinction determines what is estimated.

\begin{table}[htbp]
\centering
\small
\begin{tabular}{@{}p{0.21\textwidth}p{0.24\textwidth}p{0.19\textwidth}p{0.28\textwidth}@{}}
\toprule
\textbf{Zero} & \textbf{Treatment} & \textbf{Probability} & \textbf{Role}\\
\midrule
Whole attraction area $a$ has no supplier in $s$
& \emph{Structural}: $T_{as}=0$, $(s,a)$ outside the active set
& $1$ by construction
& None --- absorbed by the $a\times s$ fixed effect (Lemma~\ref{lem:fe}); no parameter estimated\\[4pt]
Region $r'$ inside an active area has no supplier
& \emph{Endogenous}: inherits $T_{as}>0$, own distance $d_{r'}$
& $(1-q_{r's})^{N_{s}}$
& \textbf{The identifying variation}: disciplined by distance alone\\
\bottomrule
\end{tabular}
\caption{The two classes of extensive-margin zeros. The second row is exactly the set of regions
that host firms but no supplier. Under region-level comparative advantage these had to be
assigned $T\equiv0$, making their zeros exogenous; under Assumption~\ref{as:geo} they inherit
their area's comparative advantage and their zeros become model predictions.}
\label{tab:app-hierarchy}
\end{table}

The reason the fixed effect dictates this split is the following.

\begin{lemma}[All-zero groups carry no information about the distance slope]
\label{lem:fe}
Index observations by firm $i$ within fixed-effect group $g=(a,s)$, with outcome $y_{i}$,
regressor $x_{i}$, weights $\omega_{i}>0$ and group effect $\phi_{g}$.
\begin{enumerate}[nosep,topsep=2pt,label=(\roman*)]
\item \emph{Complementary log-log.} The log-likelihood of a group in which $y\equiv0$ is
$\ell_{g}=-e^{\phi_{g}}\sum_{i}e^{\beta x_{i}}$, whose supremum over $\phi_{g}$ is attained as
$\phi_{g}\to-\infty$ with $\ell_{g}\to0$ and $\partial\ell_{g}/\partial\beta\to0$: the group
drops out of the concentrated likelihood and contributes exactly nothing to $\hat\beta$.
\item \emph{Linear probability model.} By Frisch--Waugh--Lovell an all-zero group has
$\tilde y\equiv0$, contributing zero to the numerator and a positive amount to the denominator,
so $\hat\beta=\varkappa\hat\beta^{+}$ with
$\varkappa=\big(\sum_{a\in\mathcal{A}^{+}}\sum\omega\tilde x^{2}\big)/\big(\sum_{a}\sum\omega\tilde x^{2}\big)\in(0,1]$,
where $\mathcal{A}^{+}$ is the set of areas hosting at least one supplier.
\end{enumerate}
\end{lemma}

Under the cloglog link --- the link the model implies --- $\alpha$ is therefore identified
exclusively from variation \emph{within} attraction areas that host at least one supplier. Under
a linear probability model the empty areas enter only as a mechanical dilution of the same
within-area signal. Either way, an area with no supplier at all carries no information about the
distance slope, which is why the model does not attempt to explain it.

\subsection{The firm-level reduced form}
\label{app:model-prop}

The object the estimation matches is the union event \eqref{eq:q}: whether a firm supplies the
downstream industry at all. That probability has no tractable closed form. The single-buyer case
does, and it is what gives the empirical cloglog its structural interpretation; we state it
first and then say precisely what survives when there are several buyers.

\begin{proposition}[Firm-level reduced form, single buyer]
\label{prop:firm}
Let a firm be a (region, variety) champion with productivity $z$ located in $r'$, and suppose it
faces a single downstream buyer, its area's anchor $a$. Writing
$\Phi^{-r'}_{as}=\sum_{r''\neq r'}T_{a(r'')s}(w_{r''s}\tau_{r''as})^{-\theta}$ for the
competitors' cost scale,
\begin{equation}
\Pr\big(\text{the firm supplies}\mid z\big)
=\exp\Big(-\Phi^{-r'}_{as}\big(w_{r's}\tau_{r'as}\big)^{\theta}z^{-\theta}\Big),
\label{eq:app-firm-win}
\end{equation}
so the probability that it does \emph{not} supply is a complementary log-log function
\begin{equation}
\Pr\big(\text{no supply}\mid z\big)=1-\exp\big(-\exp(\eta)\big),
\qquad
\eta=\underbrace{\ln\Phi^{-r'}_{as}}_{\text{area}\times\text{sector FE}}
+\theta\ln w_{r's}
+\theta\alpha\,\ln d_{r'a}
-\theta\,\ln z .
\label{eq:app-firm-cloglog}
\end{equation}
The variety count $N_{s}$ does not appear.
\end{proposition}

Three consequences, and they are the backbone of the estimation.

\begin{enumerate}[nosep,topsep=3pt,label=(\alph*)]
\item \textbf{The cloglog is the model's own auxiliary model.} A cloglog of the ``does not
supply'' indicator on log distance and log own productivity, with an area$\times$sector fixed
effect, is exactly \eqref{eq:app-firm-cloglog} when a firm faces one buyer, with distance
coefficient $+\theta\alpha$ and productivity coefficient $-\theta$.
\item \textbf{$\alpha$ and $N_{s}$ separate.} $N_{s}$ is absent from
\eqref{eq:app-firm-cloglog}: conditioning on the firm removes it. The regression identifies
$\alpha$; the count level identifies $N_{s}$; there is no residual coupling.
\item \textbf{The productivity coefficient is a free over-identifying test.} It equals
$-\theta$. Extending the moment vector to include it would let $\theta$ be estimated rather than
calibrated.
\end{enumerate}

\paragraph{Several buyers: what survives.} Our surveys record supply to the industry, not to a
named buyer, so the empirical outcome is the union event $\bigcup_{r}W_{r}$, where $W_{r}$
denotes winning buyer $r$. Its probability is \emph{not} $1-\prod_{r}[1-\rho_{r'rs}(z)]$: that
expression would require the $W_{r}$ to be independent, and they are not.

The dependence is positive and structural: conditional on $z$, the firm's delivered cost to
every buyer is the same number scaled by $\tau_{r'rs}$, so a firm cheap enough to beat the
competition in one market tends to beat it everywhere. The exact object is the
inclusion--exclusion sum
\begin{equation}
\Pr\Big(\bigcup_{r}W_{r}\,\Big|\,z\Big)
=\sum_{\emptyset\neq\mathcal{S}\subseteq\{1,\dots,R\}}(-1)^{|\mathcal{S}|+1}
\exp\Big(-z^{-\theta}\max_{r\in\mathcal{S}}\big\{\Phi^{-r'}_{rs}(w_{r's}\tau_{r'rs})^{\theta}\big\}\Big),
\label{eq:app-inclexcl}
\end{equation}
where the maximum appears because $\cap_{r\in\mathcal{S}}W_{r}$ requires the firm to beat the
\emph{tightest} of the thresholds it faces. Equation~\eqref{eq:app-inclexcl} has $2^{R}-1$
terms. Many are redundant --- if the threshold for buyer $r$ dominates that for $r''$ then
$W_{r}\subseteq W_{r''}$ and the pair collapses --- and pruning the dominated buyers reduces the
sum to the non-dominated set, but that set is still large enough in our geography that
evaluating \eqref{eq:app-inclexcl} inside an optimizer is not practical. The independence
approximation $1-\prod_{r}(1-\rho_{r'rs}(z))$, which is what a closed-form estimator would have
to use, is a strict \emph{under}-statement of \eqref{eq:app-inclexcl} by positive dependence, and
the error is largest exactly where the extensive margin has content, at high $z$.

This is the reason the model is estimated by simulated rather than analytical method of moments.
Simulation evaluates the union exactly --- each variety's champion costs are drawn once and
compared against every buyer, so the dependence is reproduced by construction --- at the cost of
simulation noise, which the weighting matrix accounts for. We therefore treat
\eqref{eq:app-firm-cloglog} as an auxiliary model rather than an estimating equation: the same
cloglog is run on simulated and observed data, with the union as the outcome on both sides, and
the coefficients are matched. This is consistent for $\alpha$ whether or not the link is exact,
and it is why the distance coefficient enters the moment vector rather than being inverted
analytically.

The outcome convention is dictated by the unit of observation rather than chosen. Simulating the
exact object at $\theta=1$, $\alpha=0.30$, with firms as (region $\times$ variety) champions,
the ``does not supply'' outcome recovers $\hat\beta_{\ln d}=+0.327$ and
$\hat\beta_{\ln z}=-1.051$ against truths $\theta\alpha=0.30$ and $-\theta=-1.0$; the reversed
outcome recovers neither ($-0.426$ and $+1.452$).

\subsection{The count moment and bounds on the variety count}
\label{app:model-bounds}

The zeros are a cell-level object. With $\mathcal{R}^{+}_{s}$ the regions inside active
attraction areas, the empirical moment is the share of those regions hosting at most $n$
suppliers,
\begin{equation}
\bar G_{s}(n)=\frac{1}{|\mathcal{R}^{+}_{s}|}
\#\big\{r'\in\mathcal{R}^{+}_{s}:K_{r's}\le n\big\},
\qquad
\mathbb{E}\big[\bar G_{s}(0)\big]=\frac{1}{|\mathcal{R}^{+}_{s}|}\sum_{r'}(1-q_{r's})^{N_{s}} .
\label{eq:app-count}
\end{equation}
Writing $\mu_{r's}=N_{s}q_{r's}$ and using the Poisson approximation,
$\partial\bar G_{s}(K)/\partial N_{s}
=-|\mathcal{R}^{+}_{s}|^{-1}\sum_{r'}q_{r's}e^{-\mu_{r's}}\mu_{r's}^{K}/K!$, which peaks at
$K\approx\mu_{r's}$ and vanishes as $\bar G_{s}(K)\to1$: saturated thresholds carry no
identifying variation and a poorly estimated variance, so we use the empty-region share
$\bar G_{s}(0)$, which is both the most interpretable threshold and typically near-optimal.

Because the iceberg cost never changes which firm is cheapest within a region, each variety
generates one distinct supplying firm per winning origin. So if $N^{\mathrm{obs}}_{s}$ is the
number of distinct upstream firms observed to supply the downstream industry,
$N^{\mathrm{obs}}_{s}=\sum_{\rho}\#\{\text{origins winning }\rho\text{ somewhere}\}\in[N_{s},N_{s}R^{D}]$,
where $R^{D}$ is the number of downstream buyers, hence
\begin{equation}
\frac{N^{\mathrm{obs}}_{s}}{R^{D}}\;\le\;N_{s}\;\le\;N^{\mathrm{obs}}_{s} .
\label{eq:app-bounds}
\end{equation}
The bounds are tight: the upper when every variety is globally sourced from a single origin, the
lower when each of the $R^{D}$ buyers sources it from a different origin. Since
$\mathbb{E}[\sum_{r'}K_{r's}]=N_{s}\sum_{r'}q_{r's}$, matching the observed firm count gives a second,
closed-form estimate $N^{\mathrm{count}}_{s}=N^{\mathrm{obs}}_{s}/\sum_{r'}q_{r's}$ that lies
inside \eqref{eq:app-bounds} by construction --- a free over-identifying check.

\paragraph{Degrees of freedom.} Fix a sector, let $A^{+}=|\mathcal{A}^{+}_{s}|$ and let
$R^{\mathrm{act}}_{a}$ be the number of regions in area $a$ with a strictly positive observed
share. The $T$ block has $A^{+}-1$ free elements after the reference normalization, against
$\sum_{a}R^{\mathrm{act}}_{a}$ share moments, so Assumption~\ref{as:geo} converts what were free
parameters into
\begin{equation}
\sum_{a\in\mathcal{A}^{+}}\big(R^{\mathrm{act}}_{a}-1\big)
\quad\text{over-identifying restrictions,}
\end{equation}
one per active region beyond the first in each area. Each is a within-area statement: given the
area aggregate, the split across its regions must follow the distance profile. Under
region-level comparative advantage this count is zero.

\subsection{Structural estimation: the moments}
\label{app:model-moments}

We use six sets of theoretical moments, observed in the data using a combination of IO tables
and micro-level data.

\begin{enumerate}[topsep=3pt,itemsep=3pt,label=\arabic*.]
\item \textbf{Aggregate labor share at the level of the industry:}
\begin{equation}
\tilde\Omega_{L}\equiv\frac{\sum_{r}w_{r}\ell_{r}}{\sum_{r}c_{r}A_{r}^{-1}y_{r}}
=\sum_{r}\frac{c_{r}A_{r}^{-1}y_{r}}{\sum_{r}c_{r}A_{r}^{-1}y_{r}}\,\Omega_{L}
\left(\frac{w_{r}}{c_{r}}\right)^{1-\lambda}.
\end{equation}
In the model it is a sales-weighted average of firm-level labor shares. In the data we use the
share of labor costs in total costs from IO tables.

\item \textbf{Aggregate share of sector $s$ in the industry's input purchases:}
\begin{equation}
\tilde\Omega_{s}\equiv\frac{\sum_{r}p_{rs}q_{rs}}{\sum_{r}\sum_{s}p_{rs}q_{rs}}
=\sum_{r}\frac{p_{r}q^{I}_{r}}{\sum_{r}p_{r}q^{I}_{r}}\,\Omega_{s}
\left(\frac{p_{rs}}{p_{r}}\right)^{1-\nu}.
\end{equation}
In the data we use technical coefficients for the downstream industry from IO tables.

\item \textbf{The share of each region in downstream sales:}
\begin{equation}
\pi_{r}\equiv\frac{c_{r}A_{r}^{-1}y_{r}}{\sum_{r}c_{r}A_{r}^{-1}y_{r}}
=\left(\frac{p_{r}}{P}\right)^{-\varepsilon}
=\frac{(c_{r}A_{r}^{-1})^{-\varepsilon}}{\sum_{r}(c_{r}A_{r}^{-1})^{-\varepsilon}},
\end{equation}
recovered from sales data on the population of downstream producers.

\item \textbf{The extensive-margin distance coefficient.} The cloglog of
\eqref{eq:app-firm-cloglog}, estimated on the simulated firm population with the same
specification as Equation~(1) in the data: outcome ``does not supply'', regressors log distance
to the area anchor and log own productivity, area$\times$sector fixed effects.

\item \textbf{Area-aggregated sourcing shares.} The share of the downstream industry's purchases
of sector $s$ sourced from area $a$:
\begin{equation}
\tilde\gamma_{as}\equiv\sum_{r'\in\mathcal{R}_{a}}\frac{\sum_{r}X_{r'rs}}{\sum_{r}\sum_{r''}X_{r''rs}}
=\sum_{r'\in\mathcal{R}_{a}}\sum_{r}\frac{\sum_{r''}X_{r''rs}}{\sum_{r''}\sum_{r}X_{r''rs}}\,\gamma_{r'rs}.
\label{eq:app-gamma-area}
\end{equation}
In the data we reconstitute firm-level sales to the downstream industry using sales-share data,
$x_{i(r',s)d}=\mathbf{1}(\mathrm{Sup}_{i(r',s)})\,a^{D}_{d\,i(r',s)}\,x_{i(s,r')}$, sum over all
firms in the area and normalize by the sum over all firms from sector $s$. The aggregation runs
over \emph{every} region of the area, including those hosting no supplier, on both the model and
the data side.

\item \textbf{The count moment} $\bar G_{s}(0)$ of \eqref{eq:app-count}: the share of regions
inside active attraction areas that host no supplier of sector $s$.
\end{enumerate}

Moments 1--3 and 5 are the same objects as in the continuum version of the model, with 5 now
aggregated to the area. Moments 4 and 6 are what the granular structure adds.

\paragraph{Notes on foreign parameters.}
\label{app:model-foreign}
We do not have the data to estimate parameters relating to foreign firms. Instead we take the
share of foreign inputs in sourcing as an exogenous parameter used in the simulated data, so the
vector of parameters to be estimated excludes region $R+1$. Moments 1--4 and 6 do not involve
$R+1$. For moment 5, with $w_{R+1,s}$ the share of foreign inputs in sector $s$'s input
purchases,
\begin{equation}
\tilde\gamma_{as}=(1-w_{R+1,s})\sum_{r'\in\mathcal{R}_{a}}\sum_{r}
\frac{\sum_{r''\neq R+1}X_{r''rs}}{\sum_{r''\neq R+1}\sum_{r}X_{r''rs}}\,\tilde\gamma_{r'rs},
\qquad
\tilde\gamma_{r'rs}=\frac{T_{a(r')s}(w_{r's}\tau_{r'rs})^{-\theta}}{\tilde\Phi_{rs}},
\end{equation}
with $\tilde\Phi_{rs}=\sum_{r'=1}^{R}T_{a(r')s}(w_{r's}\tau_{r'rs})^{-\theta}$. Foreign
competition still enters the win probabilities through $\Phi_{rs}$ in \eqref{eq:app-Phi}; what
is not estimated is foreign comparative advantage.

\paragraph{Calibrated parameters.}
\label{app:model-calibration}
The remaining parameters $\{\varepsilon,\theta,\lambda,\nu,\{\nu_{s}\},\{w_{r's}\}\}$ are
calibrated. We use the estimates from Fontagn\'e et al.\ (2022) for $\varepsilon$: the sales
elasticity for motor vehicles and aerospace is respectively $-9$ and $-16$. The shape $\theta$
of the Fr\'echet distribution is calibrated following Antr\`as et al.\ (2017). Elasticities of
substitution in production follow the literature, notably Baqaee and Farhi (2019): $\lambda=0.5$
for the elasticity between labor and intermediate inputs and $\nu=0.9$ across sectors. Local
wages are calibrated using labor cost information from employer--employee data in 2019
(DADS/BTS). Note that Proposition~\ref{prop:firm}(c) makes $\theta$ over-identified: the
coefficient on log own productivity in moment 4 should equal $-\theta$, and we report it.

\subsection{Estimation}
\label{app:model-estimation}

\paragraph{Simulation algorithm.} Given the calibrated elasticities and a candidate parameter
vector, we simulate the productivity champions of every (region, variety) pair, solve for all
bilateral trade flows --- from each potential supplier to each downstream region --- and
construct the simulated moments. All six blocks are computed from a \emph{single} Ricardian
solve on one set of draws. The SMM estimator is
\begin{equation}
\hat\mu=\argmin_{\mu}\;\big(\tilde M(\mu)-\hat M\big)'W\big(\tilde M(\mu)-\hat M\big),
\label{eq:app-smm}
\end{equation}
where $\tilde M$ (resp.\ $\hat M$) denotes the vector of theoretical (resp.\ empirical) moments
and $W$ is a weighting matrix. Optimal weights assign lower importance to imprecisely estimated
moments.

The variety count never enters the simulation. It is confined to two objects, both closed form,
and that confinement is the content of Proposition~\ref{prop:firm} together with:

\begin{lemma}[$q_{r's}$ is free of $N_{s}$]
\label{lem:qfree}
Winning a variety is the comparison
$w_{r's}\tau_{r'rs}/\varphi_{r'\rho s}<\min_{r''\neq r'}w_{r''s}\tau_{r''rs}/\varphi_{r''\rho s}$,
so $q_{r's}=q_{r's}(T_{\cdot s},w,\tau,\theta)$ and nothing else. The variety count, the price
index, downstream costs and expenditures do not appear.
\end{lemma}

The economics is that winning is a statement about \emph{relative costs}, whereas $N_{s}$ and
$P_{rs}$ are aggregators built after the winners are known: expenditure allocates value across
winners, it does not choose them. The one condition that would break the lemma is endogenous
upstream wages, since $w$ enters the comparison directly; wages are taken as data.

Together the two results say that $N_{s}$ reaches the moment vector only through
$\bar G_{s}(0)$ and through $\widehat N_{s}$ itself, and that both are functions of $q_{r's}$
alone. Writing $k_{r's}$ for the number of the $m$ simulated varieties won somewhere by region
$r'$ --- by Lemma~\ref{lem:qfree} an estimate of $q_{r's}$ valid whatever the variety count ---
the count moment is evaluated by the unbiased combinatorial estimator
\begin{equation}
\bar G_{s}(n)=\frac{1}{|\mathcal{R}^{+}_{s}|}\sum_{r'\in\mathcal{R}^{+}_{s}}
\frac{\binom{m-k_{r's}}{n}}{\binom{m}{n}},
\qquad
\mathbb{E}\left[\frac{\binom{m-k}{n}}{\binom{m}{n}}\right]=(1-q)^{n}\ \ \text{for } n\le m,
\label{eq:app-Gclosed}
\end{equation}
which is the probability that $n$ varieties drawn without replacement from the $m$ simulated
ones were all lost. We use \eqref{eq:app-Gclosed} rather than the plug-in $(1-\hat q_{r's})^{n}$
because the latter is convex in $\hat q$ and therefore biased upward. Cross-cell dependence ---
two regions in a sector face the same variety draws --- does not affect either expression, since
an expectation of an average is the average of expectations whatever the dependence. No economy
of size $N_{s}$ is ever drawn and no replication is needed.

\begin{remark}[What is given up]
\label{rem:binding}
Moment 4 is then the coefficient of a regression run on the full simulated sample, i.e.\ the
continuum limit $\beta(\mathbb{E}[y])$ rather than $\mathbb{E}[\hat\beta(N_{s})]$. The finite-sample bias of the
auxiliary cloglog --- which an average of $\hat\beta$ over realized economies of size $N_{s}$
would have cancelled against the same bias in the data --- is therefore not differenced out. It
is a measurable quantity, of the order of $2$--$4\%$ of $\beta$ at data-consistent group sizes,
and we report it alongside $\hat\alpha$. The alternative is not free: it would treat simulation
noise on one block by replication while every other block relies on the accuracy of the draw
design, and it would sit oddly beside Appendix~\ref{app:model-price}, where the value block is
already evaluated at the certainty-equivalent index rather than on a realized economy.
\end{remark}

\paragraph{Concentrating out the variety count.} $N_{s}$ is estimated by the same weighted
objective as every other parameter, concentrated rather than searched jointly:
\begin{equation}
\widehat N(\mu_{-N})=\argmin_{N\in\prod_{s}[N^{\mathrm{lo}}_{s},N^{\mathrm{hi}}_{s}]\cap\mathbb{N}}
\;r(\mu_{-N},N)'W\,r(\mu_{-N},N),
\label{eq:app-profile}
\end{equation}
with the bounds of \eqref{eq:app-bounds}. This is a profiled extremum estimator: $\widehat N_{s}$
is an integer at every evaluation, so integrality is never relaxed and never rounded, and by
Lemma~\ref{lem:qfree} the inner minimization needs no re-simulation. Since $\bar G_{s}(0;n)$ is
closed form and decreasing in $n$, \eqref{eq:app-profile} reduces to a monotone integer
bisection. Clamping at either bound is informative rather than benign: the upper bound binding
means the model cannot generate enough sparsity even when every variety is sourced from a single
origin, which is a rejection signal for the mechanism rather than a numerical nuisance.

\paragraph{Concentrating out comparative advantage.} Given
$(\alpha,\Omega_{L},\{\Omega_{s}\},\{A_{r}\})$, the area-level comparative advantage that
reproduces the observed sourcing shares \eqref{eq:app-gamma-area} is the solution of a
matrix-scaling problem. We solve it by the damped multiplicative iteration
\begin{equation}
T^{+}_{as}=T_{as}\left(\frac{\hat{\tilde\gamma}_{as}}{\tilde\gamma_{as}(T)}\right)^{\varsigma},
\qquad \varsigma\in(0,1],
\label{eq:app-sinkhorn}
\end{equation}
renormalizing to the reference area of each sector after every pass, and recomputing the
general-equilibrium expenditure weights that enter \eqref{eq:app-gamma-area}.

It is useful to separate the two things this iteration does, because only the first is covered
by a standard argument.

\emph{Step one: the scaling problem at fixed expenditure.} Hold the expenditure weights
$\omega_{rs}=X_{rs}/\sum_{r}X_{rs}$ of \eqref{eq:app-gamma-area} fixed. Then
$\tilde\gamma_{as}(T)$ is a ratio of positive linear forms in $T$: writing
$B_{ar}=\sum_{r'\in\mathcal{R}_{a}}(w_{r's}\tau_{r'rs})^{-\theta}$ for the (strictly positive)
bilateral access matrix, $\tilde\gamma_{as}(T)=\sum_{r}\omega_{rs}T_{as}B_{ar}/\sum_{a'}T_{a's}B_{a'r}$.
Recovering $T$ from target shares is then the classical matrix-scaling (Sinkhorn--Knopp)
problem, and \eqref{eq:app-sinkhorn} with $\varsigma=1$ is its standard iteration. Because
$B$ has full support --- every area can serve every buyer at finite cost --- the associated
positive linear map contracts the Hilbert projective metric with a modulus
$\kappa_{S}=\tanh(\Delta/4)<1$, where $\Delta$ is the projective diameter of the image cone
(Birkhoff--Hopf). Two consequences: the iteration converges geometrically from any strictly
positive start, and \textbf{the fixed point is unique up to the per-sector scale}, which is
exactly the gauge the reference normalization fixes. Given trade costs and expenditures, there
is one and only one vector of comparative advantages consistent with the observed area shares.

\emph{Step two: the general-equilibrium feedback.} Expenditure is not fixed. A change in $T$
moves the price indices $p_{rs}$, hence $c_{r}$, hence downstream sales $y_{r}$ and the weights
$\omega_{rs}$. The map actually iterated is the composition $\Psi=\mathcal{S}\circ\omega$, whose
Jacobian decomposes as
\begin{equation}
D\Psi = \underbrace{D\mathcal{S}\big|_{\omega\text{ fixed}}}_{\text{scaling channel}}
\;+\;\underbrace{J_{\mathrm{GE}}}_{\text{expenditure channel}},
\qquad
J_{\mathrm{GE}}\equiv D\Psi-D\mathcal{S}\big|_{\omega\text{ fixed}} .
\label{eq:app-jacobian-split}
\end{equation}
Birkhoff's theorem bounds the first term by $\kappa_{S}$ but says nothing about the second, so a
sufficient condition for $\Psi$ to be a contraction is
\begin{equation}
\kappa_{S}+\big\|J_{\mathrm{GE}}\big\|<1 .
\label{eq:app-contraction}
\end{equation}
The size of $J_{\mathrm{GE}}$ is governed by how far the technology is from Cobb--Douglas: the
expenditure weights are invariant to $T$ when $\lambda=\nu=1$, so $\|J_{\mathrm{GE}}\|$ vanishes
with $|1-\lambda|$ and $|1-\nu|$. Our calibration is not close to Cobb--Douglas, so
\eqref{eq:app-contraction} is an empirical question rather than an assumption, and we verify it
numerically at the estimated parameters by finite differences on the one-step map. We find
$\kappa_{S}\approx0.005$ and $\|J_{\mathrm{GE}}\|\approx0.015$, so the sufficient condition
\eqref{eq:app-contraction} holds with a wide margin ($\approx0.02\ll1$); the measured spectral
radius of the full map is $\approx0.012$, and convergence to $10^{-9}$ takes a handful of
iterations.%
% NOTE: these three numbers were measured by the Phase-0 inversion gate
% (test/test_ge_inversion.jl). Re-run it under --granular=true --ca_level=aa to
% refresh them before circulating, since the earlier measurement was taken on the
% ZE-level continuum configuration.
\ We further confirm that the inversion recovers a known $T$ from its own implied
shares to machine precision, and that ten dispersed starting values converge to the same fixed
point.

This is a local, numerically verified statement at the estimated parameters, not a global
theorem: \eqref{eq:app-contraction} is sufficient, not necessary, and we do not characterize the
region of the parameter space over which it holds. Should it fail at some candidate parameter
vector, the iteration is not used blindly --- non-convergence is detected and the candidate
retains its warm-start $T$, so the outer optimizer sees a well-defined objective throughout.

Substituting $T^{*}(\alpha,\Omega,A)$ into \eqref{eq:app-smm} collapses the outer search from
several hundred parameters to a handful and makes the sourcing-share block just-identified along
the profiled manifold, so that the extensive-margin coefficient is fit by $(\alpha,\Omega,A)$
alone. Note that just-identification here is at the \emph{area} level: within an area, the split
of sourcing across regions remains over-identified by the distance profile, which is the content
of Assumption~\ref{as:geo}.

\paragraph{Weighting matrix.} In the SMM framework noise arises from two sources: sampling
variability in the data and simulation error. We distinguish between moments for which
bootstrapping is feasible (sourcing shares, extensive-margin coefficients and count moments) and
those for which it is not (input--output coefficients and the downstream sales distribution). We
therefore approximate the variance--covariance matrix of the data moments, $\Sigma_{\mathrm{data}}$,
with a block matrix that accounts for within-group interactions; the block corresponding to
non-bootstrappable moments is set to zero. Because the same firms generate the sourcing shares,
the extensive-margin coefficients and the counts, these three blocks are estimated by a
\emph{single joint} bootstrap, which is what delivers the cross-covariance blocks the efficient
weight matrix requires. The resampling unit is the supplier firm, pooled over the geography of
active areas, with the denominator of $\bar G_{s}(0)$ held fixed --- so that a region with one
or two firms can draw zero and flip into the empty set, which is precisely the variance the
count moment needs.

\paragraph{Optimization procedure.} The estimator is computed using Particle Swarm Optimization
\citep{kennedy1995}, which explores the parameter space through parallel evaluation, with
adaptive inertia weights declining from $0.9$ to $0.4$ over iterations and cognitive and social
parameters set to $1.5$ and $2.5$. The estimation proceeds in three steps.

\begin{enumerate}[topsep=3pt,itemsep=3pt,label=\arabic*.]
\item Starting values are set as follows: $\Omega_{L}$ and $\{\Omega_{s}\}$ at their empirical
counterparts from input--output tables; $\{A_{r}\}$ proportional to $\pi_{r}^{1/\varepsilon}$;
$\alpha$ at the extensive-margin regression estimate divided by $\theta$; and $\{T_{as}\}$ at the
matrix-scaling solution \eqref{eq:app-sinkhorn} evaluated at that $\alpha$. We then run the
optimizer over the profiled parameter vector using the identity weighting matrix, refining with
an alternating block scheme in which each iteration sequentially optimizes $\{\Omega_{L},\Omega_{s}\}$,
then $\alpha$, then $\{A_{r}\}$, with a search radius progressively reduced across iterations.
This yields $\hat\mu_{1}$.
\item Using $\hat\mu_{1}$ we estimate the variance--covariance matrix of the simulated moments,
$\Sigma_{\mathrm{sim}}$, from repeated re-simulation with independent draws, and construct
$W=(\Sigma_{\mathrm{data}}+\Sigma_{\mathrm{sim}})^{-1}$.
\item We re-estimate the model with the updated weighting matrix, initializing at $\hat\mu_{1}$,
to obtain $\hat\mu_{2}$.
\end{enumerate}

\paragraph{Standard errors.} Let $G=\partial\tilde M/\partial\mu$ evaluated at $\hat\mu_{2}$ by
finite differences on common random numbers, and $\Omega=\Sigma_{\mathrm{data}}+\Sigma_{\mathrm{sim}}$.
We report the efficient variance $(G'WG)^{-1}$ alongside the sandwich form
$(G'WG)^{-1}G'W\Omega WG(G'WG)^{-1}$, which coincide when $W=\Omega^{-1}$; the ratio is a
diagnostic on the weighting matrix. Because $T$ is concentrated out through
\eqref{eq:app-sinkhorn} rather than searched, the Jacobian is taken along the profiled manifold
--- perturbing $\alpha$ and letting $T^{*}(\alpha)$ follow --- so the reported standard error on
$\alpha$ is a total derivative, and confidence intervals for $T$ are obtained by the delta
method, propagating both the uncertainty in $\hat\alpha$ and the sampling uncertainty in the
sourcing-share targets that \eqref{eq:app-sinkhorn} inverts.

\subsection{The price index, love of variety, and the normalization}
\label{app:model-price}

The one place where granularity meets the value block requires an explicit assumption. For
variety $\rho$ in market $(r,s)$, origin $r'$'s champion cost satisfies
$\Pr(c_{r'\rho s}(r)>x)=\exp(-T_{a(r')s}(w_{r's}\tau_{r'rs})^{-\theta}x^{\theta})$, so by
independence across origins the minimum $c^{*}_{\rho}$ is Weibull:
\begin{equation}
\Pr\big(c^{*}_{\rho}>x\big)=\exp\big(-\Phi_{rs}x^{\theta}\big),
\qquad
\mathbb{E}\big[(c^{*}_{\rho})^{k}\big]=\Phi_{rs}^{-k/\theta}\,\Gamma\Big(1+\tfrac{k}{\theta}\Big),
\quad k>-\theta .
\end{equation}
With the CES aggregate taken as a sum,
$\mathbb{E}[P_{rs}^{1-\nu_{s}}]=N_{s}\Gamma(1+\tfrac{1-\nu_{s}}{\theta})\Phi_{rs}^{-\frac{1-\nu_{s}}{\theta}}$
up to a constant, so the certainty-equivalent index
$\widetilde P_{rs}=(\mathbb{E}[P_{rs}^{1-\nu_{s}}])^{1/(1-\nu_{s})}$ satisfies
\begin{equation}
\widetilde P_{rs}=N_{s}^{\frac{1}{1-\nu_{s}}}\times
\underbrace{\Gamma\Big(1+\tfrac{1-\nu_{s}}{\theta}\Big)^{\frac{1}{1-\nu_{s}}}\Phi_{rs}^{-1/\theta}}_{N_{s}\text{-free}} .
\label{eq:app-loveofvariety}
\end{equation}
Since $1-\nu_{s}<0$ the factor is decreasing in $N_{s}$: more varieties lower the sector price
index. Rescaling $\widetilde T_{as}=T_{as}N_{s}^{\kappa}$ with $\kappa=\theta/(1-\nu_{s})$ maps
$\Phi_{rs}\mapsto N_{s}^{\kappa}\Phi_{rs}$ and cancels the two powers exactly, leaving
$\widetilde P_{rs}$ invariant to $N_{s}$. Equivalently, taking the CES aggregate as an
\emph{average} rather than a sum --- as in \eqref{eq:app-ces} --- implements the same
normalization. It is consistent because it acts on a gauge the sourcing shares do not see: the
share \eqref{eq:app-share} is a ratio of $T$'s within a sector, so the extensive margin, the
count moment and the within-area contrast are all unchanged. It is also a gauge the estimator
already fixes, by normalizing one reference area per sector.

Two caveats. First, $P_{rs}$ is a \emph{convex} transform of the average in \eqref{eq:app-ces},
so the realized index exceeds its certainty-equivalent value by Jensen's inequality, and the gap
closes only as fast as the average concentrates. At the calibration the regularity condition
$\theta>2(\nu_{s}-1)$ holds with little slack, and the gap is of the order of $6\%$ at
$N_{s}\approx10^{2}$, falling to $1\%$ at $N_{s}\approx10^{3}$. The estimator adopts the
certainty-equivalent reading --- downstream firms optimize against
$\mathbb{E}[P^{1-\nu_{s}}]^{1/(1-\nu_{s})}$, neglecting granular price randomness, the
general-equilibrium embedding of \citet{lmm2023}. Most of that $6\%$ would not have been
identifying variation in any case: a price-level shift is observationally equivalent to a
technology shift, so the sector technology terms absorb it and only the cross-sector pattern of
$N_{s}$ survives as signal. And by Lemma~\ref{lem:qfree} none of it touches $q_{r's}$ or the
extensive margin, so the identification of $\alpha$ and $N_{s}$ is unaffected either way.

Second, the assumption must not survive into counterfactuals. The elasticity
$\partial\ln\widetilde P_{rs}/\partial\ln N_{s}=1/(1-\nu_{s})<0$ is a genuine gain from variety
and a live propagation channel: a shock that removes suppliers in one region raises $P_{rs}$
downstream through exactly this elasticity, and that channel is absent from the continuum model.
Any counterfactual that moves the variety set must map back to the physical $T$ and use the
realized index.

\subsection{Proofs}
\label{app:model-proofs}

\paragraph{Equation \eqref{eq:app-champion}.} The number of firms in $(r',s,\rho)$ with
productivity above $z$ is Poisson with mean $T_{r's}z^{-\theta}$, so the probability that none
exceeds $z$ --- equivalently that the maximum is below $z$ --- is $\exp(-T_{r's}z^{-\theta})$,
the $\mathrm{Fr\acute{e}chet}(T_{r's},\theta)$ CDF. \hfill$\square$

\paragraph{Lemma~\ref{lem:fe}.} (i) With $\Pr(y=1)=1-\exp(-\exp(\phi_{g}+\beta x))$, a group in
which $y\equiv0$ contributes
$\ell_{g}=\sum_{i}\ln\exp(-e^{\phi_{g}+\beta x_{i}})=-e^{\phi_{g}}\sum_{i}e^{\beta x_{i}}$. This
is increasing in $-\phi_{g}$ with supremum $0$ as $\phi_{g}\to-\infty$, at which
$\partial\ell_{g}/\partial\beta=-e^{\phi_{g}}\sum_{i}x_{i}e^{\beta x_{i}}\to0$. The group
therefore drops from the concentrated likelihood.

(ii) By Frisch--Waugh--Lovell,
$\hat\beta=\sum_{g,i}\omega\tilde x\tilde y/\sum_{g,i}\omega\tilde x^{2}$ with $\tilde{\cdot}$
the weighted within-group demeaning. A group with $y\equiv0$ has $\tilde y\equiv0$, so it
contributes zero to the numerator and $\sum\omega\tilde x^{2}>0$ to the denominator. Splitting
the denominator between $\mathcal{A}^{+}$ and its complement gives
$\hat\beta=\varkappa\hat\beta^{+}$. \hfill$\square$

\paragraph{Proposition~\ref{prop:firm}.} Condition on the firm's productivity $z$; its delivered
cost to buyer $a$ is $c=w_{r's}\tau_{r'as}/z$. Each competitor $r''$ has champion cost
$w_{r''s}\tau_{r''as}/\varphi_{r''}$ with
$\varphi_{r''}\sim\mathrm{Fr\acute{e}chet}(T_{a(r'')s},\theta)$, so
$\Pr(w_{r''s}\tau_{r''as}/\varphi_{r''}>x)=\exp(-T_{a(r'')s}(w_{r''s}\tau_{r''as})^{-\theta}x^{\theta})$.
By independence across competitors the minimum competing cost $M$ satisfies
$\Pr(M>x)=\exp(-\Phi^{-r'}_{as}x^{\theta})$. The firm supplies iff $M>c$, giving
\eqref{eq:app-firm-win}. Taking the complement,
$\Pr(\text{no supply}\mid z)=1-\exp(-\exp(\eta))$ with
$\eta=\ln\Phi^{-r'}_{as}+\theta\ln(w_{r's}\tau_{r'as})-\theta\ln z$, and
$\tau_{r'as}=d_{r'a}^{\alpha}$ gives \eqref{eq:app-firm-cloglog}. No step involves the variety
count. \hfill$\square$

\paragraph{Lemma~\ref{lem:qfree}.} Origin $r'$ wins variety $\rho$ for buyer $r$ iff
$w_{r's}\tau_{r'rs}/\varphi_{r'\rho s}<\min_{r''\neq r'}w_{r''s}\tau_{r''rs}/\varphi_{r''\rho s}$,
an event measurable with respect to $\{\varphi_{\cdot\rho s}\}$, $w$ and $\tau$ alone; the
$\varphi$ have distribution governed by $(T_{\cdot s},\theta)$. The win-somewhere event is the
union over $r$ of those events, hence measurable with respect to the same objects. Neither the
variety count nor any price or expenditure aggregate enters. \hfill$\square$

\paragraph{Unbiasedness of \eqref{eq:app-Gclosed}.} Let $k$ be the number of the $m$ simulated
varieties won by $r'$, each won independently with probability $q$. Draw $n\le m$ of the $m$
varieties uniformly without replacement; conditional on $k$, the probability that all $n$ are
losses is $\binom{m-k}{n}/\binom{m}{n}$. Marginalizing over the draws, each of the $n$ selected
varieties is a loss independently with probability $1-q$, so the unconditional probability is
$(1-q)^{n}$. Hence $\mathbb{E}[\binom{m-k}{n}/\binom{m}{n}]=(1-q)^{n}$ exactly, with the convention
$\binom{a}{n}=0$ for $a<n$. \hfill$\square$
 from the main file. Requires amsmath, amssymb, amsthm,
% booktabs, enumitem and natbib. Add to the preamble if not already present:
%
%   \newtheorem{proposition}{Proposition}
%   \newtheorem{lemma}{Lemma}
%   \newtheorem{assumption}{Assumption}
%   \newtheorem{remark}{Remark}
%   \DeclareMathOperator*{\argmin}{arg\,min}
%
% Labels used by core_model.tex: app:model-derivation, app:model-foreign,
% app:model-calibration, app:model-prop, app:model-zeros, app:model-bounds.
% =====================================================================

\section{Details on the Structural Model}
\label{app:model}

\subsection{Derivation}
\label{app:model-derivation}

\paragraph{Setup.} The economy is divided into $R+1$ regions, populated by two types of firms:
producers of a final good and producers of intermediate goods used as inputs in the production
of the final good. The $(R+1)$th region represents the rest of the world. The analysis primarily
focuses on the matching of firms across domestic regions, controlling for competition from
foreign firms using their shares in total sales as reference. We assume that each domestic
region hosts at most one representative firm in the downstream sector, identified by its region
of origin. Intermediate goods are produced across $S$ sectors, with each region potentially
hosting many suppliers in each sector. We define a market $(s,r')$ as the collection of firms
located in region $r'$ that belong to upstream sector $s$.

Domestic regions are partitioned into attraction areas,
$\mathcal{R}=\bigsqcup_{a\in\mathcal{A}}\mathcal{R}_{a}$, where $\mathcal{R}_{a}$ is the set of
regions whose nearest downstream producer is the anchor of area $a$, and $a(r')$ denotes the
area of region $r'$.

\paragraph{Demand for final manufacturing products.} Firms in the downstream sector each own the
blueprint to produce a unique differentiated variety of the final good. The market for the final
good is characterized by monopolistic competition and iso-elastic demand functions:
\begin{equation}
q_{r}=\left(\frac{p_{r}}{P}\right)^{\varepsilon-1}\frac{E}{P}\,\delta_{r},
\end{equation}
where $q_{r}$ denotes the quantity demanded of the variety produced in region $r$ at price
$p_{r}$, given nominal consumption $E$ and the downstream sector's price index $P$. The
parameter $\varepsilon$ measures the price elasticity of sales, and $\delta_{r}$ is an exogenous
component of demand, varying across downstream producers and later used to simulate demand
shocks and their diffusion along production linkages. We abstract from foreign demand and assume
that the domestic market for final consumption goods is perfectly integrated.

\paragraph{Technology in the downstream sector.} Production of final goods requires the assembly
of intermediates, combined with labor, under a nested CES structure:
\begin{equation}
y_{r}=A_{r}\Big[\Omega_{L}^{\frac{1}{\lambda}}\ell_{r}^{\frac{\lambda-1}{\lambda}}
+(1-\Omega_{L})^{\frac{1}{\lambda}}\big(q^{I}_{r}\big)^{\frac{\lambda-1}{\lambda}}\Big]^{\frac{\lambda}{\lambda-1}},
\end{equation}
where $\ell_{r}$ is labor demand and $q^{I}_{r}$ the firm's intermediate input consumption. The
parameter $\Omega_{L}$ governs the share of labor in total production costs, while $\lambda$ is
the elasticity of substitution between labor and intermediates. Firm productivity $A_{r}$ is
exogenous. Intermediate consumption is itself a nested CES aggregator across sectors and, within
each sector, across varieties:
\begin{equation}
q^{I}_{r}=\Big[\sum_{s=1}^{S}\Omega_{s}^{\frac{1}{\nu}}q_{rs}^{\frac{\nu-1}{\nu}}\Big]^{\frac{\nu}{\nu-1}},
\qquad
q_{rs}=\Big[\frac{1}{N_{s}}\sum_{\rho=1}^{N_{s}}q_{rs}(\rho)^{\frac{\nu_{s}-1}{\nu_{s}}}\Big]^{\frac{\nu_{s}}{\nu_{s}-1}},
\label{eq:app-ces}
\end{equation}
where $\Omega_{s}$ governs the share of sector $s$ in intermediate consumption, $\nu$ is the
elasticity of substitution across sectors and $\nu_{s}$ that across varieties within a sector.

The only change relative to the continuum formulation is the inner aggregator: the integral over
$j\in[0,1]$ is replaced by an average over the $N_{s}$ varieties. Writing it as an
\emph{average} rather than a sum is a normalization of the love-of-variety level, discussed and
justified in Appendix~\ref{app:model-price}; it leaves every sourcing share, and hence the
entire extensive margin, unchanged.

Given the nested CES structure, the nominal demand for intermediate products is
\begin{equation}
x_{rs}(\rho)=\Omega_{s}(1-\Omega_{L})
\left(\frac{p_{rs}(\rho)}{p_{rs}}\right)^{1-\nu_{s}}
\left(\frac{p_{rs}}{p_{r}}\right)^{1-\nu}
\left(\frac{p_{r}}{c_{r}}\right)^{1-\lambda}\frac{c_{r}y_{r}}{A_{r}},
\end{equation}
with $p_{rs}\equiv\big[\tfrac{1}{N_{s}}\sum_{\rho}p_{rs}(\rho)^{1-\nu_{s}}\big]^{\frac{1}{1-\nu_{s}}}$
the cost index for intermediates in sector $s$,
$p_{r}\equiv\big[\sum_{s}\Omega_{s}p_{rs}^{1-\nu}\big]^{\frac{1}{1-\nu}}$ the cost index for
intermediates, and
$c_{r}\equiv\big[\Omega_{L}w_{r}^{1-\lambda}+(1-\Omega_{L})p_{r}^{1-\lambda}\big]^{\frac{1}{1-\lambda}}$
the unit cost of the firm. In equilibrium $y_{r}=q_{r}$.

\paragraph{Technology in upstream sectors.} Intermediate goods are produced with labor using a
constant-returns-to-scale technology. Wages in sector $s$ and region $r'$ are denoted $w_{r's}$
and taken as given. As in Eaton and Kortum (2002) we assume that varieties produced within and
across regions are perfect substitutes and that productivity draws are independent. The
departure from the continuum case is that each sector contains $N_{s}$ varieties.

\begin{assumption}[Granularity]
\label{as:granular}
Each upstream sector $s$ consists of $N_{s}\in\mathbb{N}$ horizontally differentiated varieties
$\rho=1,\dots,N_{s}$. Within each variety $\rho$ and each region $r'$ there is a continuum of
firms whose productivities above $z$ form a Poisson process with mean $T_{r's}z^{-\theta}$.
Draws are independent across varieties, regions and sectors.
\end{assumption}

The continuum of firms is what keeps the model tractable. The number of firms above $z$ in
$(r',s,\rho)$ is Poisson with mean $T_{r's}z^{-\theta}$, so the most productive one satisfies
\begin{equation}
\Pr\big(\varphi_{r'\rho s}\le z\big)=\Pr(\text{no firm above }z)=e^{-T_{r's}z^{-\theta}},
\label{eq:app-champion}
\end{equation}
that is, $\varphi_{r'\rho s}$ is exactly $\mathrm{Fr\acute{e}chet}(T_{r's},\theta)$. One
Fr\'echet draw per (region, variety) pair \emph{is} the champion of a continuum of firms, not an
approximation to it. Since the iceberg cost scales all firms of region $r'$ identically, that
champion is the unique region-$r'$ firm that ever serves any buyer sourcing $\rho$ from $r'$:
\textbf{a firm is a (region, variety) champion}, with delivered cost
$c_{r'\rho s}(r)=w_{r's}\tau_{r'rs}/\varphi_{r'\rho s}$. This is the ``fixed number of
independent Ricardian draws'' form of granular Eaton--Kortum \citep[cf.][]{eks2012,bejk2003}.

\begin{assumption}[Area-level comparative advantage]
\label{as:geo}
Comparative advantage is area-specific: $T_{r's}=T_{a(r')s}$ for every $r'\in\mathcal{R}_{a}$,
with one parameter per (sector, area). Productivity draws remain independent across regions ---
two regions in the same area share the Fr\'echet scale, not the realization of their champions.
\end{assumption}

Assumption~\ref{as:geo} does not make regions within an area interchangeable. Each runs its own
Ricardian competition for each variety, and two regions at different distances from a buyer have
different win probabilities. What it removes is unobserved technological heterogeneity within an
area, which is what allows the model to explain the same variation the extensive-margin
regression uses.

\paragraph{Equilibrium.} The resolution follows the same steps as in Eaton and Kortum (2002).
Given the Fr\'echet assumption, the minimum price at which suppliers from $(r',s)$ can serve
market $r$ is distributed
\begin{equation}
G_{r'rs}(p)=\Pr\Big(\frac{w_{r's}\tau_{r'rs}}{\varphi_{r'\rho s}}\le p\Big)
=1-e^{-\left[T_{a(r')s}(w_{r's}\tau_{r'rs})^{-\theta}\right]p^{\theta}},
\end{equation}
so the distribution of input prices for downstream firms in region $r$ is
\begin{equation}
G_{rs}(p)=1-\prod_{r'=1}^{R+1}\big[1-G_{r'rs}(p)\big]=1-e^{-\Phi_{rs}p^{\theta}},
\qquad
\Phi_{rs}=\sum_{a'}T_{a's}\sum_{r'\in\mathcal{R}_{a'}}(w_{r's}\tau_{r'rs})^{-\theta}.
\label{eq:app-Phi}
\end{equation}
The probability that a downstream firm in location $r$ picks a supplier in market $(r',s)$ is
thus
\begin{equation}
\gamma_{r'rs}=\frac{T_{a(r')s}(w_{r's}\tau_{r'rs})^{-\theta}}{\Phi_{rs}} .
\label{eq:app-share}
\end{equation}
Taking the perspective of the supplier, the probability of being the lowest-cost supplier in
market $(r,s)$, given location $r'$ and productivity $z$, is
\begin{equation}
\rho_{r'rs}(z)=e^{-\Phi_{rs}(w_{r's}\tau_{r'rs})^{\theta}z^{-\theta}}
=e^{-T_{a(r')s}\gamma^{-1}_{r'rs}z^{-\theta}} .
\label{eq:app-rho}
\end{equation}
Because the conditional distribution of prices is not origin-specific,
$G_{r'rs}(p)=G_{rs}(p)$, and $\gamma_{r'rs}$ is also the expected share of intermediates from
$r'$ purchased by a downstream firm in $r$. The expected cost indices are
\begin{equation}
p_{rs}=\Big[\Gamma\Big(\frac{\theta+1-\nu_{s}}{\theta}\Big)\Big]^{\frac{1}{1-\nu_{s}}}\Phi_{rs}^{-1/\theta},
\qquad
p_{r}=\Big[\sum_{s}\Omega_{s}p_{rs}^{1-\nu}\Big]^{\frac{1}{1-\nu}},
\qquad
c_{r}=\Big[\Omega_{L}w_{r}^{1-\lambda}+(1-\Omega_{L})p_{r}^{1-\lambda}\Big]^{\frac{1}{1-\lambda}},
\end{equation}
where the first expression is the certainty-equivalent index of Appendix~\ref{app:model-price}.
Bilateral flows are
\begin{equation}
X_{r'rs}=\gamma_{r'rs}\,x_{rs}
=\gamma_{r'rs}\,\Omega_{s}(1-\Omega_{L})
\left(\frac{p_{rs}}{p_{r}}\right)^{1-\nu}\left(\frac{p_{r}}{c_{r}}\right)^{1-\lambda}\frac{c_{r}y_{r}}{A_{r}} .
\end{equation}
Since we do not observe these bilateral flows directly, we cannot use this equation to estimate
the model's structural parameters and rely instead on simulated moments.

\subsection{The extensive margin and the two classes of zeros}
\label{app:model-zeros}

Across the $N_{s}$ varieties the win events are i.i.d., so with $K_{r'rs}$ the number of
sector-$s$ varieties buyer $r$ sources from $r'$,
\begin{equation}
K_{r'rs}\sim\mathrm{Bin}\big(N_{s},\gamma_{r'rs}\big),
\qquad
\Pr(K_{r'rs}=0)=(1-\gamma_{r'rs})^{N_{s}}>0 .
\end{equation}
Under a continuum this probability is $\lim_{N\to\infty}(1-\gamma_{r'rs})^{N}=0$ for every
origin with $T_{r's}>0$: the extensive margin is degenerate and the observed zeros are not a
model outcome. Finiteness is exactly what restores it.

\begin{remark}[Dependence across buyers]
\label{rem:dependence}
For a given variety the champion costs are common across buyers; only $\tau_{r'rs}$ differs. Win
events are therefore positively dependent across buyers, and the union probability $q_{r's}$ is
not the product of the per-market marginals. The per-market marginal above is nonetheless exact,
and $q_{r's}$ is evaluated by simulation. This matters because the survey records whether a firm
supplies the industry, not which buyer it serves: the empirical outcome is the union, and the
model counterpart is computed as the same union.
\end{remark}

Assumption~\ref{as:geo} sorts observed zeros into two economically distinct classes, and the
distinction determines what is estimated.

\begin{table}[htbp]
\centering
\small
\begin{tabular}{@{}p{0.21\textwidth}p{0.24\textwidth}p{0.19\textwidth}p{0.28\textwidth}@{}}
\toprule
\textbf{Zero} & \textbf{Treatment} & \textbf{Probability} & \textbf{Role}\\
\midrule
Whole attraction area $a$ has no supplier in $s$
& \emph{Structural}: $T_{as}=0$, $(s,a)$ outside the active set
& $1$ by construction
& None --- absorbed by the $a\times s$ fixed effect (Lemma~\ref{lem:fe}); no parameter estimated\\[4pt]
Region $r'$ inside an active area has no supplier
& \emph{Endogenous}: inherits $T_{as}>0$, own distance $d_{r'}$
& $(1-q_{r's})^{N_{s}}$
& \textbf{The identifying variation}: disciplined by distance alone\\
\bottomrule
\end{tabular}
\caption{The two classes of extensive-margin zeros. The second row is exactly the set of regions
that host firms but no supplier. Under region-level comparative advantage these had to be
assigned $T\equiv0$, making their zeros exogenous; under Assumption~\ref{as:geo} they inherit
their area's comparative advantage and their zeros become model predictions.}
\label{tab:app-hierarchy}
\end{table}

The reason the fixed effect dictates this split is the following.

\begin{lemma}[All-zero groups carry no information about the distance slope]
\label{lem:fe}
Index observations by firm $i$ within fixed-effect group $g=(a,s)$, with outcome $y_{i}$,
regressor $x_{i}$, weights $\omega_{i}>0$ and group effect $\phi_{g}$.
\begin{enumerate}[nosep,topsep=2pt,label=(\roman*)]
\item \emph{Complementary log-log.} The log-likelihood of a group in which $y\equiv0$ is
$\ell_{g}=-e^{\phi_{g}}\sum_{i}e^{\beta x_{i}}$, whose supremum over $\phi_{g}$ is attained as
$\phi_{g}\to-\infty$ with $\ell_{g}\to0$ and $\partial\ell_{g}/\partial\beta\to0$: the group
drops out of the concentrated likelihood and contributes exactly nothing to $\hat\beta$.
\item \emph{Linear probability model.} By Frisch--Waugh--Lovell an all-zero group has
$\tilde y\equiv0$, contributing zero to the numerator and a positive amount to the denominator,
so $\hat\beta=\varkappa\hat\beta^{+}$ with
$\varkappa=\big(\sum_{a\in\mathcal{A}^{+}}\sum\omega\tilde x^{2}\big)/\big(\sum_{a}\sum\omega\tilde x^{2}\big)\in(0,1]$,
where $\mathcal{A}^{+}$ is the set of areas hosting at least one supplier.
\end{enumerate}
\end{lemma}

Under the cloglog link --- the link the model implies --- $\alpha$ is therefore identified
exclusively from variation \emph{within} attraction areas that host at least one supplier. Under
a linear probability model the empty areas enter only as a mechanical dilution of the same
within-area signal. Either way, an area with no supplier at all carries no information about the
distance slope, which is why the model does not attempt to explain it.

\subsection{The firm-level reduced form}
\label{app:model-prop}

The object the estimation matches is the union event \eqref{eq:q}: whether a firm supplies the
downstream industry at all. That probability has no tractable closed form. The single-buyer case
does, and it is what gives the empirical cloglog its structural interpretation; we state it
first and then say precisely what survives when there are several buyers.

\begin{proposition}[Firm-level reduced form, single buyer]
\label{prop:firm}
Let a firm be a (region, variety) champion with productivity $z$ located in $r'$, and suppose it
faces a single downstream buyer, its area's anchor $a$. Writing
$\Phi^{-r'}_{as}=\sum_{r''\neq r'}T_{a(r'')s}(w_{r''s}\tau_{r''as})^{-\theta}$ for the
competitors' cost scale,
\begin{equation}
\Pr\big(\text{the firm supplies}\mid z\big)
=\exp\Big(-\Phi^{-r'}_{as}\big(w_{r's}\tau_{r'as}\big)^{\theta}z^{-\theta}\Big),
\label{eq:app-firm-win}
\end{equation}
so the probability that it does \emph{not} supply is a complementary log-log function
\begin{equation}
\Pr\big(\text{no supply}\mid z\big)=1-\exp\big(-\exp(\eta)\big),
\qquad
\eta=\underbrace{\ln\Phi^{-r'}_{as}}_{\text{area}\times\text{sector FE}}
+\theta\ln w_{r's}
+\theta\alpha\,\ln d_{r'a}
-\theta\,\ln z .
\label{eq:app-firm-cloglog}
\end{equation}
The variety count $N_{s}$ does not appear.
\end{proposition}

Three consequences, and they are the backbone of the estimation.

\begin{enumerate}[nosep,topsep=3pt,label=(\alph*)]
\item \textbf{The cloglog is the model's own auxiliary model.} A cloglog of the ``does not
supply'' indicator on log distance and log own productivity, with an area$\times$sector fixed
effect, is exactly \eqref{eq:app-firm-cloglog} when a firm faces one buyer, with distance
coefficient $+\theta\alpha$ and productivity coefficient $-\theta$.
\item \textbf{$\alpha$ and $N_{s}$ separate.} $N_{s}$ is absent from
\eqref{eq:app-firm-cloglog}: conditioning on the firm removes it. The regression identifies
$\alpha$; the count level identifies $N_{s}$; there is no residual coupling.
\item \textbf{The productivity coefficient is a free over-identifying test.} It equals
$-\theta$. Extending the moment vector to include it would let $\theta$ be estimated rather than
calibrated.
\end{enumerate}

\paragraph{Several buyers: what survives.} Our surveys record supply to the industry, not to a
named buyer, so the empirical outcome is the union event $\bigcup_{r}W_{r}$, where $W_{r}$
denotes winning buyer $r$. Its probability is \emph{not} $1-\prod_{r}[1-\rho_{r'rs}(z)]$: that
expression would require the $W_{r}$ to be independent, and they are not.

The dependence is positive and structural: conditional on $z$, the firm's delivered cost to
every buyer is the same number scaled by $\tau_{r'rs}$, so a firm cheap enough to beat the
competition in one market tends to beat it everywhere. The exact object is the
inclusion--exclusion sum
\begin{equation}
\Pr\Big(\bigcup_{r}W_{r}\,\Big|\,z\Big)
=\sum_{\emptyset\neq\mathcal{S}\subseteq\{1,\dots,R\}}(-1)^{|\mathcal{S}|+1}
\exp\Big(-z^{-\theta}\max_{r\in\mathcal{S}}\big\{\Phi^{-r'}_{rs}(w_{r's}\tau_{r'rs})^{\theta}\big\}\Big),
\label{eq:app-inclexcl}
\end{equation}
where the maximum appears because $\cap_{r\in\mathcal{S}}W_{r}$ requires the firm to beat the
\emph{tightest} of the thresholds it faces. Equation~\eqref{eq:app-inclexcl} has $2^{R}-1$
terms. Many are redundant --- if the threshold for buyer $r$ dominates that for $r''$ then
$W_{r}\subseteq W_{r''}$ and the pair collapses --- and pruning the dominated buyers reduces the
sum to the non-dominated set, but that set is still large enough in our geography that
evaluating \eqref{eq:app-inclexcl} inside an optimizer is not practical. The independence
approximation $1-\prod_{r}(1-\rho_{r'rs}(z))$, which is what a closed-form estimator would have
to use, is a strict \emph{under}-statement of \eqref{eq:app-inclexcl} by positive dependence, and
the error is largest exactly where the extensive margin has content, at high $z$.

This is the reason the model is estimated by simulated rather than analytical method of moments.
Simulation evaluates the union exactly --- each variety's champion costs are drawn once and
compared against every buyer, so the dependence is reproduced by construction --- at the cost of
simulation noise, which the weighting matrix accounts for. We therefore treat
\eqref{eq:app-firm-cloglog} as an auxiliary model rather than an estimating equation: the same
cloglog is run on simulated and observed data, with the union as the outcome on both sides, and
the coefficients are matched. This is consistent for $\alpha$ whether or not the link is exact,
and it is why the distance coefficient enters the moment vector rather than being inverted
analytically.

The outcome convention is dictated by the unit of observation rather than chosen. Simulating the
exact object at $\theta=1$, $\alpha=0.30$, with firms as (region $\times$ variety) champions,
the ``does not supply'' outcome recovers $\hat\beta_{\ln d}=+0.327$ and
$\hat\beta_{\ln z}=-1.051$ against truths $\theta\alpha=0.30$ and $-\theta=-1.0$; the reversed
outcome recovers neither ($-0.426$ and $+1.452$).

\subsection{The count moment and bounds on the variety count}
\label{app:model-bounds}

The zeros are a cell-level object. With $\mathcal{R}^{+}_{s}$ the regions inside active
attraction areas, the empirical moment is the share of those regions hosting at most $n$
suppliers,
\begin{equation}
\bar G_{s}(n)=\frac{1}{|\mathcal{R}^{+}_{s}|}
\#\big\{r'\in\mathcal{R}^{+}_{s}:K_{r's}\le n\big\},
\qquad
\mathbb{E}\big[\bar G_{s}(0)\big]=\frac{1}{|\mathcal{R}^{+}_{s}|}\sum_{r'}(1-q_{r's})^{N_{s}} .
\label{eq:app-count}
\end{equation}
Writing $\mu_{r's}=N_{s}q_{r's}$ and using the Poisson approximation,
$\partial\bar G_{s}(K)/\partial N_{s}
=-|\mathcal{R}^{+}_{s}|^{-1}\sum_{r'}q_{r's}e^{-\mu_{r's}}\mu_{r's}^{K}/K!$, which peaks at
$K\approx\mu_{r's}$ and vanishes as $\bar G_{s}(K)\to1$: saturated thresholds carry no
identifying variation and a poorly estimated variance, so we use the empty-region share
$\bar G_{s}(0)$, which is both the most interpretable threshold and typically near-optimal.

Because the iceberg cost never changes which firm is cheapest within a region, each variety
generates one distinct supplying firm per winning origin. So if $N^{\mathrm{obs}}_{s}$ is the
number of distinct upstream firms observed to supply the downstream industry,
$N^{\mathrm{obs}}_{s}=\sum_{\rho}\#\{\text{origins winning }\rho\text{ somewhere}\}\in[N_{s},N_{s}R^{D}]$,
where $R^{D}$ is the number of downstream buyers, hence
\begin{equation}
\frac{N^{\mathrm{obs}}_{s}}{R^{D}}\;\le\;N_{s}\;\le\;N^{\mathrm{obs}}_{s} .
\label{eq:app-bounds}
\end{equation}
The bounds are tight: the upper when every variety is globally sourced from a single origin, the
lower when each of the $R^{D}$ buyers sources it from a different origin. Since
$\mathbb{E}[\sum_{r'}K_{r's}]=N_{s}\sum_{r'}q_{r's}$, matching the observed firm count gives a second,
closed-form estimate $N^{\mathrm{count}}_{s}=N^{\mathrm{obs}}_{s}/\sum_{r'}q_{r's}$ that lies
inside \eqref{eq:app-bounds} by construction --- a free over-identifying check.

\paragraph{Degrees of freedom.} Fix a sector, let $A^{+}=|\mathcal{A}^{+}_{s}|$ and let
$R^{\mathrm{act}}_{a}$ be the number of regions in area $a$ with a strictly positive observed
share. The $T$ block has $A^{+}-1$ free elements after the reference normalization, against
$\sum_{a}R^{\mathrm{act}}_{a}$ share moments, so Assumption~\ref{as:geo} converts what were free
parameters into
\begin{equation}
\sum_{a\in\mathcal{A}^{+}}\big(R^{\mathrm{act}}_{a}-1\big)
\quad\text{over-identifying restrictions,}
\end{equation}
one per active region beyond the first in each area. Each is a within-area statement: given the
area aggregate, the split across its regions must follow the distance profile. Under
region-level comparative advantage this count is zero.

\subsection{Structural estimation: the moments}
\label{app:model-moments}

We use six sets of theoretical moments, observed in the data using a combination of IO tables
and micro-level data.

\begin{enumerate}[topsep=3pt,itemsep=3pt,label=\arabic*.]
\item \textbf{Aggregate labor share at the level of the industry:}
\begin{equation}
\tilde\Omega_{L}\equiv\frac{\sum_{r}w_{r}\ell_{r}}{\sum_{r}c_{r}A_{r}^{-1}y_{r}}
=\sum_{r}\frac{c_{r}A_{r}^{-1}y_{r}}{\sum_{r}c_{r}A_{r}^{-1}y_{r}}\,\Omega_{L}
\left(\frac{w_{r}}{c_{r}}\right)^{1-\lambda}.
\end{equation}
In the model it is a sales-weighted average of firm-level labor shares. In the data we use the
share of labor costs in total costs from IO tables.

\item \textbf{Aggregate share of sector $s$ in the industry's input purchases:}
\begin{equation}
\tilde\Omega_{s}\equiv\frac{\sum_{r}p_{rs}q_{rs}}{\sum_{r}\sum_{s}p_{rs}q_{rs}}
=\sum_{r}\frac{p_{r}q^{I}_{r}}{\sum_{r}p_{r}q^{I}_{r}}\,\Omega_{s}
\left(\frac{p_{rs}}{p_{r}}\right)^{1-\nu}.
\end{equation}
In the data we use technical coefficients for the downstream industry from IO tables.

\item \textbf{The share of each region in downstream sales:}
\begin{equation}
\pi_{r}\equiv\frac{c_{r}A_{r}^{-1}y_{r}}{\sum_{r}c_{r}A_{r}^{-1}y_{r}}
=\left(\frac{p_{r}}{P}\right)^{-\varepsilon}
=\frac{(c_{r}A_{r}^{-1})^{-\varepsilon}}{\sum_{r}(c_{r}A_{r}^{-1})^{-\varepsilon}},
\end{equation}
recovered from sales data on the population of downstream producers.

\item \textbf{The extensive-margin distance coefficient.} The cloglog of
\eqref{eq:app-firm-cloglog}, estimated on the simulated firm population with the same
specification as Equation~(1) in the data: outcome ``does not supply'', regressors log distance
to the area anchor and log own productivity, area$\times$sector fixed effects.

\item \textbf{Area-aggregated sourcing shares.} The share of the downstream industry's purchases
of sector $s$ sourced from area $a$:
\begin{equation}
\tilde\gamma_{as}\equiv\sum_{r'\in\mathcal{R}_{a}}\frac{\sum_{r}X_{r'rs}}{\sum_{r}\sum_{r''}X_{r''rs}}
=\sum_{r'\in\mathcal{R}_{a}}\sum_{r}\frac{\sum_{r''}X_{r''rs}}{\sum_{r''}\sum_{r}X_{r''rs}}\,\gamma_{r'rs}.
\label{eq:app-gamma-area}
\end{equation}
In the data we reconstitute firm-level sales to the downstream industry using sales-share data,
$x_{i(r',s)d}=\mathbf{1}(\mathrm{Sup}_{i(r',s)})\,a^{D}_{d\,i(r',s)}\,x_{i(s,r')}$, sum over all
firms in the area and normalize by the sum over all firms from sector $s$. The aggregation runs
over \emph{every} region of the area, including those hosting no supplier, on both the model and
the data side.

\item \textbf{The count moment} $\bar G_{s}(0)$ of \eqref{eq:app-count}: the share of regions
inside active attraction areas that host no supplier of sector $s$.
\end{enumerate}

Moments 1--3 and 5 are the same objects as in the continuum version of the model, with 5 now
aggregated to the area. Moments 4 and 6 are what the granular structure adds.

\paragraph{Notes on foreign parameters.}
\label{app:model-foreign}
We do not have the data to estimate parameters relating to foreign firms. Instead we take the
share of foreign inputs in sourcing as an exogenous parameter used in the simulated data, so the
vector of parameters to be estimated excludes region $R+1$. Moments 1--4 and 6 do not involve
$R+1$. For moment 5, with $w_{R+1,s}$ the share of foreign inputs in sector $s$'s input
purchases,
\begin{equation}
\tilde\gamma_{as}=(1-w_{R+1,s})\sum_{r'\in\mathcal{R}_{a}}\sum_{r}
\frac{\sum_{r''\neq R+1}X_{r''rs}}{\sum_{r''\neq R+1}\sum_{r}X_{r''rs}}\,\tilde\gamma_{r'rs},
\qquad
\tilde\gamma_{r'rs}=\frac{T_{a(r')s}(w_{r's}\tau_{r'rs})^{-\theta}}{\tilde\Phi_{rs}},
\end{equation}
with $\tilde\Phi_{rs}=\sum_{r'=1}^{R}T_{a(r')s}(w_{r's}\tau_{r'rs})^{-\theta}$. Foreign
competition still enters the win probabilities through $\Phi_{rs}$ in \eqref{eq:app-Phi}; what
is not estimated is foreign comparative advantage.

\paragraph{Calibrated parameters.}
\label{app:model-calibration}
The remaining parameters $\{\varepsilon,\theta,\lambda,\nu,\{\nu_{s}\},\{w_{r's}\}\}$ are
calibrated. We use the estimates from Fontagn\'e et al.\ (2022) for $\varepsilon$: the sales
elasticity for motor vehicles and aerospace is respectively $-9$ and $-16$. The shape $\theta$
of the Fr\'echet distribution is calibrated following Antr\`as et al.\ (2017). Elasticities of
substitution in production follow the literature, notably Baqaee and Farhi (2019): $\lambda=0.5$
for the elasticity between labor and intermediate inputs and $\nu=0.9$ across sectors. Local
wages are calibrated using labor cost information from employer--employee data in 2019
(DADS/BTS). Note that Proposition~\ref{prop:firm}(c) makes $\theta$ over-identified: the
coefficient on log own productivity in moment 4 should equal $-\theta$, and we report it.

\subsection{Estimation}
\label{app:model-estimation}

\paragraph{Simulation algorithm.} Given the calibrated elasticities and a candidate parameter
vector, we simulate the productivity champions of every (region, variety) pair, solve for all
bilateral trade flows --- from each potential supplier to each downstream region --- and
construct the simulated moments. All six blocks are computed from a \emph{single} Ricardian
solve on one set of draws. The SMM estimator is
\begin{equation}
\hat\mu=\argmin_{\mu}\;\big(\tilde M(\mu)-\hat M\big)'W\big(\tilde M(\mu)-\hat M\big),
\label{eq:app-smm}
\end{equation}
where $\tilde M$ (resp.\ $\hat M$) denotes the vector of theoretical (resp.\ empirical) moments
and $W$ is a weighting matrix. Optimal weights assign lower importance to imprecisely estimated
moments.

The variety count never enters the simulation. It is confined to two objects, both closed form,
and that confinement is the content of Proposition~\ref{prop:firm} together with:

\begin{lemma}[$q_{r's}$ is free of $N_{s}$]
\label{lem:qfree}
Winning a variety is the comparison
$w_{r's}\tau_{r'rs}/\varphi_{r'\rho s}<\min_{r''\neq r'}w_{r''s}\tau_{r''rs}/\varphi_{r''\rho s}$,
so $q_{r's}=q_{r's}(T_{\cdot s},w,\tau,\theta)$ and nothing else. The variety count, the price
index, downstream costs and expenditures do not appear.
\end{lemma}

The economics is that winning is a statement about \emph{relative costs}, whereas $N_{s}$ and
$P_{rs}$ are aggregators built after the winners are known: expenditure allocates value across
winners, it does not choose them. The one condition that would break the lemma is endogenous
upstream wages, since $w$ enters the comparison directly; wages are taken as data.

Together the two results say that $N_{s}$ reaches the moment vector only through
$\bar G_{s}(0)$ and through $\widehat N_{s}$ itself, and that both are functions of $q_{r's}$
alone. Writing $k_{r's}$ for the number of the $m$ simulated varieties won somewhere by region
$r'$ --- by Lemma~\ref{lem:qfree} an estimate of $q_{r's}$ valid whatever the variety count ---
the count moment is evaluated by the unbiased combinatorial estimator
\begin{equation}
\bar G_{s}(n)=\frac{1}{|\mathcal{R}^{+}_{s}|}\sum_{r'\in\mathcal{R}^{+}_{s}}
\frac{\binom{m-k_{r's}}{n}}{\binom{m}{n}},
\qquad
\mathbb{E}\left[\frac{\binom{m-k}{n}}{\binom{m}{n}}\right]=(1-q)^{n}\ \ \text{for } n\le m,
\label{eq:app-Gclosed}
\end{equation}
which is the probability that $n$ varieties drawn without replacement from the $m$ simulated
ones were all lost. We use \eqref{eq:app-Gclosed} rather than the plug-in $(1-\hat q_{r's})^{n}$
because the latter is convex in $\hat q$ and therefore biased upward. Cross-cell dependence ---
two regions in a sector face the same variety draws --- does not affect either expression, since
an expectation of an average is the average of expectations whatever the dependence. No economy
of size $N_{s}$ is ever drawn and no replication is needed.

\begin{remark}[What is given up]
\label{rem:binding}
Moment 4 is then the coefficient of a regression run on the full simulated sample, i.e.\ the
continuum limit $\beta(\mathbb{E}[y])$ rather than $\mathbb{E}[\hat\beta(N_{s})]$. The finite-sample bias of the
auxiliary cloglog --- which an average of $\hat\beta$ over realized economies of size $N_{s}$
would have cancelled against the same bias in the data --- is therefore not differenced out. It
is a measurable quantity, of the order of $2$--$4\%$ of $\beta$ at data-consistent group sizes,
and we report it alongside $\hat\alpha$. The alternative is not free: it would treat simulation
noise on one block by replication while every other block relies on the accuracy of the draw
design, and it would sit oddly beside Appendix~\ref{app:model-price}, where the value block is
already evaluated at the certainty-equivalent index rather than on a realized economy.
\end{remark}

\paragraph{Concentrating out the variety count.} $N_{s}$ is estimated by the same weighted
objective as every other parameter, concentrated rather than searched jointly:
\begin{equation}
\widehat N(\mu_{-N})=\argmin_{N\in\prod_{s}[N^{\mathrm{lo}}_{s},N^{\mathrm{hi}}_{s}]\cap\mathbb{N}}
\;r(\mu_{-N},N)'W\,r(\mu_{-N},N),
\label{eq:app-profile}
\end{equation}
with the bounds of \eqref{eq:app-bounds}. This is a profiled extremum estimator: $\widehat N_{s}$
is an integer at every evaluation, so integrality is never relaxed and never rounded, and by
Lemma~\ref{lem:qfree} the inner minimization needs no re-simulation. Since $\bar G_{s}(0;n)$ is
closed form and decreasing in $n$, \eqref{eq:app-profile} reduces to a monotone integer
bisection. Clamping at either bound is informative rather than benign: the upper bound binding
means the model cannot generate enough sparsity even when every variety is sourced from a single
origin, which is a rejection signal for the mechanism rather than a numerical nuisance.

\paragraph{Concentrating out comparative advantage.} Given
$(\alpha,\Omega_{L},\{\Omega_{s}\},\{A_{r}\})$, the area-level comparative advantage that
reproduces the observed sourcing shares \eqref{eq:app-gamma-area} is the solution of a
matrix-scaling problem. We solve it by the damped multiplicative iteration
\begin{equation}
T^{+}_{as}=T_{as}\left(\frac{\hat{\tilde\gamma}_{as}}{\tilde\gamma_{as}(T)}\right)^{\varsigma},
\qquad \varsigma\in(0,1],
\label{eq:app-sinkhorn}
\end{equation}
renormalizing to the reference area of each sector after every pass, and recomputing the
general-equilibrium expenditure weights that enter \eqref{eq:app-gamma-area}.

It is useful to separate the two things this iteration does, because only the first is covered
by a standard argument.

\emph{Step one: the scaling problem at fixed expenditure.} Hold the expenditure weights
$\omega_{rs}=X_{rs}/\sum_{r}X_{rs}$ of \eqref{eq:app-gamma-area} fixed. Then
$\tilde\gamma_{as}(T)$ is a ratio of positive linear forms in $T$: writing
$B_{ar}=\sum_{r'\in\mathcal{R}_{a}}(w_{r's}\tau_{r'rs})^{-\theta}$ for the (strictly positive)
bilateral access matrix, $\tilde\gamma_{as}(T)=\sum_{r}\omega_{rs}T_{as}B_{ar}/\sum_{a'}T_{a's}B_{a'r}$.
Recovering $T$ from target shares is then the classical matrix-scaling (Sinkhorn--Knopp)
problem, and \eqref{eq:app-sinkhorn} with $\varsigma=1$ is its standard iteration. Because
$B$ has full support --- every area can serve every buyer at finite cost --- the associated
positive linear map contracts the Hilbert projective metric with a modulus
$\kappa_{S}=\tanh(\Delta/4)<1$, where $\Delta$ is the projective diameter of the image cone
(Birkhoff--Hopf). Two consequences: the iteration converges geometrically from any strictly
positive start, and \textbf{the fixed point is unique up to the per-sector scale}, which is
exactly the gauge the reference normalization fixes. Given trade costs and expenditures, there
is one and only one vector of comparative advantages consistent with the observed area shares.

\emph{Step two: the general-equilibrium feedback.} Expenditure is not fixed. A change in $T$
moves the price indices $p_{rs}$, hence $c_{r}$, hence downstream sales $y_{r}$ and the weights
$\omega_{rs}$. The map actually iterated is the composition $\Psi=\mathcal{S}\circ\omega$, whose
Jacobian decomposes as
\begin{equation}
D\Psi = \underbrace{D\mathcal{S}\big|_{\omega\text{ fixed}}}_{\text{scaling channel}}
\;+\;\underbrace{J_{\mathrm{GE}}}_{\text{expenditure channel}},
\qquad
J_{\mathrm{GE}}\equiv D\Psi-D\mathcal{S}\big|_{\omega\text{ fixed}} .
\label{eq:app-jacobian-split}
\end{equation}
Birkhoff's theorem bounds the first term by $\kappa_{S}$ but says nothing about the second, so a
sufficient condition for $\Psi$ to be a contraction is
\begin{equation}
\kappa_{S}+\big\|J_{\mathrm{GE}}\big\|<1 .
\label{eq:app-contraction}
\end{equation}
The size of $J_{\mathrm{GE}}$ is governed by how far the technology is from Cobb--Douglas: the
expenditure weights are invariant to $T$ when $\lambda=\nu=1$, so $\|J_{\mathrm{GE}}\|$ vanishes
with $|1-\lambda|$ and $|1-\nu|$. Our calibration is not close to Cobb--Douglas, so
\eqref{eq:app-contraction} is an empirical question rather than an assumption, and we verify it
numerically at the estimated parameters by finite differences on the one-step map. We find
$\kappa_{S}\approx0.005$ and $\|J_{\mathrm{GE}}\|\approx0.015$, so the sufficient condition
\eqref{eq:app-contraction} holds with a wide margin ($\approx0.02\ll1$); the measured spectral
radius of the full map is $\approx0.012$, and convergence to $10^{-9}$ takes a handful of
iterations.%
% NOTE: these three numbers were measured by the Phase-0 inversion gate
% (test/test_ge_inversion.jl). Re-run it under --granular=true --ca_level=aa to
% refresh them before circulating, since the earlier measurement was taken on the
% ZE-level continuum configuration.
\ We further confirm that the inversion recovers a known $T$ from its own implied
shares to machine precision, and that ten dispersed starting values converge to the same fixed
point.

This is a local, numerically verified statement at the estimated parameters, not a global
theorem: \eqref{eq:app-contraction} is sufficient, not necessary, and we do not characterize the
region of the parameter space over which it holds. Should it fail at some candidate parameter
vector, the iteration is not used blindly --- non-convergence is detected and the candidate
retains its warm-start $T$, so the outer optimizer sees a well-defined objective throughout.

Substituting $T^{*}(\alpha,\Omega,A)$ into \eqref{eq:app-smm} collapses the outer search from
several hundred parameters to a handful and makes the sourcing-share block just-identified along
the profiled manifold, so that the extensive-margin coefficient is fit by $(\alpha,\Omega,A)$
alone. Note that just-identification here is at the \emph{area} level: within an area, the split
of sourcing across regions remains over-identified by the distance profile, which is the content
of Assumption~\ref{as:geo}.

\paragraph{Weighting matrix.} In the SMM framework noise arises from two sources: sampling
variability in the data and simulation error. We distinguish between moments for which
bootstrapping is feasible (sourcing shares, extensive-margin coefficients and count moments) and
those for which it is not (input--output coefficients and the downstream sales distribution). We
therefore approximate the variance--covariance matrix of the data moments, $\Sigma_{\mathrm{data}}$,
with a block matrix that accounts for within-group interactions; the block corresponding to
non-bootstrappable moments is set to zero. Because the same firms generate the sourcing shares,
the extensive-margin coefficients and the counts, these three blocks are estimated by a
\emph{single joint} bootstrap, which is what delivers the cross-covariance blocks the efficient
weight matrix requires. The resampling unit is the supplier firm, pooled over the geography of
active areas, with the denominator of $\bar G_{s}(0)$ held fixed --- so that a region with one
or two firms can draw zero and flip into the empty set, which is precisely the variance the
count moment needs.

\paragraph{Optimization procedure.} The estimator is computed using Particle Swarm Optimization
\citep{kennedy1995}, which explores the parameter space through parallel evaluation, with
adaptive inertia weights declining from $0.9$ to $0.4$ over iterations and cognitive and social
parameters set to $1.5$ and $2.5$. The estimation proceeds in three steps.

\begin{enumerate}[topsep=3pt,itemsep=3pt,label=\arabic*.]
\item Starting values are set as follows: $\Omega_{L}$ and $\{\Omega_{s}\}$ at their empirical
counterparts from input--output tables; $\{A_{r}\}$ proportional to $\pi_{r}^{1/\varepsilon}$;
$\alpha$ at the extensive-margin regression estimate divided by $\theta$; and $\{T_{as}\}$ at the
matrix-scaling solution \eqref{eq:app-sinkhorn} evaluated at that $\alpha$. We then run the
optimizer over the profiled parameter vector using the identity weighting matrix, refining with
an alternating block scheme in which each iteration sequentially optimizes $\{\Omega_{L},\Omega_{s}\}$,
then $\alpha$, then $\{A_{r}\}$, with a search radius progressively reduced across iterations.
This yields $\hat\mu_{1}$.
\item Using $\hat\mu_{1}$ we estimate the variance--covariance matrix of the simulated moments,
$\Sigma_{\mathrm{sim}}$, from repeated re-simulation with independent draws, and construct
$W=(\Sigma_{\mathrm{data}}+\Sigma_{\mathrm{sim}})^{-1}$.
\item We re-estimate the model with the updated weighting matrix, initializing at $\hat\mu_{1}$,
to obtain $\hat\mu_{2}$.
\end{enumerate}

\paragraph{Standard errors.} Let $G=\partial\tilde M/\partial\mu$ evaluated at $\hat\mu_{2}$ by
finite differences on common random numbers, and $\Omega=\Sigma_{\mathrm{data}}+\Sigma_{\mathrm{sim}}$.
We report the efficient variance $(G'WG)^{-1}$ alongside the sandwich form
$(G'WG)^{-1}G'W\Omega WG(G'WG)^{-1}$, which coincide when $W=\Omega^{-1}$; the ratio is a
diagnostic on the weighting matrix. Because $T$ is concentrated out through
\eqref{eq:app-sinkhorn} rather than searched, the Jacobian is taken along the profiled manifold
--- perturbing $\alpha$ and letting $T^{*}(\alpha)$ follow --- so the reported standard error on
$\alpha$ is a total derivative, and confidence intervals for $T$ are obtained by the delta
method, propagating both the uncertainty in $\hat\alpha$ and the sampling uncertainty in the
sourcing-share targets that \eqref{eq:app-sinkhorn} inverts.

\subsection{The price index, love of variety, and the normalization}
\label{app:model-price}

The one place where granularity meets the value block requires an explicit assumption. For
variety $\rho$ in market $(r,s)$, origin $r'$'s champion cost satisfies
$\Pr(c_{r'\rho s}(r)>x)=\exp(-T_{a(r')s}(w_{r's}\tau_{r'rs})^{-\theta}x^{\theta})$, so by
independence across origins the minimum $c^{*}_{\rho}$ is Weibull:
\begin{equation}
\Pr\big(c^{*}_{\rho}>x\big)=\exp\big(-\Phi_{rs}x^{\theta}\big),
\qquad
\mathbb{E}\big[(c^{*}_{\rho})^{k}\big]=\Phi_{rs}^{-k/\theta}\,\Gamma\Big(1+\tfrac{k}{\theta}\Big),
\quad k>-\theta .
\end{equation}
With the CES aggregate taken as a sum,
$\mathbb{E}[P_{rs}^{1-\nu_{s}}]=N_{s}\Gamma(1+\tfrac{1-\nu_{s}}{\theta})\Phi_{rs}^{-\frac{1-\nu_{s}}{\theta}}$
up to a constant, so the certainty-equivalent index
$\widetilde P_{rs}=(\mathbb{E}[P_{rs}^{1-\nu_{s}}])^{1/(1-\nu_{s})}$ satisfies
\begin{equation}
\widetilde P_{rs}=N_{s}^{\frac{1}{1-\nu_{s}}}\times
\underbrace{\Gamma\Big(1+\tfrac{1-\nu_{s}}{\theta}\Big)^{\frac{1}{1-\nu_{s}}}\Phi_{rs}^{-1/\theta}}_{N_{s}\text{-free}} .
\label{eq:app-loveofvariety}
\end{equation}
Since $1-\nu_{s}<0$ the factor is decreasing in $N_{s}$: more varieties lower the sector price
index. Rescaling $\widetilde T_{as}=T_{as}N_{s}^{\kappa}$ with $\kappa=\theta/(1-\nu_{s})$ maps
$\Phi_{rs}\mapsto N_{s}^{\kappa}\Phi_{rs}$ and cancels the two powers exactly, leaving
$\widetilde P_{rs}$ invariant to $N_{s}$. Equivalently, taking the CES aggregate as an
\emph{average} rather than a sum --- as in \eqref{eq:app-ces} --- implements the same
normalization. It is consistent because it acts on a gauge the sourcing shares do not see: the
share \eqref{eq:app-share} is a ratio of $T$'s within a sector, so the extensive margin, the
count moment and the within-area contrast are all unchanged. It is also a gauge the estimator
already fixes, by normalizing one reference area per sector.

Two caveats. First, $P_{rs}$ is a \emph{convex} transform of the average in \eqref{eq:app-ces},
so the realized index exceeds its certainty-equivalent value by Jensen's inequality, and the gap
closes only as fast as the average concentrates. At the calibration the regularity condition
$\theta>2(\nu_{s}-1)$ holds with little slack, and the gap is of the order of $6\%$ at
$N_{s}\approx10^{2}$, falling to $1\%$ at $N_{s}\approx10^{3}$. The estimator adopts the
certainty-equivalent reading --- downstream firms optimize against
$\mathbb{E}[P^{1-\nu_{s}}]^{1/(1-\nu_{s})}$, neglecting granular price randomness, the
general-equilibrium embedding of \citet{lmm2023}. Most of that $6\%$ would not have been
identifying variation in any case: a price-level shift is observationally equivalent to a
technology shift, so the sector technology terms absorb it and only the cross-sector pattern of
$N_{s}$ survives as signal. And by Lemma~\ref{lem:qfree} none of it touches $q_{r's}$ or the
extensive margin, so the identification of $\alpha$ and $N_{s}$ is unaffected either way.

Second, the assumption must not survive into counterfactuals. The elasticity
$\partial\ln\widetilde P_{rs}/\partial\ln N_{s}=1/(1-\nu_{s})<0$ is a genuine gain from variety
and a live propagation channel: a shock that removes suppliers in one region raises $P_{rs}$
downstream through exactly this elasticity, and that channel is absent from the continuum model.
Any counterfactual that moves the variety set must map back to the physical $T$ and use the
realized index.

\subsection{Proofs}
\label{app:model-proofs}

\paragraph{Equation \eqref{eq:app-champion}.} The number of firms in $(r',s,\rho)$ with
productivity above $z$ is Poisson with mean $T_{r's}z^{-\theta}$, so the probability that none
exceeds $z$ --- equivalently that the maximum is below $z$ --- is $\exp(-T_{r's}z^{-\theta})$,
the $\mathrm{Fr\acute{e}chet}(T_{r's},\theta)$ CDF. \hfill$\square$

\paragraph{Lemma~\ref{lem:fe}.} (i) With $\Pr(y=1)=1-\exp(-\exp(\phi_{g}+\beta x))$, a group in
which $y\equiv0$ contributes
$\ell_{g}=\sum_{i}\ln\exp(-e^{\phi_{g}+\beta x_{i}})=-e^{\phi_{g}}\sum_{i}e^{\beta x_{i}}$. This
is increasing in $-\phi_{g}$ with supremum $0$ as $\phi_{g}\to-\infty$, at which
$\partial\ell_{g}/\partial\beta=-e^{\phi_{g}}\sum_{i}x_{i}e^{\beta x_{i}}\to0$. The group
therefore drops from the concentrated likelihood.

(ii) By Frisch--Waugh--Lovell,
$\hat\beta=\sum_{g,i}\omega\tilde x\tilde y/\sum_{g,i}\omega\tilde x^{2}$ with $\tilde{\cdot}$
the weighted within-group demeaning. A group with $y\equiv0$ has $\tilde y\equiv0$, so it
contributes zero to the numerator and $\sum\omega\tilde x^{2}>0$ to the denominator. Splitting
the denominator between $\mathcal{A}^{+}$ and its complement gives
$\hat\beta=\varkappa\hat\beta^{+}$. \hfill$\square$

\paragraph{Proposition~\ref{prop:firm}.} Condition on the firm's productivity $z$; its delivered
cost to buyer $a$ is $c=w_{r's}\tau_{r'as}/z$. Each competitor $r''$ has champion cost
$w_{r''s}\tau_{r''as}/\varphi_{r''}$ with
$\varphi_{r''}\sim\mathrm{Fr\acute{e}chet}(T_{a(r'')s},\theta)$, so
$\Pr(w_{r''s}\tau_{r''as}/\varphi_{r''}>x)=\exp(-T_{a(r'')s}(w_{r''s}\tau_{r''as})^{-\theta}x^{\theta})$.
By independence across competitors the minimum competing cost $M$ satisfies
$\Pr(M>x)=\exp(-\Phi^{-r'}_{as}x^{\theta})$. The firm supplies iff $M>c$, giving
\eqref{eq:app-firm-win}. Taking the complement,
$\Pr(\text{no supply}\mid z)=1-\exp(-\exp(\eta))$ with
$\eta=\ln\Phi^{-r'}_{as}+\theta\ln(w_{r's}\tau_{r'as})-\theta\ln z$, and
$\tau_{r'as}=d_{r'a}^{\alpha}$ gives \eqref{eq:app-firm-cloglog}. No step involves the variety
count. \hfill$\square$

\paragraph{Lemma~\ref{lem:qfree}.} Origin $r'$ wins variety $\rho$ for buyer $r$ iff
$w_{r's}\tau_{r'rs}/\varphi_{r'\rho s}<\min_{r''\neq r'}w_{r''s}\tau_{r''rs}/\varphi_{r''\rho s}$,
an event measurable with respect to $\{\varphi_{\cdot\rho s}\}$, $w$ and $\tau$ alone; the
$\varphi$ have distribution governed by $(T_{\cdot s},\theta)$. The win-somewhere event is the
union over $r$ of those events, hence measurable with respect to the same objects. Neither the
variety count nor any price or expenditure aggregate enters. \hfill$\square$

\paragraph{Unbiasedness of \eqref{eq:app-Gclosed}.} Let $k$ be the number of the $m$ simulated
varieties won by $r'$, each won independently with probability $q$. Draw $n\le m$ of the $m$
varieties uniformly without replacement; conditional on $k$, the probability that all $n$ are
losses is $\binom{m-k}{n}/\binom{m}{n}$. Marginalizing over the draws, each of the $n$ selected
varieties is a loss independently with probability $1-q$, so the unconditional probability is
$(1-q)^{n}$. Hence $\mathbb{E}[\binom{m-k}{n}/\binom{m}{n}]=(1-q)^{n}$ exactly, with the convention
$\binom{a}{n}=0$ for $a<n$. \hfill$\square$
 from the main file. Requires amsmath, amssymb, amsthm,
% booktabs, enumitem and natbib. Add to the preamble if not already present:
%
%   \newtheorem{proposition}{Proposition}
%   \newtheorem{lemma}{Lemma}
%   \newtheorem{assumption}{Assumption}
%   \newtheorem{remark}{Remark}
%   \DeclareMathOperator*{\argmin}{arg\,min}
%
% Labels used by core_model.tex: app:model-derivation, app:model-foreign,
% app:model-calibration, app:model-prop, app:model-zeros, app:model-bounds.
% =====================================================================

\section{Details on the Structural Model}
\label{app:model}

\subsection{Derivation}
\label{app:model-derivation}

\paragraph{Setup.} The economy is divided into $R+1$ regions, populated by two types of firms:
producers of a final good and producers of intermediate goods used as inputs in the production
of the final good. The $(R+1)$th region represents the rest of the world. The analysis primarily
focuses on the matching of firms across domestic regions, controlling for competition from
foreign firms using their shares in total sales as reference. We assume that each domestic
region hosts at most one representative firm in the downstream sector, identified by its region
of origin. Intermediate goods are produced across $S$ sectors, with each region potentially
hosting many suppliers in each sector. We define a market $(s,r')$ as the collection of firms
located in region $r'$ that belong to upstream sector $s$.

Domestic regions are partitioned into attraction areas,
$\mathcal{R}=\bigsqcup_{a\in\mathcal{A}}\mathcal{R}_{a}$, where $\mathcal{R}_{a}$ is the set of
regions whose nearest downstream producer is the anchor of area $a$, and $a(r')$ denotes the
area of region $r'$.

\paragraph{Demand for final manufacturing products.} Firms in the downstream sector each own the
blueprint to produce a unique differentiated variety of the final good. The market for the final
good is characterized by monopolistic competition and iso-elastic demand functions:
\begin{equation}
q_{r}=\left(\frac{p_{r}}{P}\right)^{\varepsilon-1}\frac{E}{P}\,\delta_{r},
\end{equation}
where $q_{r}$ denotes the quantity demanded of the variety produced in region $r$ at price
$p_{r}$, given nominal consumption $E$ and the downstream sector's price index $P$. The
parameter $\varepsilon$ measures the price elasticity of sales, and $\delta_{r}$ is an exogenous
component of demand, varying across downstream producers and later used to simulate demand
shocks and their diffusion along production linkages. We abstract from foreign demand and assume
that the domestic market for final consumption goods is perfectly integrated.

\paragraph{Technology in the downstream sector.} Production of final goods requires the assembly
of intermediates, combined with labor, under a nested CES structure:
\begin{equation}
y_{r}=A_{r}\Big[\Omega_{L}^{\frac{1}{\lambda}}\ell_{r}^{\frac{\lambda-1}{\lambda}}
+(1-\Omega_{L})^{\frac{1}{\lambda}}\big(q^{I}_{r}\big)^{\frac{\lambda-1}{\lambda}}\Big]^{\frac{\lambda}{\lambda-1}},
\end{equation}
where $\ell_{r}$ is labor demand and $q^{I}_{r}$ the firm's intermediate input consumption. The
parameter $\Omega_{L}$ governs the share of labor in total production costs, while $\lambda$ is
the elasticity of substitution between labor and intermediates. Firm productivity $A_{r}$ is
exogenous. Intermediate consumption is itself a nested CES aggregator across sectors and, within
each sector, across varieties:
\begin{equation}
q^{I}_{r}=\Big[\sum_{s=1}^{S}\Omega_{s}^{\frac{1}{\nu}}q_{rs}^{\frac{\nu-1}{\nu}}\Big]^{\frac{\nu}{\nu-1}},
\qquad
q_{rs}=\Big[\frac{1}{N_{s}}\sum_{\rho=1}^{N_{s}}q_{rs}(\rho)^{\frac{\nu_{s}-1}{\nu_{s}}}\Big]^{\frac{\nu_{s}}{\nu_{s}-1}},
\label{eq:app-ces}
\end{equation}
where $\Omega_{s}$ governs the share of sector $s$ in intermediate consumption, $\nu$ is the
elasticity of substitution across sectors and $\nu_{s}$ that across varieties within a sector.

The only change relative to the continuum formulation is the inner aggregator: the integral over
$j\in[0,1]$ is replaced by an average over the $N_{s}$ varieties. Writing it as an
\emph{average} rather than a sum is a normalization of the love-of-variety level, discussed and
justified in Appendix~\ref{app:model-price}; it leaves every sourcing share, and hence the
entire extensive margin, unchanged.

Given the nested CES structure, the nominal demand for intermediate products is
\begin{equation}
x_{rs}(\rho)=\Omega_{s}(1-\Omega_{L})
\left(\frac{p_{rs}(\rho)}{p_{rs}}\right)^{1-\nu_{s}}
\left(\frac{p_{rs}}{p_{r}}\right)^{1-\nu}
\left(\frac{p_{r}}{c_{r}}\right)^{1-\lambda}\frac{c_{r}y_{r}}{A_{r}},
\end{equation}
with $p_{rs}\equiv\big[\tfrac{1}{N_{s}}\sum_{\rho}p_{rs}(\rho)^{1-\nu_{s}}\big]^{\frac{1}{1-\nu_{s}}}$
the cost index for intermediates in sector $s$,
$p_{r}\equiv\big[\sum_{s}\Omega_{s}p_{rs}^{1-\nu}\big]^{\frac{1}{1-\nu}}$ the cost index for
intermediates, and
$c_{r}\equiv\big[\Omega_{L}w_{r}^{1-\lambda}+(1-\Omega_{L})p_{r}^{1-\lambda}\big]^{\frac{1}{1-\lambda}}$
the unit cost of the firm. In equilibrium $y_{r}=q_{r}$.

\paragraph{Technology in upstream sectors.} Intermediate goods are produced with labor using a
constant-returns-to-scale technology. Wages in sector $s$ and region $r'$ are denoted $w_{r's}$
and taken as given. As in Eaton and Kortum (2002) we assume that varieties produced within and
across regions are perfect substitutes and that productivity draws are independent. The
departure from the continuum case is that each sector contains $N_{s}$ varieties.

\begin{assumption}[Granularity]
\label{as:granular}
Each upstream sector $s$ consists of $N_{s}\in\mathbb{N}$ horizontally differentiated varieties
$\rho=1,\dots,N_{s}$. Within each variety $\rho$ and each region $r'$ there is a continuum of
firms whose productivities above $z$ form a Poisson process with mean $T_{r's}z^{-\theta}$.
Draws are independent across varieties, regions and sectors.
\end{assumption}

The continuum of firms is what keeps the model tractable. The number of firms above $z$ in
$(r',s,\rho)$ is Poisson with mean $T_{r's}z^{-\theta}$, so the most productive one satisfies
\begin{equation}
\Pr\big(\varphi_{r'\rho s}\le z\big)=\Pr(\text{no firm above }z)=e^{-T_{r's}z^{-\theta}},
\label{eq:app-champion}
\end{equation}
that is, $\varphi_{r'\rho s}$ is exactly $\mathrm{Fr\acute{e}chet}(T_{r's},\theta)$. One
Fr\'echet draw per (region, variety) pair \emph{is} the champion of a continuum of firms, not an
approximation to it. Since the iceberg cost scales all firms of region $r'$ identically, that
champion is the unique region-$r'$ firm that ever serves any buyer sourcing $\rho$ from $r'$:
\textbf{a firm is a (region, variety) champion}, with delivered cost
$c_{r'\rho s}(r)=w_{r's}\tau_{r'rs}/\varphi_{r'\rho s}$. This is the ``fixed number of
independent Ricardian draws'' form of granular Eaton--Kortum \citep[cf.][]{eks2012,bejk2003}.

\begin{assumption}[Area-level comparative advantage]
\label{as:geo}
Comparative advantage is area-specific: $T_{r's}=T_{a(r')s}$ for every $r'\in\mathcal{R}_{a}$,
with one parameter per (sector, area). Productivity draws remain independent across regions ---
two regions in the same area share the Fr\'echet scale, not the realization of their champions.
\end{assumption}

Assumption~\ref{as:geo} does not make regions within an area interchangeable. Each runs its own
Ricardian competition for each variety, and two regions at different distances from a buyer have
different win probabilities. What it removes is unobserved technological heterogeneity within an
area, which is what allows the model to explain the same variation the extensive-margin
regression uses.

\paragraph{Equilibrium.} The resolution follows the same steps as in Eaton and Kortum (2002).
Given the Fr\'echet assumption, the minimum price at which suppliers from $(r',s)$ can serve
market $r$ is distributed
\begin{equation}
G_{r'rs}(p)=\Pr\Big(\frac{w_{r's}\tau_{r'rs}}{\varphi_{r'\rho s}}\le p\Big)
=1-e^{-\left[T_{a(r')s}(w_{r's}\tau_{r'rs})^{-\theta}\right]p^{\theta}},
\end{equation}
so the distribution of input prices for downstream firms in region $r$ is
\begin{equation}
G_{rs}(p)=1-\prod_{r'=1}^{R+1}\big[1-G_{r'rs}(p)\big]=1-e^{-\Phi_{rs}p^{\theta}},
\qquad
\Phi_{rs}=\sum_{a'}T_{a's}\sum_{r'\in\mathcal{R}_{a'}}(w_{r's}\tau_{r'rs})^{-\theta}.
\label{eq:app-Phi}
\end{equation}
The probability that a downstream firm in location $r$ picks a supplier in market $(r',s)$ is
thus
\begin{equation}
\gamma_{r'rs}=\frac{T_{a(r')s}(w_{r's}\tau_{r'rs})^{-\theta}}{\Phi_{rs}} .
\label{eq:app-share}
\end{equation}
Taking the perspective of the supplier, the probability of being the lowest-cost supplier in
market $(r,s)$, given location $r'$ and productivity $z$, is
\begin{equation}
\rho_{r'rs}(z)=e^{-\Phi_{rs}(w_{r's}\tau_{r'rs})^{\theta}z^{-\theta}}
=e^{-T_{a(r')s}\gamma^{-1}_{r'rs}z^{-\theta}} .
\label{eq:app-rho}
\end{equation}
Because the conditional distribution of prices is not origin-specific,
$G_{r'rs}(p)=G_{rs}(p)$, and $\gamma_{r'rs}$ is also the expected share of intermediates from
$r'$ purchased by a downstream firm in $r$. The expected cost indices are
\begin{equation}
p_{rs}=\Big[\Gamma\Big(\frac{\theta+1-\nu_{s}}{\theta}\Big)\Big]^{\frac{1}{1-\nu_{s}}}\Phi_{rs}^{-1/\theta},
\qquad
p_{r}=\Big[\sum_{s}\Omega_{s}p_{rs}^{1-\nu}\Big]^{\frac{1}{1-\nu}},
\qquad
c_{r}=\Big[\Omega_{L}w_{r}^{1-\lambda}+(1-\Omega_{L})p_{r}^{1-\lambda}\Big]^{\frac{1}{1-\lambda}},
\end{equation}
where the first expression is the certainty-equivalent index of Appendix~\ref{app:model-price}.
Bilateral flows are
\begin{equation}
X_{r'rs}=\gamma_{r'rs}\,x_{rs}
=\gamma_{r'rs}\,\Omega_{s}(1-\Omega_{L})
\left(\frac{p_{rs}}{p_{r}}\right)^{1-\nu}\left(\frac{p_{r}}{c_{r}}\right)^{1-\lambda}\frac{c_{r}y_{r}}{A_{r}} .
\end{equation}
Since we do not observe these bilateral flows directly, we cannot use this equation to estimate
the model's structural parameters and rely instead on simulated moments.

\subsection{The extensive margin and the two classes of zeros}
\label{app:model-zeros}

Across the $N_{s}$ varieties the win events are i.i.d., so with $K_{r'rs}$ the number of
sector-$s$ varieties buyer $r$ sources from $r'$,
\begin{equation}
K_{r'rs}\sim\mathrm{Bin}\big(N_{s},\gamma_{r'rs}\big),
\qquad
\Pr(K_{r'rs}=0)=(1-\gamma_{r'rs})^{N_{s}}>0 .
\end{equation}
Under a continuum this probability is $\lim_{N\to\infty}(1-\gamma_{r'rs})^{N}=0$ for every
origin with $T_{r's}>0$: the extensive margin is degenerate and the observed zeros are not a
model outcome. Finiteness is exactly what restores it.

\begin{remark}[Dependence across buyers]
\label{rem:dependence}
For a given variety the champion costs are common across buyers; only $\tau_{r'rs}$ differs. Win
events are therefore positively dependent across buyers, and the union probability $q_{r's}$ is
not the product of the per-market marginals. The per-market marginal above is nonetheless exact,
and $q_{r's}$ is evaluated by simulation. This matters because the survey records whether a firm
supplies the industry, not which buyer it serves: the empirical outcome is the union, and the
model counterpart is computed as the same union.
\end{remark}

Assumption~\ref{as:geo} sorts observed zeros into two economically distinct classes, and the
distinction determines what is estimated.

\begin{table}[htbp]
\centering
\small
\begin{tabular}{@{}p{0.21\textwidth}p{0.24\textwidth}p{0.19\textwidth}p{0.28\textwidth}@{}}
\toprule
\textbf{Zero} & \textbf{Treatment} & \textbf{Probability} & \textbf{Role}\\
\midrule
Whole attraction area $a$ has no supplier in $s$
& \emph{Structural}: $T_{as}=0$, $(s,a)$ outside the active set
& $1$ by construction
& None --- absorbed by the $a\times s$ fixed effect (Lemma~\ref{lem:fe}); no parameter estimated\\[4pt]
Region $r'$ inside an active area has no supplier
& \emph{Endogenous}: inherits $T_{as}>0$, own distance $d_{r'}$
& $(1-q_{r's})^{N_{s}}$
& \textbf{The identifying variation}: disciplined by distance alone\\
\bottomrule
\end{tabular}
\caption{The two classes of extensive-margin zeros. The second row is exactly the set of regions
that host firms but no supplier. Under region-level comparative advantage these had to be
assigned $T\equiv0$, making their zeros exogenous; under Assumption~\ref{as:geo} they inherit
their area's comparative advantage and their zeros become model predictions.}
\label{tab:app-hierarchy}
\end{table}

The reason the fixed effect dictates this split is the following.

\begin{lemma}[All-zero groups carry no information about the distance slope]
\label{lem:fe}
Index observations by firm $i$ within fixed-effect group $g=(a,s)$, with outcome $y_{i}$,
regressor $x_{i}$, weights $\omega_{i}>0$ and group effect $\phi_{g}$.
\begin{enumerate}[nosep,topsep=2pt,label=(\roman*)]
\item \emph{Complementary log-log.} The log-likelihood of a group in which $y\equiv0$ is
$\ell_{g}=-e^{\phi_{g}}\sum_{i}e^{\beta x_{i}}$, whose supremum over $\phi_{g}$ is attained as
$\phi_{g}\to-\infty$ with $\ell_{g}\to0$ and $\partial\ell_{g}/\partial\beta\to0$: the group
drops out of the concentrated likelihood and contributes exactly nothing to $\hat\beta$.
\item \emph{Linear probability model.} By Frisch--Waugh--Lovell an all-zero group has
$\tilde y\equiv0$, contributing zero to the numerator and a positive amount to the denominator,
so $\hat\beta=\varkappa\hat\beta^{+}$ with
$\varkappa=\big(\sum_{a\in\mathcal{A}^{+}}\sum\omega\tilde x^{2}\big)/\big(\sum_{a}\sum\omega\tilde x^{2}\big)\in(0,1]$,
where $\mathcal{A}^{+}$ is the set of areas hosting at least one supplier.
\end{enumerate}
\end{lemma}

Under the cloglog link --- the link the model implies --- $\alpha$ is therefore identified
exclusively from variation \emph{within} attraction areas that host at least one supplier. Under
a linear probability model the empty areas enter only as a mechanical dilution of the same
within-area signal. Either way, an area with no supplier at all carries no information about the
distance slope, which is why the model does not attempt to explain it.

\subsection{The firm-level reduced form}
\label{app:model-prop}

The object the estimation matches is the union event \eqref{eq:q}: whether a firm supplies the
downstream industry at all. That probability has no tractable closed form. The single-buyer case
does, and it is what gives the empirical cloglog its structural interpretation; we state it
first and then say precisely what survives when there are several buyers.

\begin{proposition}[Firm-level reduced form, single buyer]
\label{prop:firm}
Let a firm be a (region, variety) champion with productivity $z$ located in $r'$, and suppose it
faces a single downstream buyer, its area's anchor $a$. Writing
$\Phi^{-r'}_{as}=\sum_{r''\neq r'}T_{a(r'')s}(w_{r''s}\tau_{r''as})^{-\theta}$ for the
competitors' cost scale,
\begin{equation}
\Pr\big(\text{the firm supplies}\mid z\big)
=\exp\Big(-\Phi^{-r'}_{as}\big(w_{r's}\tau_{r'as}\big)^{\theta}z^{-\theta}\Big),
\label{eq:app-firm-win}
\end{equation}
so the probability that it does \emph{not} supply is a complementary log-log function
\begin{equation}
\Pr\big(\text{no supply}\mid z\big)=1-\exp\big(-\exp(\eta)\big),
\qquad
\eta=\underbrace{\ln\Phi^{-r'}_{as}}_{\text{area}\times\text{sector FE}}
+\theta\ln w_{r's}
+\theta\alpha\,\ln d_{r'a}
-\theta\,\ln z .
\label{eq:app-firm-cloglog}
\end{equation}
The variety count $N_{s}$ does not appear.
\end{proposition}

Three consequences, and they are the backbone of the estimation.

\begin{enumerate}[nosep,topsep=3pt,label=(\alph*)]
\item \textbf{The cloglog is the model's own auxiliary model.} A cloglog of the ``does not
supply'' indicator on log distance and log own productivity, with an area$\times$sector fixed
effect, is exactly \eqref{eq:app-firm-cloglog} when a firm faces one buyer, with distance
coefficient $+\theta\alpha$ and productivity coefficient $-\theta$.
\item \textbf{$\alpha$ and $N_{s}$ separate.} $N_{s}$ is absent from
\eqref{eq:app-firm-cloglog}: conditioning on the firm removes it. The regression identifies
$\alpha$; the count level identifies $N_{s}$; there is no residual coupling.
\item \textbf{The productivity coefficient is a free over-identifying test.} It equals
$-\theta$. Extending the moment vector to include it would let $\theta$ be estimated rather than
calibrated.
\end{enumerate}

\paragraph{Several buyers: what survives.} Our surveys record supply to the industry, not to a
named buyer, so the empirical outcome is the union event $\bigcup_{r}W_{r}$, where $W_{r}$
denotes winning buyer $r$. Its probability is \emph{not} $1-\prod_{r}[1-\rho_{r'rs}(z)]$: that
expression would require the $W_{r}$ to be independent, and they are not.

The dependence is positive and structural: conditional on $z$, the firm's delivered cost to
every buyer is the same number scaled by $\tau_{r'rs}$, so a firm cheap enough to beat the
competition in one market tends to beat it everywhere. The exact object is the
inclusion--exclusion sum
\begin{equation}
\Pr\Big(\bigcup_{r}W_{r}\,\Big|\,z\Big)
=\sum_{\emptyset\neq\mathcal{S}\subseteq\{1,\dots,R\}}(-1)^{|\mathcal{S}|+1}
\exp\Big(-z^{-\theta}\max_{r\in\mathcal{S}}\big\{\Phi^{-r'}_{rs}(w_{r's}\tau_{r'rs})^{\theta}\big\}\Big),
\label{eq:app-inclexcl}
\end{equation}
where the maximum appears because $\cap_{r\in\mathcal{S}}W_{r}$ requires the firm to beat the
\emph{tightest} of the thresholds it faces. Equation~\eqref{eq:app-inclexcl} has $2^{R}-1$
terms. Many are redundant --- if the threshold for buyer $r$ dominates that for $r''$ then
$W_{r}\subseteq W_{r''}$ and the pair collapses --- and pruning the dominated buyers reduces the
sum to the non-dominated set, but that set is still large enough in our geography that
evaluating \eqref{eq:app-inclexcl} inside an optimizer is not practical. The independence
approximation $1-\prod_{r}(1-\rho_{r'rs}(z))$, which is what a closed-form estimator would have
to use, is a strict \emph{under}-statement of \eqref{eq:app-inclexcl} by positive dependence, and
the error is largest exactly where the extensive margin has content, at high $z$.

This is the reason the model is estimated by simulated rather than analytical method of moments.
Simulation evaluates the union exactly --- each variety's champion costs are drawn once and
compared against every buyer, so the dependence is reproduced by construction --- at the cost of
simulation noise, which the weighting matrix accounts for. We therefore treat
\eqref{eq:app-firm-cloglog} as an auxiliary model rather than an estimating equation: the same
cloglog is run on simulated and observed data, with the union as the outcome on both sides, and
the coefficients are matched. This is consistent for $\alpha$ whether or not the link is exact,
and it is why the distance coefficient enters the moment vector rather than being inverted
analytically.

The outcome convention is dictated by the unit of observation rather than chosen. Simulating the
exact object at $\theta=1$, $\alpha=0.30$, with firms as (region $\times$ variety) champions,
the ``does not supply'' outcome recovers $\hat\beta_{\ln d}=+0.327$ and
$\hat\beta_{\ln z}=-1.051$ against truths $\theta\alpha=0.30$ and $-\theta=-1.0$; the reversed
outcome recovers neither ($-0.426$ and $+1.452$).

\subsection{The count moment and bounds on the variety count}
\label{app:model-bounds}

The zeros are a cell-level object. With $\mathcal{R}^{+}_{s}$ the regions inside active
attraction areas, the empirical moment is the share of those regions hosting at most $n$
suppliers,
\begin{equation}
\bar G_{s}(n)=\frac{1}{|\mathcal{R}^{+}_{s}|}
\#\big\{r'\in\mathcal{R}^{+}_{s}:K_{r's}\le n\big\},
\qquad
\mathbb{E}\big[\bar G_{s}(0)\big]=\frac{1}{|\mathcal{R}^{+}_{s}|}\sum_{r'}(1-q_{r's})^{N_{s}} .
\label{eq:app-count}
\end{equation}
Writing $\mu_{r's}=N_{s}q_{r's}$ and using the Poisson approximation,
$\partial\bar G_{s}(K)/\partial N_{s}
=-|\mathcal{R}^{+}_{s}|^{-1}\sum_{r'}q_{r's}e^{-\mu_{r's}}\mu_{r's}^{K}/K!$, which peaks at
$K\approx\mu_{r's}$ and vanishes as $\bar G_{s}(K)\to1$: saturated thresholds carry no
identifying variation and a poorly estimated variance, so we use the empty-region share
$\bar G_{s}(0)$, which is both the most interpretable threshold and typically near-optimal.

Because the iceberg cost never changes which firm is cheapest within a region, each variety
generates one distinct supplying firm per winning origin. So if $N^{\mathrm{obs}}_{s}$ is the
number of distinct upstream firms observed to supply the downstream industry,
$N^{\mathrm{obs}}_{s}=\sum_{\rho}\#\{\text{origins winning }\rho\text{ somewhere}\}\in[N_{s},N_{s}R^{D}]$,
where $R^{D}$ is the number of downstream buyers, hence
\begin{equation}
\frac{N^{\mathrm{obs}}_{s}}{R^{D}}\;\le\;N_{s}\;\le\;N^{\mathrm{obs}}_{s} .
\label{eq:app-bounds}
\end{equation}
The bounds are tight: the upper when every variety is globally sourced from a single origin, the
lower when each of the $R^{D}$ buyers sources it from a different origin. Since
$\mathbb{E}[\sum_{r'}K_{r's}]=N_{s}\sum_{r'}q_{r's}$, matching the observed firm count gives a second,
closed-form estimate $N^{\mathrm{count}}_{s}=N^{\mathrm{obs}}_{s}/\sum_{r'}q_{r's}$ that lies
inside \eqref{eq:app-bounds} by construction --- a free over-identifying check.

\paragraph{Degrees of freedom.} Fix a sector, let $A^{+}=|\mathcal{A}^{+}_{s}|$ and let
$R^{\mathrm{act}}_{a}$ be the number of regions in area $a$ with a strictly positive observed
share. The $T$ block has $A^{+}-1$ free elements after the reference normalization, against
$\sum_{a}R^{\mathrm{act}}_{a}$ share moments, so Assumption~\ref{as:geo} converts what were free
parameters into
\begin{equation}
\sum_{a\in\mathcal{A}^{+}}\big(R^{\mathrm{act}}_{a}-1\big)
\quad\text{over-identifying restrictions,}
\end{equation}
one per active region beyond the first in each area. Each is a within-area statement: given the
area aggregate, the split across its regions must follow the distance profile. Under
region-level comparative advantage this count is zero.

\subsection{Structural estimation: the moments}
\label{app:model-moments}

We use six sets of theoretical moments, observed in the data using a combination of IO tables
and micro-level data.

\begin{enumerate}[topsep=3pt,itemsep=3pt,label=\arabic*.]
\item \textbf{Aggregate labor share at the level of the industry:}
\begin{equation}
\tilde\Omega_{L}\equiv\frac{\sum_{r}w_{r}\ell_{r}}{\sum_{r}c_{r}A_{r}^{-1}y_{r}}
=\sum_{r}\frac{c_{r}A_{r}^{-1}y_{r}}{\sum_{r}c_{r}A_{r}^{-1}y_{r}}\,\Omega_{L}
\left(\frac{w_{r}}{c_{r}}\right)^{1-\lambda}.
\end{equation}
In the model it is a sales-weighted average of firm-level labor shares. In the data we use the
share of labor costs in total costs from IO tables.

\item \textbf{Aggregate share of sector $s$ in the industry's input purchases:}
\begin{equation}
\tilde\Omega_{s}\equiv\frac{\sum_{r}p_{rs}q_{rs}}{\sum_{r}\sum_{s}p_{rs}q_{rs}}
=\sum_{r}\frac{p_{r}q^{I}_{r}}{\sum_{r}p_{r}q^{I}_{r}}\,\Omega_{s}
\left(\frac{p_{rs}}{p_{r}}\right)^{1-\nu}.
\end{equation}
In the data we use technical coefficients for the downstream industry from IO tables.

\item \textbf{The share of each region in downstream sales:}
\begin{equation}
\pi_{r}\equiv\frac{c_{r}A_{r}^{-1}y_{r}}{\sum_{r}c_{r}A_{r}^{-1}y_{r}}
=\left(\frac{p_{r}}{P}\right)^{-\varepsilon}
=\frac{(c_{r}A_{r}^{-1})^{-\varepsilon}}{\sum_{r}(c_{r}A_{r}^{-1})^{-\varepsilon}},
\end{equation}
recovered from sales data on the population of downstream producers.

\item \textbf{The extensive-margin distance coefficient.} The cloglog of
\eqref{eq:app-firm-cloglog}, estimated on the simulated firm population with the same
specification as Equation~(1) in the data: outcome ``does not supply'', regressors log distance
to the area anchor and log own productivity, area$\times$sector fixed effects.

\item \textbf{Area-aggregated sourcing shares.} The share of the downstream industry's purchases
of sector $s$ sourced from area $a$:
\begin{equation}
\tilde\gamma_{as}\equiv\sum_{r'\in\mathcal{R}_{a}}\frac{\sum_{r}X_{r'rs}}{\sum_{r}\sum_{r''}X_{r''rs}}
=\sum_{r'\in\mathcal{R}_{a}}\sum_{r}\frac{\sum_{r''}X_{r''rs}}{\sum_{r''}\sum_{r}X_{r''rs}}\,\gamma_{r'rs}.
\label{eq:app-gamma-area}
\end{equation}
In the data we reconstitute firm-level sales to the downstream industry using sales-share data,
$x_{i(r',s)d}=\mathbf{1}(\mathrm{Sup}_{i(r',s)})\,a^{D}_{d\,i(r',s)}\,x_{i(s,r')}$, sum over all
firms in the area and normalize by the sum over all firms from sector $s$. The aggregation runs
over \emph{every} region of the area, including those hosting no supplier, on both the model and
the data side.

\item \textbf{The count moment} $\bar G_{s}(0)$ of \eqref{eq:app-count}: the share of regions
inside active attraction areas that host no supplier of sector $s$.
\end{enumerate}

Moments 1--3 and 5 are the same objects as in the continuum version of the model, with 5 now
aggregated to the area. Moments 4 and 6 are what the granular structure adds.

\paragraph{Notes on foreign parameters.}
\label{app:model-foreign}
We do not have the data to estimate parameters relating to foreign firms. Instead we take the
share of foreign inputs in sourcing as an exogenous parameter used in the simulated data, so the
vector of parameters to be estimated excludes region $R+1$. Moments 1--4 and 6 do not involve
$R+1$. For moment 5, with $w_{R+1,s}$ the share of foreign inputs in sector $s$'s input
purchases,
\begin{equation}
\tilde\gamma_{as}=(1-w_{R+1,s})\sum_{r'\in\mathcal{R}_{a}}\sum_{r}
\frac{\sum_{r''\neq R+1}X_{r''rs}}{\sum_{r''\neq R+1}\sum_{r}X_{r''rs}}\,\tilde\gamma_{r'rs},
\qquad
\tilde\gamma_{r'rs}=\frac{T_{a(r')s}(w_{r's}\tau_{r'rs})^{-\theta}}{\tilde\Phi_{rs}},
\end{equation}
with $\tilde\Phi_{rs}=\sum_{r'=1}^{R}T_{a(r')s}(w_{r's}\tau_{r'rs})^{-\theta}$. Foreign
competition still enters the win probabilities through $\Phi_{rs}$ in \eqref{eq:app-Phi}; what
is not estimated is foreign comparative advantage.

\paragraph{Calibrated parameters.}
\label{app:model-calibration}
The remaining parameters $\{\varepsilon,\theta,\lambda,\nu,\{\nu_{s}\},\{w_{r's}\}\}$ are
calibrated. We use the estimates from Fontagn\'e et al.\ (2022) for $\varepsilon$: the sales
elasticity for motor vehicles and aerospace is respectively $-9$ and $-16$. The shape $\theta$
of the Fr\'echet distribution is calibrated following Antr\`as et al.\ (2017). Elasticities of
substitution in production follow the literature, notably Baqaee and Farhi (2019): $\lambda=0.5$
for the elasticity between labor and intermediate inputs and $\nu=0.9$ across sectors. Local
wages are calibrated using labor cost information from employer--employee data in 2019
(DADS/BTS). Note that Proposition~\ref{prop:firm}(c) makes $\theta$ over-identified: the
coefficient on log own productivity in moment 4 should equal $-\theta$, and we report it.

\subsection{Estimation}
\label{app:model-estimation}

\paragraph{Simulation algorithm.} Given the calibrated elasticities and a candidate parameter
vector, we simulate the productivity champions of every (region, variety) pair, solve for all
bilateral trade flows --- from each potential supplier to each downstream region --- and
construct the simulated moments. All six blocks are computed from a \emph{single} Ricardian
solve on one set of draws. The SMM estimator is
\begin{equation}
\hat\mu=\argmin_{\mu}\;\big(\tilde M(\mu)-\hat M\big)'W\big(\tilde M(\mu)-\hat M\big),
\label{eq:app-smm}
\end{equation}
where $\tilde M$ (resp.\ $\hat M$) denotes the vector of theoretical (resp.\ empirical) moments
and $W$ is a weighting matrix. Optimal weights assign lower importance to imprecisely estimated
moments.

The variety count never enters the simulation. It is confined to two objects, both closed form,
and that confinement is the content of Proposition~\ref{prop:firm} together with:

\begin{lemma}[$q_{r's}$ is free of $N_{s}$]
\label{lem:qfree}
Winning a variety is the comparison
$w_{r's}\tau_{r'rs}/\varphi_{r'\rho s}<\min_{r''\neq r'}w_{r''s}\tau_{r''rs}/\varphi_{r''\rho s}$,
so $q_{r's}=q_{r's}(T_{\cdot s},w,\tau,\theta)$ and nothing else. The variety count, the price
index, downstream costs and expenditures do not appear.
\end{lemma}

The economics is that winning is a statement about \emph{relative costs}, whereas $N_{s}$ and
$P_{rs}$ are aggregators built after the winners are known: expenditure allocates value across
winners, it does not choose them. The one condition that would break the lemma is endogenous
upstream wages, since $w$ enters the comparison directly; wages are taken as data.

Together the two results say that $N_{s}$ reaches the moment vector only through
$\bar G_{s}(0)$ and through $\widehat N_{s}$ itself, and that both are functions of $q_{r's}$
alone. Writing $k_{r's}$ for the number of the $m$ simulated varieties won somewhere by region
$r'$ --- by Lemma~\ref{lem:qfree} an estimate of $q_{r's}$ valid whatever the variety count ---
the count moment is evaluated by the unbiased combinatorial estimator
\begin{equation}
\bar G_{s}(n)=\frac{1}{|\mathcal{R}^{+}_{s}|}\sum_{r'\in\mathcal{R}^{+}_{s}}
\frac{\binom{m-k_{r's}}{n}}{\binom{m}{n}},
\qquad
\mathbb{E}\left[\frac{\binom{m-k}{n}}{\binom{m}{n}}\right]=(1-q)^{n}\ \ \text{for } n\le m,
\label{eq:app-Gclosed}
\end{equation}
which is the probability that $n$ varieties drawn without replacement from the $m$ simulated
ones were all lost. We use \eqref{eq:app-Gclosed} rather than the plug-in $(1-\hat q_{r's})^{n}$
because the latter is convex in $\hat q$ and therefore biased upward. Cross-cell dependence ---
two regions in a sector face the same variety draws --- does not affect either expression, since
an expectation of an average is the average of expectations whatever the dependence. No economy
of size $N_{s}$ is ever drawn and no replication is needed.

\begin{remark}[What is given up]
\label{rem:binding}
Moment 4 is then the coefficient of a regression run on the full simulated sample, i.e.\ the
continuum limit $\beta(\mathbb{E}[y])$ rather than $\mathbb{E}[\hat\beta(N_{s})]$. The finite-sample bias of the
auxiliary cloglog --- which an average of $\hat\beta$ over realized economies of size $N_{s}$
would have cancelled against the same bias in the data --- is therefore not differenced out. It
is a measurable quantity, of the order of $2$--$4\%$ of $\beta$ at data-consistent group sizes,
and we report it alongside $\hat\alpha$. The alternative is not free: it would treat simulation
noise on one block by replication while every other block relies on the accuracy of the draw
design, and it would sit oddly beside Appendix~\ref{app:model-price}, where the value block is
already evaluated at the certainty-equivalent index rather than on a realized economy.
\end{remark}

\paragraph{Concentrating out the variety count.} $N_{s}$ is estimated by the same weighted
objective as every other parameter, concentrated rather than searched jointly:
\begin{equation}
\widehat N(\mu_{-N})=\argmin_{N\in\prod_{s}[N^{\mathrm{lo}}_{s},N^{\mathrm{hi}}_{s}]\cap\mathbb{N}}
\;r(\mu_{-N},N)'W\,r(\mu_{-N},N),
\label{eq:app-profile}
\end{equation}
with the bounds of \eqref{eq:app-bounds}. This is a profiled extremum estimator: $\widehat N_{s}$
is an integer at every evaluation, so integrality is never relaxed and never rounded, and by
Lemma~\ref{lem:qfree} the inner minimization needs no re-simulation. Since $\bar G_{s}(0;n)$ is
closed form and decreasing in $n$, \eqref{eq:app-profile} reduces to a monotone integer
bisection. Clamping at either bound is informative rather than benign: the upper bound binding
means the model cannot generate enough sparsity even when every variety is sourced from a single
origin, which is a rejection signal for the mechanism rather than a numerical nuisance.

\paragraph{Concentrating out comparative advantage.} Given
$(\alpha,\Omega_{L},\{\Omega_{s}\},\{A_{r}\})$, the area-level comparative advantage that
reproduces the observed sourcing shares \eqref{eq:app-gamma-area} is the solution of a
matrix-scaling problem. We solve it by the damped multiplicative iteration
\begin{equation}
T^{+}_{as}=T_{as}\left(\frac{\hat{\tilde\gamma}_{as}}{\tilde\gamma_{as}(T)}\right)^{\varsigma},
\qquad \varsigma\in(0,1],
\label{eq:app-sinkhorn}
\end{equation}
renormalizing to the reference area of each sector after every pass, and recomputing the
general-equilibrium expenditure weights that enter \eqref{eq:app-gamma-area}.

It is useful to separate the two things this iteration does, because only the first is covered
by a standard argument.

\emph{Step one: the scaling problem at fixed expenditure.} Hold the expenditure weights
$\omega_{rs}=X_{rs}/\sum_{r}X_{rs}$ of \eqref{eq:app-gamma-area} fixed. Then
$\tilde\gamma_{as}(T)$ is a ratio of positive linear forms in $T$: writing
$B_{ar}=\sum_{r'\in\mathcal{R}_{a}}(w_{r's}\tau_{r'rs})^{-\theta}$ for the (strictly positive)
bilateral access matrix, $\tilde\gamma_{as}(T)=\sum_{r}\omega_{rs}T_{as}B_{ar}/\sum_{a'}T_{a's}B_{a'r}$.
Recovering $T$ from target shares is then the classical matrix-scaling (Sinkhorn--Knopp)
problem, and \eqref{eq:app-sinkhorn} with $\varsigma=1$ is its standard iteration. Because
$B$ has full support --- every area can serve every buyer at finite cost --- the associated
positive linear map contracts the Hilbert projective metric with a modulus
$\kappa_{S}=\tanh(\Delta/4)<1$, where $\Delta$ is the projective diameter of the image cone
(Birkhoff--Hopf). Two consequences: the iteration converges geometrically from any strictly
positive start, and \textbf{the fixed point is unique up to the per-sector scale}, which is
exactly the gauge the reference normalization fixes. Given trade costs and expenditures, there
is one and only one vector of comparative advantages consistent with the observed area shares.

\emph{Step two: the general-equilibrium feedback.} Expenditure is not fixed. A change in $T$
moves the price indices $p_{rs}$, hence $c_{r}$, hence downstream sales $y_{r}$ and the weights
$\omega_{rs}$. The map actually iterated is the composition $\Psi=\mathcal{S}\circ\omega$, whose
Jacobian decomposes as
\begin{equation}
D\Psi = \underbrace{D\mathcal{S}\big|_{\omega\text{ fixed}}}_{\text{scaling channel}}
\;+\;\underbrace{J_{\mathrm{GE}}}_{\text{expenditure channel}},
\qquad
J_{\mathrm{GE}}\equiv D\Psi-D\mathcal{S}\big|_{\omega\text{ fixed}} .
\label{eq:app-jacobian-split}
\end{equation}
Birkhoff's theorem bounds the first term by $\kappa_{S}$ but says nothing about the second, so a
sufficient condition for $\Psi$ to be a contraction is
\begin{equation}
\kappa_{S}+\big\|J_{\mathrm{GE}}\big\|<1 .
\label{eq:app-contraction}
\end{equation}
The size of $J_{\mathrm{GE}}$ is governed by how far the technology is from Cobb--Douglas: the
expenditure weights are invariant to $T$ when $\lambda=\nu=1$, so $\|J_{\mathrm{GE}}\|$ vanishes
with $|1-\lambda|$ and $|1-\nu|$. Our calibration is not close to Cobb--Douglas, so
\eqref{eq:app-contraction} is an empirical question rather than an assumption, and we verify it
numerically at the estimated parameters by finite differences on the one-step map. We find
$\kappa_{S}\approx0.005$ and $\|J_{\mathrm{GE}}\|\approx0.015$, so the sufficient condition
\eqref{eq:app-contraction} holds with a wide margin ($\approx0.02\ll1$); the measured spectral
radius of the full map is $\approx0.012$, and convergence to $10^{-9}$ takes a handful of
iterations.%
% NOTE: these three numbers were measured by the Phase-0 inversion gate
% (test/test_ge_inversion.jl). Re-run it under --granular=true --ca_level=aa to
% refresh them before circulating, since the earlier measurement was taken on the
% ZE-level continuum configuration.
\ We further confirm that the inversion recovers a known $T$ from its own implied
shares to machine precision, and that ten dispersed starting values converge to the same fixed
point.

This is a local, numerically verified statement at the estimated parameters, not a global
theorem: \eqref{eq:app-contraction} is sufficient, not necessary, and we do not characterize the
region of the parameter space over which it holds. Should it fail at some candidate parameter
vector, the iteration is not used blindly --- non-convergence is detected and the candidate
retains its warm-start $T$, so the outer optimizer sees a well-defined objective throughout.

Substituting $T^{*}(\alpha,\Omega,A)$ into \eqref{eq:app-smm} collapses the outer search from
several hundred parameters to a handful and makes the sourcing-share block just-identified along
the profiled manifold, so that the extensive-margin coefficient is fit by $(\alpha,\Omega,A)$
alone. Note that just-identification here is at the \emph{area} level: within an area, the split
of sourcing across regions remains over-identified by the distance profile, which is the content
of Assumption~\ref{as:geo}.

\paragraph{Weighting matrix.} In the SMM framework noise arises from two sources: sampling
variability in the data and simulation error. We distinguish between moments for which
bootstrapping is feasible (sourcing shares, extensive-margin coefficients and count moments) and
those for which it is not (input--output coefficients and the downstream sales distribution). We
therefore approximate the variance--covariance matrix of the data moments, $\Sigma_{\mathrm{data}}$,
with a block matrix that accounts for within-group interactions; the block corresponding to
non-bootstrappable moments is set to zero. Because the same firms generate the sourcing shares,
the extensive-margin coefficients and the counts, these three blocks are estimated by a
\emph{single joint} bootstrap, which is what delivers the cross-covariance blocks the efficient
weight matrix requires. The resampling unit is the supplier firm, pooled over the geography of
active areas, with the denominator of $\bar G_{s}(0)$ held fixed --- so that a region with one
or two firms can draw zero and flip into the empty set, which is precisely the variance the
count moment needs.

\paragraph{Optimization procedure.} The estimator is computed using Particle Swarm Optimization
\citep{kennedy1995}, which explores the parameter space through parallel evaluation, with
adaptive inertia weights declining from $0.9$ to $0.4$ over iterations and cognitive and social
parameters set to $1.5$ and $2.5$. The estimation proceeds in three steps.

\begin{enumerate}[topsep=3pt,itemsep=3pt,label=\arabic*.]
\item Starting values are set as follows: $\Omega_{L}$ and $\{\Omega_{s}\}$ at their empirical
counterparts from input--output tables; $\{A_{r}\}$ proportional to $\pi_{r}^{1/\varepsilon}$;
$\alpha$ at the extensive-margin regression estimate divided by $\theta$; and $\{T_{as}\}$ at the
matrix-scaling solution \eqref{eq:app-sinkhorn} evaluated at that $\alpha$. We then run the
optimizer over the profiled parameter vector using the identity weighting matrix, refining with
an alternating block scheme in which each iteration sequentially optimizes $\{\Omega_{L},\Omega_{s}\}$,
then $\alpha$, then $\{A_{r}\}$, with a search radius progressively reduced across iterations.
This yields $\hat\mu_{1}$.
\item Using $\hat\mu_{1}$ we estimate the variance--covariance matrix of the simulated moments,
$\Sigma_{\mathrm{sim}}$, from repeated re-simulation with independent draws, and construct
$W=(\Sigma_{\mathrm{data}}+\Sigma_{\mathrm{sim}})^{-1}$.
\item We re-estimate the model with the updated weighting matrix, initializing at $\hat\mu_{1}$,
to obtain $\hat\mu_{2}$.
\end{enumerate}

\paragraph{Standard errors.} Let $G=\partial\tilde M/\partial\mu$ evaluated at $\hat\mu_{2}$ by
finite differences on common random numbers, and $\Omega=\Sigma_{\mathrm{data}}+\Sigma_{\mathrm{sim}}$.
We report the efficient variance $(G'WG)^{-1}$ alongside the sandwich form
$(G'WG)^{-1}G'W\Omega WG(G'WG)^{-1}$, which coincide when $W=\Omega^{-1}$; the ratio is a
diagnostic on the weighting matrix. Because $T$ is concentrated out through
\eqref{eq:app-sinkhorn} rather than searched, the Jacobian is taken along the profiled manifold
--- perturbing $\alpha$ and letting $T^{*}(\alpha)$ follow --- so the reported standard error on
$\alpha$ is a total derivative, and confidence intervals for $T$ are obtained by the delta
method, propagating both the uncertainty in $\hat\alpha$ and the sampling uncertainty in the
sourcing-share targets that \eqref{eq:app-sinkhorn} inverts.

\subsection{The price index, love of variety, and the normalization}
\label{app:model-price}

The one place where granularity meets the value block requires an explicit assumption. For
variety $\rho$ in market $(r,s)$, origin $r'$'s champion cost satisfies
$\Pr(c_{r'\rho s}(r)>x)=\exp(-T_{a(r')s}(w_{r's}\tau_{r'rs})^{-\theta}x^{\theta})$, so by
independence across origins the minimum $c^{*}_{\rho}$ is Weibull:
\begin{equation}
\Pr\big(c^{*}_{\rho}>x\big)=\exp\big(-\Phi_{rs}x^{\theta}\big),
\qquad
\mathbb{E}\big[(c^{*}_{\rho})^{k}\big]=\Phi_{rs}^{-k/\theta}\,\Gamma\Big(1+\tfrac{k}{\theta}\Big),
\quad k>-\theta .
\end{equation}
With the CES aggregate taken as a sum,
$\mathbb{E}[P_{rs}^{1-\nu_{s}}]=N_{s}\Gamma(1+\tfrac{1-\nu_{s}}{\theta})\Phi_{rs}^{-\frac{1-\nu_{s}}{\theta}}$
up to a constant, so the certainty-equivalent index
$\widetilde P_{rs}=(\mathbb{E}[P_{rs}^{1-\nu_{s}}])^{1/(1-\nu_{s})}$ satisfies
\begin{equation}
\widetilde P_{rs}=N_{s}^{\frac{1}{1-\nu_{s}}}\times
\underbrace{\Gamma\Big(1+\tfrac{1-\nu_{s}}{\theta}\Big)^{\frac{1}{1-\nu_{s}}}\Phi_{rs}^{-1/\theta}}_{N_{s}\text{-free}} .
\label{eq:app-loveofvariety}
\end{equation}
Since $1-\nu_{s}<0$ the factor is decreasing in $N_{s}$: more varieties lower the sector price
index. Rescaling $\widetilde T_{as}=T_{as}N_{s}^{\kappa}$ with $\kappa=\theta/(1-\nu_{s})$ maps
$\Phi_{rs}\mapsto N_{s}^{\kappa}\Phi_{rs}$ and cancels the two powers exactly, leaving
$\widetilde P_{rs}$ invariant to $N_{s}$. Equivalently, taking the CES aggregate as an
\emph{average} rather than a sum --- as in \eqref{eq:app-ces} --- implements the same
normalization. It is consistent because it acts on a gauge the sourcing shares do not see: the
share \eqref{eq:app-share} is a ratio of $T$'s within a sector, so the extensive margin, the
count moment and the within-area contrast are all unchanged. It is also a gauge the estimator
already fixes, by normalizing one reference area per sector.

Two caveats. First, $P_{rs}$ is a \emph{convex} transform of the average in \eqref{eq:app-ces},
so the realized index exceeds its certainty-equivalent value by Jensen's inequality, and the gap
closes only as fast as the average concentrates. At the calibration the regularity condition
$\theta>2(\nu_{s}-1)$ holds with little slack, and the gap is of the order of $6\%$ at
$N_{s}\approx10^{2}$, falling to $1\%$ at $N_{s}\approx10^{3}$. The estimator adopts the
certainty-equivalent reading --- downstream firms optimize against
$\mathbb{E}[P^{1-\nu_{s}}]^{1/(1-\nu_{s})}$, neglecting granular price randomness, the
general-equilibrium embedding of \citet{lmm2023}. Most of that $6\%$ would not have been
identifying variation in any case: a price-level shift is observationally equivalent to a
technology shift, so the sector technology terms absorb it and only the cross-sector pattern of
$N_{s}$ survives as signal. And by Lemma~\ref{lem:qfree} none of it touches $q_{r's}$ or the
extensive margin, so the identification of $\alpha$ and $N_{s}$ is unaffected either way.

Second, the assumption must not survive into counterfactuals. The elasticity
$\partial\ln\widetilde P_{rs}/\partial\ln N_{s}=1/(1-\nu_{s})<0$ is a genuine gain from variety
and a live propagation channel: a shock that removes suppliers in one region raises $P_{rs}$
downstream through exactly this elasticity, and that channel is absent from the continuum model.
Any counterfactual that moves the variety set must map back to the physical $T$ and use the
realized index.

\subsection{Proofs}
\label{app:model-proofs}

\paragraph{Equation \eqref{eq:app-champion}.} The number of firms in $(r',s,\rho)$ with
productivity above $z$ is Poisson with mean $T_{r's}z^{-\theta}$, so the probability that none
exceeds $z$ --- equivalently that the maximum is below $z$ --- is $\exp(-T_{r's}z^{-\theta})$,
the $\mathrm{Fr\acute{e}chet}(T_{r's},\theta)$ CDF. \hfill$\square$

\paragraph{Lemma~\ref{lem:fe}.} (i) With $\Pr(y=1)=1-\exp(-\exp(\phi_{g}+\beta x))$, a group in
which $y\equiv0$ contributes
$\ell_{g}=\sum_{i}\ln\exp(-e^{\phi_{g}+\beta x_{i}})=-e^{\phi_{g}}\sum_{i}e^{\beta x_{i}}$. This
is increasing in $-\phi_{g}$ with supremum $0$ as $\phi_{g}\to-\infty$, at which
$\partial\ell_{g}/\partial\beta=-e^{\phi_{g}}\sum_{i}x_{i}e^{\beta x_{i}}\to0$. The group
therefore drops from the concentrated likelihood.

(ii) By Frisch--Waugh--Lovell,
$\hat\beta=\sum_{g,i}\omega\tilde x\tilde y/\sum_{g,i}\omega\tilde x^{2}$ with $\tilde{\cdot}$
the weighted within-group demeaning. A group with $y\equiv0$ has $\tilde y\equiv0$, so it
contributes zero to the numerator and $\sum\omega\tilde x^{2}>0$ to the denominator. Splitting
the denominator between $\mathcal{A}^{+}$ and its complement gives
$\hat\beta=\varkappa\hat\beta^{+}$. \hfill$\square$

\paragraph{Proposition~\ref{prop:firm}.} Condition on the firm's productivity $z$; its delivered
cost to buyer $a$ is $c=w_{r's}\tau_{r'as}/z$. Each competitor $r''$ has champion cost
$w_{r''s}\tau_{r''as}/\varphi_{r''}$ with
$\varphi_{r''}\sim\mathrm{Fr\acute{e}chet}(T_{a(r'')s},\theta)$, so
$\Pr(w_{r''s}\tau_{r''as}/\varphi_{r''}>x)=\exp(-T_{a(r'')s}(w_{r''s}\tau_{r''as})^{-\theta}x^{\theta})$.
By independence across competitors the minimum competing cost $M$ satisfies
$\Pr(M>x)=\exp(-\Phi^{-r'}_{as}x^{\theta})$. The firm supplies iff $M>c$, giving
\eqref{eq:app-firm-win}. Taking the complement,
$\Pr(\text{no supply}\mid z)=1-\exp(-\exp(\eta))$ with
$\eta=\ln\Phi^{-r'}_{as}+\theta\ln(w_{r's}\tau_{r'as})-\theta\ln z$, and
$\tau_{r'as}=d_{r'a}^{\alpha}$ gives \eqref{eq:app-firm-cloglog}. No step involves the variety
count. \hfill$\square$

\paragraph{Lemma~\ref{lem:qfree}.} Origin $r'$ wins variety $\rho$ for buyer $r$ iff
$w_{r's}\tau_{r'rs}/\varphi_{r'\rho s}<\min_{r''\neq r'}w_{r''s}\tau_{r''rs}/\varphi_{r''\rho s}$,
an event measurable with respect to $\{\varphi_{\cdot\rho s}\}$, $w$ and $\tau$ alone; the
$\varphi$ have distribution governed by $(T_{\cdot s},\theta)$. The win-somewhere event is the
union over $r$ of those events, hence measurable with respect to the same objects. Neither the
variety count nor any price or expenditure aggregate enters. \hfill$\square$

\paragraph{Unbiasedness of \eqref{eq:app-Gclosed}.} Let $k$ be the number of the $m$ simulated
varieties won by $r'$, each won independently with probability $q$. Draw $n\le m$ of the $m$
varieties uniformly without replacement; conditional on $k$, the probability that all $n$ are
losses is $\binom{m-k}{n}/\binom{m}{n}$. Marginalizing over the draws, each of the $n$ selected
varieties is a loss independently with probability $1-q$, so the unconditional probability is
$(1-q)^{n}$. Hence $\mathbb{E}[\binom{m-k}{n}/\binom{m}{n}]=(1-q)^{n}$ exactly, with the convention
$\binom{a}{n}=0$ for $a<n$. \hfill$\square$
 from the main file. Requires amsmath, amssymb, amsthm,
% booktabs, enumitem and natbib. Add to the preamble if not already present:
%
%   \newtheorem{proposition}{Proposition}
%   \newtheorem{lemma}{Lemma}
%   \newtheorem{assumption}{Assumption}
%   \newtheorem{remark}{Remark}
%   \DeclareMathOperator*{\argmin}{arg\,min}
%
% Labels used by core_model.tex: app:model-derivation, app:model-foreign,
% app:model-calibration, app:model-prop, app:model-zeros, app:model-bounds.
% =====================================================================

\section{Details on the Structural Model}
\label{app:model}

\subsection{Derivation}
\label{app:model-derivation}

\paragraph{Setup.} The economy is divided into $R+1$ regions, populated by two types of firms:
producers of a final good and producers of intermediate goods used as inputs in the production
of the final good. The $(R+1)$th region represents the rest of the world. The analysis primarily
focuses on the matching of firms across domestic regions, controlling for competition from
foreign firms using their shares in total sales as reference. We assume that each domestic
region hosts at most one representative firm in the downstream sector, identified by its region
of origin. Intermediate goods are produced across $S$ sectors, with each region potentially
hosting many suppliers in each sector. We define a market $(s,r')$ as the collection of firms
located in region $r'$ that belong to upstream sector $s$.

Domestic regions are partitioned into attraction areas,
$\mathcal{R}=\bigsqcup_{a\in\mathcal{A}}\mathcal{R}_{a}$, where $\mathcal{R}_{a}$ is the set of
regions whose nearest downstream producer is the anchor of area $a$, and $a(r')$ denotes the
area of region $r'$.

\paragraph{Demand for final manufacturing products.} Firms in the downstream sector each own the
blueprint to produce a unique differentiated variety of the final good. The market for the final
good is characterized by monopolistic competition and iso-elastic demand functions:
\begin{equation}
q_{r}=\left(\frac{p_{r}}{P}\right)^{\varepsilon-1}\frac{E}{P}\,\delta_{r},
\end{equation}
where $q_{r}$ denotes the quantity demanded of the variety produced in region $r$ at price
$p_{r}$, given nominal consumption $E$ and the downstream sector's price index $P$. The
parameter $\varepsilon$ measures the price elasticity of sales, and $\delta_{r}$ is an exogenous
component of demand, varying across downstream producers and later used to simulate demand
shocks and their diffusion along production linkages. We abstract from foreign demand and assume
that the domestic market for final consumption goods is perfectly integrated.

\paragraph{Technology in the downstream sector.} Production of final goods requires the assembly
of intermediates, combined with labor, under a nested CES structure:
\begin{equation}
y_{r}=A_{r}\Big[\Omega_{L}^{\frac{1}{\lambda}}\ell_{r}^{\frac{\lambda-1}{\lambda}}
+(1-\Omega_{L})^{\frac{1}{\lambda}}\big(q^{I}_{r}\big)^{\frac{\lambda-1}{\lambda}}\Big]^{\frac{\lambda}{\lambda-1}},
\end{equation}
where $\ell_{r}$ is labor demand and $q^{I}_{r}$ the firm's intermediate input consumption. The
parameter $\Omega_{L}$ governs the share of labor in total production costs, while $\lambda$ is
the elasticity of substitution between labor and intermediates. Firm productivity $A_{r}$ is
exogenous. Intermediate consumption is itself a nested CES aggregator across sectors and, within
each sector, across varieties:
\begin{equation}
q^{I}_{r}=\Big[\sum_{s=1}^{S}\Omega_{s}^{\frac{1}{\nu}}q_{rs}^{\frac{\nu-1}{\nu}}\Big]^{\frac{\nu}{\nu-1}},
\qquad
q_{rs}=\Big[\frac{1}{N_{s}}\sum_{\rho=1}^{N_{s}}q_{rs}(\rho)^{\frac{\nu_{s}-1}{\nu_{s}}}\Big]^{\frac{\nu_{s}}{\nu_{s}-1}},
\label{eq:app-ces}
\end{equation}
where $\Omega_{s}$ governs the share of sector $s$ in intermediate consumption, $\nu$ is the
elasticity of substitution across sectors and $\nu_{s}$ that across varieties within a sector.

The only change relative to the continuum formulation is the inner aggregator: the integral over
$j\in[0,1]$ is replaced by an average over the $N_{s}$ varieties. Writing it as an
\emph{average} rather than a sum is a normalization of the love-of-variety level, discussed and
justified in Appendix~\ref{app:model-price}; it leaves every sourcing share, and hence the
entire extensive margin, unchanged.

Given the nested CES structure, the nominal demand for intermediate products is
\begin{equation}
x_{rs}(\rho)=\Omega_{s}(1-\Omega_{L})
\left(\frac{p_{rs}(\rho)}{p_{rs}}\right)^{1-\nu_{s}}
\left(\frac{p_{rs}}{p_{r}}\right)^{1-\nu}
\left(\frac{p_{r}}{c_{r}}\right)^{1-\lambda}\frac{c_{r}y_{r}}{A_{r}},
\end{equation}
with $p_{rs}\equiv\big[\tfrac{1}{N_{s}}\sum_{\rho}p_{rs}(\rho)^{1-\nu_{s}}\big]^{\frac{1}{1-\nu_{s}}}$
the cost index for intermediates in sector $s$,
$p_{r}\equiv\big[\sum_{s}\Omega_{s}p_{rs}^{1-\nu}\big]^{\frac{1}{1-\nu}}$ the cost index for
intermediates, and
$c_{r}\equiv\big[\Omega_{L}w_{r}^{1-\lambda}+(1-\Omega_{L})p_{r}^{1-\lambda}\big]^{\frac{1}{1-\lambda}}$
the unit cost of the firm. In equilibrium $y_{r}=q_{r}$.

\paragraph{Technology in upstream sectors.} Intermediate goods are produced with labor using a
constant-returns-to-scale technology. Wages in sector $s$ and region $r'$ are denoted $w_{r's}$
and taken as given. As in Eaton and Kortum (2002) we assume that varieties produced within and
across regions are perfect substitutes and that productivity draws are independent. The
departure from the continuum case is that each sector contains $N_{s}$ varieties.

\begin{assumption}[Granularity]
\label{as:granular}
Each upstream sector $s$ consists of $N_{s}\in\mathbb{N}$ horizontally differentiated varieties
$\rho=1,\dots,N_{s}$. Within each variety $\rho$ and each region $r'$ there is a continuum of
firms whose productivities above $z$ form a Poisson process with mean $T_{r's}z^{-\theta}$.
Draws are independent across varieties, regions and sectors.
\end{assumption}

The continuum of firms is what keeps the model tractable. The number of firms above $z$ in
$(r',s,\rho)$ is Poisson with mean $T_{r's}z^{-\theta}$, so the most productive one satisfies
\begin{equation}
\Pr\big(\varphi_{r'\rho s}\le z\big)=\Pr(\text{no firm above }z)=e^{-T_{r's}z^{-\theta}},
\label{eq:app-champion}
\end{equation}
that is, $\varphi_{r'\rho s}$ is exactly $\mathrm{Fr\acute{e}chet}(T_{r's},\theta)$. One
Fr\'echet draw per (region, variety) pair \emph{is} the champion of a continuum of firms, not an
approximation to it. Since the iceberg cost scales all firms of region $r'$ identically, that
champion is the unique region-$r'$ firm that ever serves any buyer sourcing $\rho$ from $r'$:
\textbf{a firm is a (region, variety) champion}, with delivered cost
$c_{r'\rho s}(r)=w_{r's}\tau_{r'rs}/\varphi_{r'\rho s}$. This is the ``fixed number of
independent Ricardian draws'' form of granular Eaton--Kortum \citep[cf.][]{eks2012,bejk2003}.

\begin{assumption}[Area-level comparative advantage]
\label{as:geo}
Comparative advantage is area-specific: $T_{r's}=T_{a(r')s}$ for every $r'\in\mathcal{R}_{a}$,
with one parameter per (sector, area). Productivity draws remain independent across regions ---
two regions in the same area share the Fr\'echet scale, not the realization of their champions.
\end{assumption}

Assumption~\ref{as:geo} does not make regions within an area interchangeable. Each runs its own
Ricardian competition for each variety, and two regions at different distances from a buyer have
different win probabilities. What it removes is unobserved technological heterogeneity within an
area, which is what allows the model to explain the same variation the extensive-margin
regression uses.

\paragraph{Equilibrium.} The resolution follows the same steps as in Eaton and Kortum (2002).
Given the Fr\'echet assumption, the minimum price at which suppliers from $(r',s)$ can serve
market $r$ is distributed
\begin{equation}
G_{r'rs}(p)=\Pr\Big(\frac{w_{r's}\tau_{r'rs}}{\varphi_{r'\rho s}}\le p\Big)
=1-e^{-\left[T_{a(r')s}(w_{r's}\tau_{r'rs})^{-\theta}\right]p^{\theta}},
\end{equation}
so the distribution of input prices for downstream firms in region $r$ is
\begin{equation}
G_{rs}(p)=1-\prod_{r'=1}^{R+1}\big[1-G_{r'rs}(p)\big]=1-e^{-\Phi_{rs}p^{\theta}},
\qquad
\Phi_{rs}=\sum_{a'}T_{a's}\sum_{r'\in\mathcal{R}_{a'}}(w_{r's}\tau_{r'rs})^{-\theta}.
\label{eq:app-Phi}
\end{equation}
The probability that a downstream firm in location $r$ picks a supplier in market $(r',s)$ is
thus
\begin{equation}
\gamma_{r'rs}=\frac{T_{a(r')s}(w_{r's}\tau_{r'rs})^{-\theta}}{\Phi_{rs}} .
\label{eq:app-share}
\end{equation}
Taking the perspective of the supplier, the probability of being the lowest-cost supplier in
market $(r,s)$, given location $r'$ and productivity $z$, is
\begin{equation}
\rho_{r'rs}(z)=e^{-\Phi_{rs}(w_{r's}\tau_{r'rs})^{\theta}z^{-\theta}}
=e^{-T_{a(r')s}\gamma^{-1}_{r'rs}z^{-\theta}} .
\label{eq:app-rho}
\end{equation}
Because the conditional distribution of prices is not origin-specific,
$G_{r'rs}(p)=G_{rs}(p)$, and $\gamma_{r'rs}$ is also the expected share of intermediates from
$r'$ purchased by a downstream firm in $r$. The expected cost indices are
\begin{equation}
p_{rs}=\Big[\Gamma\Big(\frac{\theta+1-\nu_{s}}{\theta}\Big)\Big]^{\frac{1}{1-\nu_{s}}}\Phi_{rs}^{-1/\theta},
\qquad
p_{r}=\Big[\sum_{s}\Omega_{s}p_{rs}^{1-\nu}\Big]^{\frac{1}{1-\nu}},
\qquad
c_{r}=\Big[\Omega_{L}w_{r}^{1-\lambda}+(1-\Omega_{L})p_{r}^{1-\lambda}\Big]^{\frac{1}{1-\lambda}},
\end{equation}
where the first expression is the certainty-equivalent index of Appendix~\ref{app:model-price}.
Bilateral flows are
\begin{equation}
X_{r'rs}=\gamma_{r'rs}\,x_{rs}
=\gamma_{r'rs}\,\Omega_{s}(1-\Omega_{L})
\left(\frac{p_{rs}}{p_{r}}\right)^{1-\nu}\left(\frac{p_{r}}{c_{r}}\right)^{1-\lambda}\frac{c_{r}y_{r}}{A_{r}} .
\end{equation}
Since we do not observe these bilateral flows directly, we cannot use this equation to estimate
the model's structural parameters and rely instead on simulated moments.

\subsection{The extensive margin and the two classes of zeros}
\label{app:model-zeros}

Across the $N_{s}$ varieties the win events are i.i.d., so with $K_{r'rs}$ the number of
sector-$s$ varieties buyer $r$ sources from $r'$,
\begin{equation}
K_{r'rs}\sim\mathrm{Bin}\big(N_{s},\gamma_{r'rs}\big),
\qquad
\Pr(K_{r'rs}=0)=(1-\gamma_{r'rs})^{N_{s}}>0 .
\end{equation}
Under a continuum this probability is $\lim_{N\to\infty}(1-\gamma_{r'rs})^{N}=0$ for every
origin with $T_{r's}>0$: the extensive margin is degenerate and the observed zeros are not a
model outcome. Finiteness is exactly what restores it.

\begin{remark}[Dependence across buyers]
\label{rem:dependence}
For a given variety the champion costs are common across buyers; only $\tau_{r'rs}$ differs. Win
events are therefore positively dependent across buyers, and the union probability $q_{r's}$ is
not the product of the per-market marginals. The per-market marginal above is nonetheless exact,
and $q_{r's}$ is evaluated by simulation. This matters because the survey records whether a firm
supplies the industry, not which buyer it serves: the empirical outcome is the union, and the
model counterpart is computed as the same union.
\end{remark}

Assumption~\ref{as:geo} sorts observed zeros into two economically distinct classes, and the
distinction determines what is estimated.

\begin{table}[htbp]
\centering
\small
\begin{tabular}{@{}p{0.21\textwidth}p{0.24\textwidth}p{0.19\textwidth}p{0.28\textwidth}@{}}
\toprule
\textbf{Zero} & \textbf{Treatment} & \textbf{Probability} & \textbf{Role}\\
\midrule
Whole attraction area $a$ has no supplier in $s$
& \emph{Structural}: $T_{as}=0$, $(s,a)$ outside the active set
& $1$ by construction
& None --- absorbed by the $a\times s$ fixed effect (Lemma~\ref{lem:fe}); no parameter estimated\\[4pt]
Region $r'$ inside an active area has no supplier
& \emph{Endogenous}: inherits $T_{as}>0$, own distance $d_{r'}$
& $(1-q_{r's})^{N_{s}}$
& \textbf{The identifying variation}: disciplined by distance alone\\
\bottomrule
\end{tabular}
\caption{The two classes of extensive-margin zeros. The second row is exactly the set of regions
that host firms but no supplier. Under region-level comparative advantage these had to be
assigned $T\equiv0$, making their zeros exogenous; under Assumption~\ref{as:geo} they inherit
their area's comparative advantage and their zeros become model predictions.}
\label{tab:app-hierarchy}
\end{table}

The reason the fixed effect dictates this split is the following.

\begin{lemma}[All-zero groups carry no information about the distance slope]
\label{lem:fe}
Index observations by firm $i$ within fixed-effect group $g=(a,s)$, with outcome $y_{i}$,
regressor $x_{i}$, weights $\omega_{i}>0$ and group effect $\phi_{g}$.
\begin{enumerate}[nosep,topsep=2pt,label=(\roman*)]
\item \emph{Complementary log-log.} The log-likelihood of a group in which $y\equiv0$ is
$\ell_{g}=-e^{\phi_{g}}\sum_{i}e^{\beta x_{i}}$, whose supremum over $\phi_{g}$ is attained as
$\phi_{g}\to-\infty$ with $\ell_{g}\to0$ and $\partial\ell_{g}/\partial\beta\to0$: the group
drops out of the concentrated likelihood and contributes exactly nothing to $\hat\beta$.
\item \emph{Linear probability model.} By Frisch--Waugh--Lovell an all-zero group has
$\tilde y\equiv0$, contributing zero to the numerator and a positive amount to the denominator,
so $\hat\beta=\varkappa\hat\beta^{+}$ with
$\varkappa=\big(\sum_{a\in\mathcal{A}^{+}}\sum\omega\tilde x^{2}\big)/\big(\sum_{a}\sum\omega\tilde x^{2}\big)\in(0,1]$,
where $\mathcal{A}^{+}$ is the set of areas hosting at least one supplier.
\end{enumerate}
\end{lemma}

Under the cloglog link --- the link the model implies --- $\alpha$ is therefore identified
exclusively from variation \emph{within} attraction areas that host at least one supplier. Under
a linear probability model the empty areas enter only as a mechanical dilution of the same
within-area signal. Either way, an area with no supplier at all carries no information about the
distance slope, which is why the model does not attempt to explain it.

\subsection{The firm-level reduced form}
\label{app:model-prop}

The object the estimation matches is the union event \eqref{eq:q}: whether a firm supplies the
downstream industry at all. That probability has no tractable closed form. The single-buyer case
does, and it is what gives the empirical cloglog its structural interpretation; we state it
first and then say precisely what survives when there are several buyers.

\begin{proposition}[Firm-level reduced form, single buyer]
\label{prop:firm}
Let a firm be a (region, variety) champion with productivity $z$ located in $r'$, and suppose it
faces a single downstream buyer, its area's anchor $a$. Writing
$\Phi^{-r'}_{as}=\sum_{r''\neq r'}T_{a(r'')s}(w_{r''s}\tau_{r''as})^{-\theta}$ for the
competitors' cost scale,
\begin{equation}
\Pr\big(\text{the firm supplies}\mid z\big)
=\exp\Big(-\Phi^{-r'}_{as}\big(w_{r's}\tau_{r'as}\big)^{\theta}z^{-\theta}\Big),
\label{eq:app-firm-win}
\end{equation}
so the probability that it does \emph{not} supply is a complementary log-log function
\begin{equation}
\Pr\big(\text{no supply}\mid z\big)=1-\exp\big(-\exp(\eta)\big),
\qquad
\eta=\underbrace{\ln\Phi^{-r'}_{as}}_{\text{area}\times\text{sector FE}}
+\theta\ln w_{r's}
+\theta\alpha\,\ln d_{r'a}
-\theta\,\ln z .
\label{eq:app-firm-cloglog}
\end{equation}
The variety count $N_{s}$ does not appear.
\end{proposition}

Three consequences, and they are the backbone of the estimation.

\begin{enumerate}[nosep,topsep=3pt,label=(\alph*)]
\item \textbf{The cloglog is the model's own auxiliary model.} A cloglog of the ``does not
supply'' indicator on log distance and log own productivity, with an area$\times$sector fixed
effect, is exactly \eqref{eq:app-firm-cloglog} when a firm faces one buyer, with distance
coefficient $+\theta\alpha$ and productivity coefficient $-\theta$.
\item \textbf{$\alpha$ and $N_{s}$ separate.} $N_{s}$ is absent from
\eqref{eq:app-firm-cloglog}: conditioning on the firm removes it. The regression identifies
$\alpha$; the count level identifies $N_{s}$; there is no residual coupling.
\item \textbf{The productivity coefficient is a free over-identifying test.} It equals
$-\theta$. Extending the moment vector to include it would let $\theta$ be estimated rather than
calibrated.
\end{enumerate}

\paragraph{Several buyers: what survives.} Our surveys record supply to the industry, not to a
named buyer, so the empirical outcome is the union event $\bigcup_{r}W_{r}$, where $W_{r}$
denotes winning buyer $r$. Its probability is \emph{not} $1-\prod_{r}[1-\rho_{r'rs}(z)]$: that
expression would require the $W_{r}$ to be independent, and they are not.

The dependence is positive and structural: conditional on $z$, the firm's delivered cost to
every buyer is the same number scaled by $\tau_{r'rs}$, so a firm cheap enough to beat the
competition in one market tends to beat it everywhere. The exact object is the
inclusion--exclusion sum
\begin{equation}
\Pr\Big(\bigcup_{r}W_{r}\,\Big|\,z\Big)
=\sum_{\emptyset\neq\mathcal{S}\subseteq\{1,\dots,R\}}(-1)^{|\mathcal{S}|+1}
\exp\Big(-z^{-\theta}\max_{r\in\mathcal{S}}\big\{\Phi^{-r'}_{rs}(w_{r's}\tau_{r'rs})^{\theta}\big\}\Big),
\label{eq:app-inclexcl}
\end{equation}
where the maximum appears because $\cap_{r\in\mathcal{S}}W_{r}$ requires the firm to beat the
\emph{tightest} of the thresholds it faces. Equation~\eqref{eq:app-inclexcl} has $2^{R}-1$
terms. Many are redundant --- if the threshold for buyer $r$ dominates that for $r''$ then
$W_{r}\subseteq W_{r''}$ and the pair collapses --- and pruning the dominated buyers reduces the
sum to the non-dominated set, but that set is still large enough in our geography that
evaluating \eqref{eq:app-inclexcl} inside an optimizer is not practical. The independence
approximation $1-\prod_{r}(1-\rho_{r'rs}(z))$, which is what a closed-form estimator would have
to use, is a strict \emph{under}-statement of \eqref{eq:app-inclexcl} by positive dependence, and
the error is largest exactly where the extensive margin has content, at high $z$.

This is the reason the model is estimated by simulated rather than analytical method of moments.
Simulation evaluates the union exactly --- each variety's champion costs are drawn once and
compared against every buyer, so the dependence is reproduced by construction --- at the cost of
simulation noise, which the weighting matrix accounts for. We therefore treat
\eqref{eq:app-firm-cloglog} as an auxiliary model rather than an estimating equation: the same
cloglog is run on simulated and observed data, with the union as the outcome on both sides, and
the coefficients are matched. This is consistent for $\alpha$ whether or not the link is exact,
and it is why the distance coefficient enters the moment vector rather than being inverted
analytically.

The outcome convention is dictated by the unit of observation rather than chosen. Simulating the
exact object at $\theta=1$, $\alpha=0.30$, with firms as (region $\times$ variety) champions,
the ``does not supply'' outcome recovers $\hat\beta_{\ln d}=+0.327$ and
$\hat\beta_{\ln z}=-1.051$ against truths $\theta\alpha=0.30$ and $-\theta=-1.0$; the reversed
outcome recovers neither ($-0.426$ and $+1.452$).

\subsection{The count moment and bounds on the variety count}
\label{app:model-bounds}

The zeros are a cell-level object. With $\mathcal{R}^{+}_{s}$ the regions inside active
attraction areas, the empirical moment is the share of those regions hosting at most $n$
suppliers,
\begin{equation}
\bar G_{s}(n)=\frac{1}{|\mathcal{R}^{+}_{s}|}
\#\big\{r'\in\mathcal{R}^{+}_{s}:K_{r's}\le n\big\},
\qquad
\mathbb{E}\big[\bar G_{s}(0)\big]=\frac{1}{|\mathcal{R}^{+}_{s}|}\sum_{r'}(1-q_{r's})^{N_{s}} .
\label{eq:app-count}
\end{equation}
Writing $\mu_{r's}=N_{s}q_{r's}$ and using the Poisson approximation,
$\partial\bar G_{s}(K)/\partial N_{s}
=-|\mathcal{R}^{+}_{s}|^{-1}\sum_{r'}q_{r's}e^{-\mu_{r's}}\mu_{r's}^{K}/K!$, which peaks at
$K\approx\mu_{r's}$ and vanishes as $\bar G_{s}(K)\to1$: saturated thresholds carry no
identifying variation and a poorly estimated variance, so we use the empty-region share
$\bar G_{s}(0)$, which is both the most interpretable threshold and typically near-optimal.

Because the iceberg cost never changes which firm is cheapest within a region, each variety
generates one distinct supplying firm per winning origin. So if $N^{\mathrm{obs}}_{s}$ is the
number of distinct upstream firms observed to supply the downstream industry,
$N^{\mathrm{obs}}_{s}=\sum_{\rho}\#\{\text{origins winning }\rho\text{ somewhere}\}\in[N_{s},N_{s}R^{D}]$,
where $R^{D}$ is the number of downstream buyers, hence
\begin{equation}
\frac{N^{\mathrm{obs}}_{s}}{R^{D}}\;\le\;N_{s}\;\le\;N^{\mathrm{obs}}_{s} .
\label{eq:app-bounds}
\end{equation}
The bounds are tight: the upper when every variety is globally sourced from a single origin, the
lower when each of the $R^{D}$ buyers sources it from a different origin. Since
$\mathbb{E}[\sum_{r'}K_{r's}]=N_{s}\sum_{r'}q_{r's}$, matching the observed firm count gives a second,
closed-form estimate $N^{\mathrm{count}}_{s}=N^{\mathrm{obs}}_{s}/\sum_{r'}q_{r's}$ that lies
inside \eqref{eq:app-bounds} by construction --- a free over-identifying check.

\paragraph{Degrees of freedom.} Fix a sector, let $A^{+}=|\mathcal{A}^{+}_{s}|$ and let
$R^{\mathrm{act}}_{a}$ be the number of regions in area $a$ with a strictly positive observed
share. The $T$ block has $A^{+}-1$ free elements after the reference normalization, against
$\sum_{a}R^{\mathrm{act}}_{a}$ share moments, so Assumption~\ref{as:geo} converts what were free
parameters into
\begin{equation}
\sum_{a\in\mathcal{A}^{+}}\big(R^{\mathrm{act}}_{a}-1\big)
\quad\text{over-identifying restrictions,}
\end{equation}
one per active region beyond the first in each area. Each is a within-area statement: given the
area aggregate, the split across its regions must follow the distance profile. Under
region-level comparative advantage this count is zero.

\subsection{Structural estimation: the moments}
\label{app:model-moments}

We use six sets of theoretical moments, observed in the data using a combination of IO tables
and micro-level data.

\begin{enumerate}[topsep=3pt,itemsep=3pt,label=\arabic*.]
\item \textbf{Aggregate labor share at the level of the industry:}
\begin{equation}
\tilde\Omega_{L}\equiv\frac{\sum_{r}w_{r}\ell_{r}}{\sum_{r}c_{r}A_{r}^{-1}y_{r}}
=\sum_{r}\frac{c_{r}A_{r}^{-1}y_{r}}{\sum_{r}c_{r}A_{r}^{-1}y_{r}}\,\Omega_{L}
\left(\frac{w_{r}}{c_{r}}\right)^{1-\lambda}.
\end{equation}
In the model it is a sales-weighted average of firm-level labor shares. In the data we use the
share of labor costs in total costs from IO tables.

\item \textbf{Aggregate share of sector $s$ in the industry's input purchases:}
\begin{equation}
\tilde\Omega_{s}\equiv\frac{\sum_{r}p_{rs}q_{rs}}{\sum_{r}\sum_{s}p_{rs}q_{rs}}
=\sum_{r}\frac{p_{r}q^{I}_{r}}{\sum_{r}p_{r}q^{I}_{r}}\,\Omega_{s}
\left(\frac{p_{rs}}{p_{r}}\right)^{1-\nu}.
\end{equation}
In the data we use technical coefficients for the downstream industry from IO tables.

\item \textbf{The share of each region in downstream sales:}
\begin{equation}
\pi_{r}\equiv\frac{c_{r}A_{r}^{-1}y_{r}}{\sum_{r}c_{r}A_{r}^{-1}y_{r}}
=\left(\frac{p_{r}}{P}\right)^{-\varepsilon}
=\frac{(c_{r}A_{r}^{-1})^{-\varepsilon}}{\sum_{r}(c_{r}A_{r}^{-1})^{-\varepsilon}},
\end{equation}
recovered from sales data on the population of downstream producers.

\item \textbf{The extensive-margin distance coefficient.} The cloglog of
\eqref{eq:app-firm-cloglog}, estimated on the simulated firm population with the same
specification as Equation~(1) in the data: outcome ``does not supply'', regressors log distance
to the area anchor and log own productivity, area$\times$sector fixed effects.

\item \textbf{Area-aggregated sourcing shares.} The share of the downstream industry's purchases
of sector $s$ sourced from area $a$:
\begin{equation}
\tilde\gamma_{as}\equiv\sum_{r'\in\mathcal{R}_{a}}\frac{\sum_{r}X_{r'rs}}{\sum_{r}\sum_{r''}X_{r''rs}}
=\sum_{r'\in\mathcal{R}_{a}}\sum_{r}\frac{\sum_{r''}X_{r''rs}}{\sum_{r''}\sum_{r}X_{r''rs}}\,\gamma_{r'rs}.
\label{eq:app-gamma-area}
\end{equation}
In the data we reconstitute firm-level sales to the downstream industry using sales-share data,
$x_{i(r',s)d}=\mathbf{1}(\mathrm{Sup}_{i(r',s)})\,a^{D}_{d\,i(r',s)}\,x_{i(s,r')}$, sum over all
firms in the area and normalize by the sum over all firms from sector $s$. The aggregation runs
over \emph{every} region of the area, including those hosting no supplier, on both the model and
the data side.

\item \textbf{The count moment} $\bar G_{s}(0)$ of \eqref{eq:app-count}: the share of regions
inside active attraction areas that host no supplier of sector $s$.
\end{enumerate}

Moments 1--3 and 5 are the same objects as in the continuum version of the model, with 5 now
aggregated to the area. Moments 4 and 6 are what the granular structure adds.

\paragraph{Notes on foreign parameters.}
\label{app:model-foreign}
We do not have the data to estimate parameters relating to foreign firms. Instead we take the
share of foreign inputs in sourcing as an exogenous parameter used in the simulated data, so the
vector of parameters to be estimated excludes region $R+1$. Moments 1--4 and 6 do not involve
$R+1$. For moment 5, with $w_{R+1,s}$ the share of foreign inputs in sector $s$'s input
purchases,
\begin{equation}
\tilde\gamma_{as}=(1-w_{R+1,s})\sum_{r'\in\mathcal{R}_{a}}\sum_{r}
\frac{\sum_{r''\neq R+1}X_{r''rs}}{\sum_{r''\neq R+1}\sum_{r}X_{r''rs}}\,\tilde\gamma_{r'rs},
\qquad
\tilde\gamma_{r'rs}=\frac{T_{a(r')s}(w_{r's}\tau_{r'rs})^{-\theta}}{\tilde\Phi_{rs}},
\end{equation}
with $\tilde\Phi_{rs}=\sum_{r'=1}^{R}T_{a(r')s}(w_{r's}\tau_{r'rs})^{-\theta}$. Foreign
competition still enters the win probabilities through $\Phi_{rs}$ in \eqref{eq:app-Phi}; what
is not estimated is foreign comparative advantage.

\paragraph{Calibrated parameters.}
\label{app:model-calibration}
The remaining parameters $\{\varepsilon,\theta,\lambda,\nu,\{\nu_{s}\},\{w_{r's}\}\}$ are
calibrated. We use the estimates from Fontagn\'e et al.\ (2022) for $\varepsilon$: the sales
elasticity for motor vehicles and aerospace is respectively $-9$ and $-16$. The shape $\theta$
of the Fr\'echet distribution is calibrated following Antr\`as et al.\ (2017). Elasticities of
substitution in production follow the literature, notably Baqaee and Farhi (2019): $\lambda=0.5$
for the elasticity between labor and intermediate inputs and $\nu=0.9$ across sectors. Local
wages are calibrated using labor cost information from employer--employee data in 2019
(DADS/BTS). Note that Proposition~\ref{prop:firm}(c) makes $\theta$ over-identified: the
coefficient on log own productivity in moment 4 should equal $-\theta$, and we report it.

\subsection{Estimation}
\label{app:model-estimation}

\paragraph{Simulation algorithm.} Given the calibrated elasticities and a candidate parameter
vector, we simulate the productivity champions of every (region, variety) pair, solve for all
bilateral trade flows --- from each potential supplier to each downstream region --- and
construct the simulated moments. All six blocks are computed from a \emph{single} Ricardian
solve on one set of draws. The SMM estimator is
\begin{equation}
\hat\mu=\argmin_{\mu}\;\big(\tilde M(\mu)-\hat M\big)'W\big(\tilde M(\mu)-\hat M\big),
\label{eq:app-smm}
\end{equation}
where $\tilde M$ (resp.\ $\hat M$) denotes the vector of theoretical (resp.\ empirical) moments
and $W$ is a weighting matrix. Optimal weights assign lower importance to imprecisely estimated
moments.

The variety count never enters the simulation. It is confined to two objects, both closed form,
and that confinement is the content of Proposition~\ref{prop:firm} together with:

\begin{lemma}[$q_{r's}$ is free of $N_{s}$]
\label{lem:qfree}
Winning a variety is the comparison
$w_{r's}\tau_{r'rs}/\varphi_{r'\rho s}<\min_{r''\neq r'}w_{r''s}\tau_{r''rs}/\varphi_{r''\rho s}$,
so $q_{r's}=q_{r's}(T_{\cdot s},w,\tau,\theta)$ and nothing else. The variety count, the price
index, downstream costs and expenditures do not appear.
\end{lemma}

The economics is that winning is a statement about \emph{relative costs}, whereas $N_{s}$ and
$P_{rs}$ are aggregators built after the winners are known: expenditure allocates value across
winners, it does not choose them. The one condition that would break the lemma is endogenous
upstream wages, since $w$ enters the comparison directly; wages are taken as data.

Together the two results say that $N_{s}$ reaches the moment vector only through
$\bar G_{s}(0)$ and through $\widehat N_{s}$ itself, and that both are functions of $q_{r's}$
alone. Writing $k_{r's}$ for the number of the $m$ simulated varieties won somewhere by region
$r'$ --- by Lemma~\ref{lem:qfree} an estimate of $q_{r's}$ valid whatever the variety count ---
the count moment is evaluated by the unbiased combinatorial estimator
\begin{equation}
\bar G_{s}(n)=\frac{1}{|\mathcal{R}^{+}_{s}|}\sum_{r'\in\mathcal{R}^{+}_{s}}
\frac{\binom{m-k_{r's}}{n}}{\binom{m}{n}},
\qquad
\mathbb{E}\left[\frac{\binom{m-k}{n}}{\binom{m}{n}}\right]=(1-q)^{n}\ \ \text{for } n\le m,
\label{eq:app-Gclosed}
\end{equation}
which is the probability that $n$ varieties drawn without replacement from the $m$ simulated
ones were all lost. We use \eqref{eq:app-Gclosed} rather than the plug-in $(1-\hat q_{r's})^{n}$
because the latter is convex in $\hat q$ and therefore biased upward. Cross-cell dependence ---
two regions in a sector face the same variety draws --- does not affect either expression, since
an expectation of an average is the average of expectations whatever the dependence. No economy
of size $N_{s}$ is ever drawn and no replication is needed.

\begin{remark}[What is given up]
\label{rem:binding}
Moment 4 is then the coefficient of a regression run on the full simulated sample, i.e.\ the
continuum limit $\beta(\mathbb{E}[y])$ rather than $\mathbb{E}[\hat\beta(N_{s})]$. The finite-sample bias of the
auxiliary cloglog --- which an average of $\hat\beta$ over realized economies of size $N_{s}$
would have cancelled against the same bias in the data --- is therefore not differenced out. It
is a measurable quantity, of the order of $2$--$4\%$ of $\beta$ at data-consistent group sizes,
and we report it alongside $\hat\alpha$. The alternative is not free: it would treat simulation
noise on one block by replication while every other block relies on the accuracy of the draw
design, and it would sit oddly beside Appendix~\ref{app:model-price}, where the value block is
already evaluated at the certainty-equivalent index rather than on a realized economy.
\end{remark}

\paragraph{Concentrating out the variety count.} $N_{s}$ is estimated by the same weighted
objective as every other parameter, concentrated rather than searched jointly:
\begin{equation}
\widehat N(\mu_{-N})=\argmin_{N\in\prod_{s}[N^{\mathrm{lo}}_{s},N^{\mathrm{hi}}_{s}]\cap\mathbb{N}}
\;r(\mu_{-N},N)'W\,r(\mu_{-N},N),
\label{eq:app-profile}
\end{equation}
with the bounds of \eqref{eq:app-bounds}. This is a profiled extremum estimator: $\widehat N_{s}$
is an integer at every evaluation, so integrality is never relaxed and never rounded, and by
Lemma~\ref{lem:qfree} the inner minimization needs no re-simulation. Since $\bar G_{s}(0;n)$ is
closed form and decreasing in $n$, \eqref{eq:app-profile} reduces to a monotone integer
bisection. Clamping at either bound is informative rather than benign: the upper bound binding
means the model cannot generate enough sparsity even when every variety is sourced from a single
origin, which is a rejection signal for the mechanism rather than a numerical nuisance.

\paragraph{Concentrating out comparative advantage.} Given
$(\alpha,\Omega_{L},\{\Omega_{s}\},\{A_{r}\})$, the area-level comparative advantage that
reproduces the observed sourcing shares \eqref{eq:app-gamma-area} is the solution of a
matrix-scaling problem. We solve it by the damped multiplicative iteration
\begin{equation}
T^{+}_{as}=T_{as}\left(\frac{\hat{\tilde\gamma}_{as}}{\tilde\gamma_{as}(T)}\right)^{\varsigma},
\qquad \varsigma\in(0,1],
\label{eq:app-sinkhorn}
\end{equation}
renormalizing to the reference area of each sector after every pass, and recomputing the
general-equilibrium expenditure weights that enter \eqref{eq:app-gamma-area}.

It is useful to separate the two things this iteration does, because only the first is covered
by a standard argument.

\emph{Step one: the scaling problem at fixed expenditure.} Hold the expenditure weights
$\omega_{rs}=X_{rs}/\sum_{r}X_{rs}$ of \eqref{eq:app-gamma-area} fixed. Then
$\tilde\gamma_{as}(T)$ is a ratio of positive linear forms in $T$: writing
$B_{ar}=\sum_{r'\in\mathcal{R}_{a}}(w_{r's}\tau_{r'rs})^{-\theta}$ for the (strictly positive)
bilateral access matrix, $\tilde\gamma_{as}(T)=\sum_{r}\omega_{rs}T_{as}B_{ar}/\sum_{a'}T_{a's}B_{a'r}$.
Recovering $T$ from target shares is then the classical matrix-scaling (Sinkhorn--Knopp)
problem, and \eqref{eq:app-sinkhorn} with $\varsigma=1$ is its standard iteration. Because
$B$ has full support --- every area can serve every buyer at finite cost --- the associated
positive linear map contracts the Hilbert projective metric with a modulus
$\kappa_{S}=\tanh(\Delta/4)<1$, where $\Delta$ is the projective diameter of the image cone
(Birkhoff--Hopf). Two consequences: the iteration converges geometrically from any strictly
positive start, and \textbf{the fixed point is unique up to the per-sector scale}, which is
exactly the gauge the reference normalization fixes. Given trade costs and expenditures, there
is one and only one vector of comparative advantages consistent with the observed area shares.

\emph{Step two: the general-equilibrium feedback.} Expenditure is not fixed. A change in $T$
moves the price indices $p_{rs}$, hence $c_{r}$, hence downstream sales $y_{r}$ and the weights
$\omega_{rs}$. The map actually iterated is the composition $\Psi=\mathcal{S}\circ\omega$, whose
Jacobian decomposes as
\begin{equation}
D\Psi = \underbrace{D\mathcal{S}\big|_{\omega\text{ fixed}}}_{\text{scaling channel}}
\;+\;\underbrace{J_{\mathrm{GE}}}_{\text{expenditure channel}},
\qquad
J_{\mathrm{GE}}\equiv D\Psi-D\mathcal{S}\big|_{\omega\text{ fixed}} .
\label{eq:app-jacobian-split}
\end{equation}
Birkhoff's theorem bounds the first term by $\kappa_{S}$ but says nothing about the second, so a
sufficient condition for $\Psi$ to be a contraction is
\begin{equation}
\kappa_{S}+\big\|J_{\mathrm{GE}}\big\|<1 .
\label{eq:app-contraction}
\end{equation}
The size of $J_{\mathrm{GE}}$ is governed by how far the technology is from Cobb--Douglas: the
expenditure weights are invariant to $T$ when $\lambda=\nu=1$, so $\|J_{\mathrm{GE}}\|$ vanishes
with $|1-\lambda|$ and $|1-\nu|$. Our calibration is not close to Cobb--Douglas, so
\eqref{eq:app-contraction} is an empirical question rather than an assumption, and we verify it
numerically at the estimated parameters by finite differences on the one-step map. We find
$\kappa_{S}\approx0.005$ and $\|J_{\mathrm{GE}}\|\approx0.015$, so the sufficient condition
\eqref{eq:app-contraction} holds with a wide margin ($\approx0.02\ll1$); the measured spectral
radius of the full map is $\approx0.012$, and convergence to $10^{-9}$ takes a handful of
iterations.%
% NOTE: these three numbers were measured by the Phase-0 inversion gate
% (test/test_ge_inversion.jl). Re-run it under --granular=true --ca_level=aa to
% refresh them before circulating, since the earlier measurement was taken on the
% ZE-level continuum configuration.
\ We further confirm that the inversion recovers a known $T$ from its own implied
shares to machine precision, and that ten dispersed starting values converge to the same fixed
point.

This is a local, numerically verified statement at the estimated parameters, not a global
theorem: \eqref{eq:app-contraction} is sufficient, not necessary, and we do not characterize the
region of the parameter space over which it holds. Should it fail at some candidate parameter
vector, the iteration is not used blindly --- non-convergence is detected and the candidate
retains its warm-start $T$, so the outer optimizer sees a well-defined objective throughout.

Substituting $T^{*}(\alpha,\Omega,A)$ into \eqref{eq:app-smm} collapses the outer search from
several hundred parameters to a handful and makes the sourcing-share block just-identified along
the profiled manifold, so that the extensive-margin coefficient is fit by $(\alpha,\Omega,A)$
alone. Note that just-identification here is at the \emph{area} level: within an area, the split
of sourcing across regions remains over-identified by the distance profile, which is the content
of Assumption~\ref{as:geo}.

\paragraph{Weighting matrix.} In the SMM framework noise arises from two sources: sampling
variability in the data and simulation error. We distinguish between moments for which
bootstrapping is feasible (sourcing shares, extensive-margin coefficients and count moments) and
those for which it is not (input--output coefficients and the downstream sales distribution). We
therefore approximate the variance--covariance matrix of the data moments, $\Sigma_{\mathrm{data}}$,
with a block matrix that accounts for within-group interactions; the block corresponding to
non-bootstrappable moments is set to zero. Because the same firms generate the sourcing shares,
the extensive-margin coefficients and the counts, these three blocks are estimated by a
\emph{single joint} bootstrap, which is what delivers the cross-covariance blocks the efficient
weight matrix requires. The resampling unit is the supplier firm, pooled over the geography of
active areas, with the denominator of $\bar G_{s}(0)$ held fixed --- so that a region with one
or two firms can draw zero and flip into the empty set, which is precisely the variance the
count moment needs.

\paragraph{Optimization procedure.} The estimator is computed using Particle Swarm Optimization
\citep{kennedy1995}, which explores the parameter space through parallel evaluation, with
adaptive inertia weights declining from $0.9$ to $0.4$ over iterations and cognitive and social
parameters set to $1.5$ and $2.5$. The estimation proceeds in three steps.

\begin{enumerate}[topsep=3pt,itemsep=3pt,label=\arabic*.]
\item Starting values are set as follows: $\Omega_{L}$ and $\{\Omega_{s}\}$ at their empirical
counterparts from input--output tables; $\{A_{r}\}$ proportional to $\pi_{r}^{1/\varepsilon}$;
$\alpha$ at the extensive-margin regression estimate divided by $\theta$; and $\{T_{as}\}$ at the
matrix-scaling solution \eqref{eq:app-sinkhorn} evaluated at that $\alpha$. We then run the
optimizer over the profiled parameter vector using the identity weighting matrix, refining with
an alternating block scheme in which each iteration sequentially optimizes $\{\Omega_{L},\Omega_{s}\}$,
then $\alpha$, then $\{A_{r}\}$, with a search radius progressively reduced across iterations.
This yields $\hat\mu_{1}$.
\item Using $\hat\mu_{1}$ we estimate the variance--covariance matrix of the simulated moments,
$\Sigma_{\mathrm{sim}}$, from repeated re-simulation with independent draws, and construct
$W=(\Sigma_{\mathrm{data}}+\Sigma_{\mathrm{sim}})^{-1}$.
\item We re-estimate the model with the updated weighting matrix, initializing at $\hat\mu_{1}$,
to obtain $\hat\mu_{2}$.
\end{enumerate}

\paragraph{Standard errors.} Let $G=\partial\tilde M/\partial\mu$ evaluated at $\hat\mu_{2}$ by
finite differences on common random numbers, and $\Omega=\Sigma_{\mathrm{data}}+\Sigma_{\mathrm{sim}}$.
We report the efficient variance $(G'WG)^{-1}$ alongside the sandwich form
$(G'WG)^{-1}G'W\Omega WG(G'WG)^{-1}$, which coincide when $W=\Omega^{-1}$; the ratio is a
diagnostic on the weighting matrix. Because $T$ is concentrated out through
\eqref{eq:app-sinkhorn} rather than searched, the Jacobian is taken along the profiled manifold
--- perturbing $\alpha$ and letting $T^{*}(\alpha)$ follow --- so the reported standard error on
$\alpha$ is a total derivative, and confidence intervals for $T$ are obtained by the delta
method, propagating both the uncertainty in $\hat\alpha$ and the sampling uncertainty in the
sourcing-share targets that \eqref{eq:app-sinkhorn} inverts.

\subsection{The price index, love of variety, and the normalization}
\label{app:model-price}

The one place where granularity meets the value block requires an explicit assumption. For
variety $\rho$ in market $(r,s)$, origin $r'$'s champion cost satisfies
$\Pr(c_{r'\rho s}(r)>x)=\exp(-T_{a(r')s}(w_{r's}\tau_{r'rs})^{-\theta}x^{\theta})$, so by
independence across origins the minimum $c^{*}_{\rho}$ is Weibull:
\begin{equation}
\Pr\big(c^{*}_{\rho}>x\big)=\exp\big(-\Phi_{rs}x^{\theta}\big),
\qquad
\mathbb{E}\big[(c^{*}_{\rho})^{k}\big]=\Phi_{rs}^{-k/\theta}\,\Gamma\Big(1+\tfrac{k}{\theta}\Big),
\quad k>-\theta .
\end{equation}
With the CES aggregate taken as a sum,
$\mathbb{E}[P_{rs}^{1-\nu_{s}}]=N_{s}\Gamma(1+\tfrac{1-\nu_{s}}{\theta})\Phi_{rs}^{-\frac{1-\nu_{s}}{\theta}}$
up to a constant, so the certainty-equivalent index
$\widetilde P_{rs}=(\mathbb{E}[P_{rs}^{1-\nu_{s}}])^{1/(1-\nu_{s})}$ satisfies
\begin{equation}
\widetilde P_{rs}=N_{s}^{\frac{1}{1-\nu_{s}}}\times
\underbrace{\Gamma\Big(1+\tfrac{1-\nu_{s}}{\theta}\Big)^{\frac{1}{1-\nu_{s}}}\Phi_{rs}^{-1/\theta}}_{N_{s}\text{-free}} .
\label{eq:app-loveofvariety}
\end{equation}
Since $1-\nu_{s}<0$ the factor is decreasing in $N_{s}$: more varieties lower the sector price
index. Rescaling $\widetilde T_{as}=T_{as}N_{s}^{\kappa}$ with $\kappa=\theta/(1-\nu_{s})$ maps
$\Phi_{rs}\mapsto N_{s}^{\kappa}\Phi_{rs}$ and cancels the two powers exactly, leaving
$\widetilde P_{rs}$ invariant to $N_{s}$. Equivalently, taking the CES aggregate as an
\emph{average} rather than a sum --- as in \eqref{eq:app-ces} --- implements the same
normalization. It is consistent because it acts on a gauge the sourcing shares do not see: the
share \eqref{eq:app-share} is a ratio of $T$'s within a sector, so the extensive margin, the
count moment and the within-area contrast are all unchanged. It is also a gauge the estimator
already fixes, by normalizing one reference area per sector.

Two caveats. First, $P_{rs}$ is a \emph{convex} transform of the average in \eqref{eq:app-ces},
so the realized index exceeds its certainty-equivalent value by Jensen's inequality, and the gap
closes only as fast as the average concentrates. At the calibration the regularity condition
$\theta>2(\nu_{s}-1)$ holds with little slack, and the gap is of the order of $6\%$ at
$N_{s}\approx10^{2}$, falling to $1\%$ at $N_{s}\approx10^{3}$. The estimator adopts the
certainty-equivalent reading --- downstream firms optimize against
$\mathbb{E}[P^{1-\nu_{s}}]^{1/(1-\nu_{s})}$, neglecting granular price randomness, the
general-equilibrium embedding of \citet{lmm2023}. Most of that $6\%$ would not have been
identifying variation in any case: a price-level shift is observationally equivalent to a
technology shift, so the sector technology terms absorb it and only the cross-sector pattern of
$N_{s}$ survives as signal. And by Lemma~\ref{lem:qfree} none of it touches $q_{r's}$ or the
extensive margin, so the identification of $\alpha$ and $N_{s}$ is unaffected either way.

Second, the assumption must not survive into counterfactuals. The elasticity
$\partial\ln\widetilde P_{rs}/\partial\ln N_{s}=1/(1-\nu_{s})<0$ is a genuine gain from variety
and a live propagation channel: a shock that removes suppliers in one region raises $P_{rs}$
downstream through exactly this elasticity, and that channel is absent from the continuum model.
Any counterfactual that moves the variety set must map back to the physical $T$ and use the
realized index.

\subsection{Proofs}
\label{app:model-proofs}

\paragraph{Equation \eqref{eq:app-champion}.} The number of firms in $(r',s,\rho)$ with
productivity above $z$ is Poisson with mean $T_{r's}z^{-\theta}$, so the probability that none
exceeds $z$ --- equivalently that the maximum is below $z$ --- is $\exp(-T_{r's}z^{-\theta})$,
the $\mathrm{Fr\acute{e}chet}(T_{r's},\theta)$ CDF. \hfill$\square$

\paragraph{Lemma~\ref{lem:fe}.} (i) With $\Pr(y=1)=1-\exp(-\exp(\phi_{g}+\beta x))$, a group in
which $y\equiv0$ contributes
$\ell_{g}=\sum_{i}\ln\exp(-e^{\phi_{g}+\beta x_{i}})=-e^{\phi_{g}}\sum_{i}e^{\beta x_{i}}$. This
is increasing in $-\phi_{g}$ with supremum $0$ as $\phi_{g}\to-\infty$, at which
$\partial\ell_{g}/\partial\beta=-e^{\phi_{g}}\sum_{i}x_{i}e^{\beta x_{i}}\to0$. The group
therefore drops from the concentrated likelihood.

(ii) By Frisch--Waugh--Lovell,
$\hat\beta=\sum_{g,i}\omega\tilde x\tilde y/\sum_{g,i}\omega\tilde x^{2}$ with $\tilde{\cdot}$
the weighted within-group demeaning. A group with $y\equiv0$ has $\tilde y\equiv0$, so it
contributes zero to the numerator and $\sum\omega\tilde x^{2}>0$ to the denominator. Splitting
the denominator between $\mathcal{A}^{+}$ and its complement gives
$\hat\beta=\varkappa\hat\beta^{+}$. \hfill$\square$

\paragraph{Proposition~\ref{prop:firm}.} Condition on the firm's productivity $z$; its delivered
cost to buyer $a$ is $c=w_{r's}\tau_{r'as}/z$. Each competitor $r''$ has champion cost
$w_{r''s}\tau_{r''as}/\varphi_{r''}$ with
$\varphi_{r''}\sim\mathrm{Fr\acute{e}chet}(T_{a(r'')s},\theta)$, so
$\Pr(w_{r''s}\tau_{r''as}/\varphi_{r''}>x)=\exp(-T_{a(r'')s}(w_{r''s}\tau_{r''as})^{-\theta}x^{\theta})$.
By independence across competitors the minimum competing cost $M$ satisfies
$\Pr(M>x)=\exp(-\Phi^{-r'}_{as}x^{\theta})$. The firm supplies iff $M>c$, giving
\eqref{eq:app-firm-win}. Taking the complement,
$\Pr(\text{no supply}\mid z)=1-\exp(-\exp(\eta))$ with
$\eta=\ln\Phi^{-r'}_{as}+\theta\ln(w_{r's}\tau_{r'as})-\theta\ln z$, and
$\tau_{r'as}=d_{r'a}^{\alpha}$ gives \eqref{eq:app-firm-cloglog}. No step involves the variety
count. \hfill$\square$

\paragraph{Lemma~\ref{lem:qfree}.} Origin $r'$ wins variety $\rho$ for buyer $r$ iff
$w_{r's}\tau_{r'rs}/\varphi_{r'\rho s}<\min_{r''\neq r'}w_{r''s}\tau_{r''rs}/\varphi_{r''\rho s}$,
an event measurable with respect to $\{\varphi_{\cdot\rho s}\}$, $w$ and $\tau$ alone; the
$\varphi$ have distribution governed by $(T_{\cdot s},\theta)$. The win-somewhere event is the
union over $r$ of those events, hence measurable with respect to the same objects. Neither the
variety count nor any price or expenditure aggregate enters. \hfill$\square$

\paragraph{Unbiasedness of \eqref{eq:app-Gclosed}.} Let $k$ be the number of the $m$ simulated
varieties won by $r'$, each won independently with probability $q$. Draw $n\le m$ of the $m$
varieties uniformly without replacement; conditional on $k$, the probability that all $n$ are
losses is $\binom{m-k}{n}/\binom{m}{n}$. Marginalizing over the draws, each of the $n$ selected
varieties is a loss independently with probability $1-q$, so the unconditional probability is
$(1-q)^{n}$. Hence $\mathbb{E}[\binom{m-k}{n}/\binom{m}{n}]=(1-q)^{n}$ exactly, with the convention
$\binom{a}{n}=0$ for $a<n$. \hfill$\square$
